% Recovered core of the spacetime and Einstein-dynamics paper: premises,
% collar and overlap-gluing theorems, the Lorentz chain and the Einstein chain.
% The Standard Model gauge paper cites its overlap-gluing section. The file
% name is cited under this spelling by pinned manifests in code/a5_closure
% and by a frozen custody record, so it stays as it is.
\section{Recovered Core: Foundations and Structural Branches}\label{observer-patch-holography}

Observer-Patch Holography is a reconstruction program for fundamental physics. On the declared
screen-first global-support branch, physical data are charted by a horizon screen \(S^2\), and the
effective world is encoded by the mutual consistency of overlapping patch descriptions. This
support screen is distinct from a local carrier boundary and a federation overlap nerve. On the
producer branch, \(S^2\) enters only after the carrier-to-support map satisfies the spherical
incidence, mesh, cross-ratio, normalization, and refinement receipts. A finite reference
fixture without a composable seam triangle does not supply that map, since its
higher-overlap condition is vacuous. The source-derived producer supplies a
twelve-chart, thirty-seam, twenty-face incidence nerve, nonvacuous seam-triangle cocycles,
confluent phase repair, and a controlled oriented-\(S^2\) support limit on one declared tower.
Event, modular, scale, and laboratory attachments require separate premises.

This recovered-core section develops fixed-cutoff overlap repair, collar recovery, overlap
gluing and its transport obstructions, and the geometry-readout route to Lorentz frame
kinematics and conditional Einstein dynamics.

\subsection{Overview}\label{overview}

Theorem, lemma and section numbers inside this recovered core, such as
Theorem 4.2 or Section 2.3, follow the core's own internal numbering, not the
section numbering of the surrounding paper.

Overlap consistency turns local patch states into a constrained gluing problem. The
fixed-cutoff branch gives collar decompositions, edge centers and finite repair dynamics on the
declared physical quotient (Sections 2--3). Modular flow on caps supplies Lorentz-frame
kinematics and the three-dimensional observer-frame hyperboloid (Section 4). Certified
read-from provenance generates a finite informational poset and its longest
authenticated-parent-chain height; the rank-three source Gram quotient and a real axis define
the ambient \(1+3\) target carrier, and source height enters only through the temporal
coordinate of a scaled event placement. Positive null translations supply directional charges,
finite null tomography reconstructs a symmetric tensor, and fixed-cap generalized-entropy
stationarity with the stress, small-ball, coupling, vacuum and scale premises implies the
Einstein equation (Section 5). Bare finite consensus stops at quotient normal forms, and the
cap-normal \(H^3\) chart stops at observer frames; neither populates a spacetime.

The selected no-\(G\) scale certificate enters only the Newton-constant readout. The
source-causal continuum and Einstein theorems do not depend on the matter or sector-completion
branches.

\subsection{Model and Premises}\label{1-model-and-axioms}

\subsubsection{Observers and access model}\label{11-observers-and-access-model}

The support-local algebra-state-record reduct of an observer patch is
\(O_{\mathrm{red}}=(P_O, \mathcal{A}(P_O), \rho_O, R_O)\), where:

\begin{itemize}
\tightlist
\item
  \(P_O \subset S^2\) is a connected screen patch (the observer\textquotesingle s access region).
\item
  \(\mathcal{A}(P_O)\) is the von Neumann algebra associated to \(P_O\).
\item
  \(\rho_O\) is the local state, obtained by restricting the global state to \(\mathcal{A}(P_O)\).
\item
  \(R_O\) is a finite record algebra generated by stable internal record projectors within \(P_O\).
\end{itemize}

An operational observer additionally carries overlap interface algebras and restriction maps,
allowed update and repair instruments, and checkpoint data used for continuation. Observers are
internal self-reading patterns in the global state. A patch and its compatible marginal supply a
geometric or algebra-state chart. An operational observer requires the interface, record, repair,
and checkpoint conditions.

\subsubsection{Screen, patches, and algebra net}\label{12-screen-patches-and-algebra-net}

We work in a single static patch with a horizon screen \(S^2\). Each connected subregion \(P \subset S^2\) is assigned a von Neumann algebra \(\mathcal{A}(P)\). The net satisfies isotony:

\[
P \subset Q \implies \mathcal{A}(P) \subset \mathcal{A}(Q).
\]

A global state \(\omega\) is a positive linear functional on the inductive-limit algebra. Overlap consistency is imposed algebraically: for overlaps \(P_1 \cap P_2\), \(\omega\) restricted to \(\mathcal{A}(P_1 \cap P_2)\) is the same from either side.

\subsubsection{The premises A1--A3}\label{13-premises}

This paper states the premises A1--A3 of the observer-overlap foundation. The flagship synthesis~\cite{ophflagship} gives the account of the theory, and the executable observer paper~\cite{ophmodel} records which of these premises the construction supplies. Branch-local hypotheses are written at the theorem that consumes
them.

The support-visible BW scaling theorem handles the scaling step. The source-causal continuum and Einstein
theorems use the additional collar-recovery, edge-normalization, generalized-entropy,
common-domain, stress, scaling, universal-coupling, vacuum-reference, and scale-readout
interfaces stated in their own hypotheses; none of those interfaces belongs to A1--A3.

The screen premise A1 below defines the screen-first support branch. The producer branch reaches the
same premise only after its carrier-to-support receipts establish the global support screen.

% Shared panel of the theorem premises A1--A3. Included by papers that state
% the premises in full. Canonical machine source: claims/axiom_registry.yaml;
% reviewed reference: docs/AXIOM_REFERENCE.md. Do not hand-copy a competing
% formal block into a paper; include this fragment.

\paragraph{A1 (oriented observer-patch federation and spherical support).}
There exists an observer patch net on an oriented spherical screen: at every
finite resolution each local carrier has twelve primitive boundary ports
forming the vertices of an oriented triangular boundary with 30 edges and 20
faces, combinatorially the boundary of an icosahedron, and carriers join
through typed seams and coherent triple overlaps, refine to an oriented
spherical support, and expose local state, readback, records, repair moves,
and checkpoints. Their complete infinitesimal port response is represented
faithfully on a finite-dimensional unitary response space and is closed under
the commutators generated by ordered response composition. Formally, for
every regulator \(r\) in a
directed system there is a typed object
\(\mathfrak N_r=(\mathcal P_r,\mathcal A_r,\mathcal R_r,\mathcal I_r,
\mathcal U_r,\mathcal C_r,N_r,S_r,b_r)\):
a finite patch/overlap category with an isotone local algebra net and
central record algebras; a federation of finite carriers, each with twelve
primitive pairwise-orthogonal central port projections summing to one and a
boundary packet \(K=(P,E,F,o)\) with \(|P|=12\), \(|E|=30\), \(|F|=20\),
degree-five ports, five-cycle links, and coherently oriented two-face edges,
isomorphic to the icosahedral boundary complex with no preferred labels;
seam algebras with unital restrictions and coherent triple-overlap cocycles
forming the nerve \(N_r\); a finite edge-midpoint refinement support
\(S_r\) realized orientation-preservingly in \(S^2\) with mesh tending to
zero; and a source-bound degree-one bridge \(b_r:N_r\to S_r\), all commuting
with refinement. Each carrier also has
\(V_{r,i}=\mathbb R[P_{r,i}]\), a finite-dimensional complex Hilbert space
\(H_{r,i}\), and an injective real-linear map
\[
D_{r,i}:V_{r,i}\longrightarrow\mathfrak u(H_{r,i}).
\]
The image \(\mathfrak g_{r,i}=D_{r,i}(V_{r,i})\) is closed under
commutators and complete for the declared quotient-visible infinitesimal
port response. The primitive port probes span \(V_{r,i}\), no additional
public response direction is omitted, and
\(-\operatorname{Tr}(D_{r,i}(v)D_{r,i}(w))\) is positive definite. The
response construction is natural under admissible presentation equivalence
and refinement. A1 does not select a particular response table,
inverse-port law, Lie type, global group, particle interpretation, or metric
content.

\paragraph{A2 (observer agreement).}
Observers operating on the screen agree on the meaning of the data they
jointly interpret. Formally, the interpretation map
\(\mathcal J_r:\mathsf{Data}_r\to\mathsf{Meaning}_r\) from
observer-accessible data to operational meanings is natural with respect to
every visible overlap restriction, recharting, seam translation,
higher-overlap map, federation map, and refinement map: every declared
data-access diagram for accepted public data commutes after interpretation,
and \(\mathcal J_r\circ c_{s\to r}=C_{s\to r}\circ\mathcal J_s\) across
resolutions. A2 constrains accepted shared data only, and the domain
it quantifies over is supplied by the A1 interfaces rather than by A2
itself.

Accepted reversible overlap transports form a groupoid \(\mathcal O_r\).
For a complete carrier chart \(o\), its closed paths
\(\operatorname{Hol}_r(o)\) map onto the orientation-preserving incidence
automorphisms \(\operatorname{Aut}^{+}(K_{r,i})\). For every proper carrier
automorphism \(a\), A2 supplies a closed path \(\gamma_a\) and one
projective implementer
\([U_a]\in\operatorname{PU}(H_{r,i})\) satisfying
\[
\operatorname{Ad}_{[U_a]}D_{r,i}(v)=D_{r,i}(a\cdot v).
\]
Writing
\[
G_{D,r,i}^{0}
=
\left\langle
\exp(tD_{r,i}(v)):t\in\mathbb R,\ v\in V_{r,i}
\right\rangle^{0},
\]
the implementer is endogenous:
\[
[U_a]=[g_a c_a],
\qquad
g_a\in G_{D,r,i}^{0},
\quad
c_a\in C_{U(H_{r,i})}(\mathfrak g_{r,i}).
\]
Closed paths compose projectively and the construction commutes with
refinement. A unitary lift need not be unique. Scalar phases and centralizer
actions are permitted because they act trivially on the response algebra.
An independent spectator carrying a matching projective action cannot
replace the response-generated factor or contribute a public response
direction. The closed path, port action, and implementer must arise from the
same accepted source history in an executable realization.

The two response clauses support an abstract target-free compact Lie-type
theorem. They do not certify that an executable producer emitted a particular
current fixture. Such a certification must reconstruct the port generators,
commutator, and projective implementers from ordered source histories without
a named current model upstream. A2 does not imply global state
extension, repair termination, confluence, unique normal forms, record
durability, Byzantine safety, dissemination, or commutation of state
optimization with refinement; raw mismatch and repair precede acceptance.

\paragraph{A3 (conditional maximum randomness).}
Everything that observer agreement leaves unconstrained is maximally random.
Formally, at finite regulator \(r\) a state is a compatible family of local
normalized states on the accessible algebra net; \(\mathcal K_r\) is the
nonempty convex set of such families satisfying the finite observer-visible
constraints supplied by A1 and A2 through an A1-generated constraint
grammar with a factorization theorem; and, relative to an exact reference
family \(\tau_r\), an A1-generated observer cover \(\mathcal G_r\), and
strictly positive exact weights \(w_{r,P}\) constructed from quotient-visible
A1 data by a rule natural under admissible presentation equivalence, the
realized state is the
information projection. The cover restriction map is injective on
\(\mathcal K_r\); equivalently, at smooth points no nonzero feasible
tangent is invisible to the whole cover. The realized state is
\[
\rho_r=\operatorname*{arg\,min}_{\rho\in\mathcal K_r}
\sum_{P\in\mathcal G_r} w_{r,P}\,D(\rho_{r,P}\Vert\tau_{r,P}).
\]
The same projection rule applies separately to each declared optimizer
object type (ontic state, inference state, or transition distribution), and
statements about one type do not transfer to another.
This is equivalent to weighted local entropy maximization when every local
reference density is identity-proportional in its declared trace. "Random"
means least informative relative to the
declared reference and the complete agreement constraint set. A3
selects one state inside one fixed feasible space. It does not compare
unrelated state spaces or field contents, does not imply optimizer
compatibility across refinement, and does not imply recovery, alignment, or
mixing; each such statement carries its own theorem, interface, or
countermodel.

The complete formal statement of A1--A3, including its constraint map and
non-implications, is given in the
\href{https://github.com/FloatingPragma/observer-patch-holography/blob/r2053/docs/AXIOM_REFERENCE.md}
{premise reference}.

% Shared closure panel. Included by papers that state the premises A1--A3,
% immediately after PREMISE_BASIS.tex. It states the closure hypothesis
% of docs/BASIS.md, which no theorem takes as a premise. Reviewed
% references: docs/BASIS.md and docs/AXIOM_REFERENCE.md.

\paragraph{The closure hypothesis.}
The theorem premises A1--A3 state the library's hypotheses on the carrier and
its response, on internal transport and on local state selection. The closure hypothesis
states that the Universe is a timeless, self-referential structure whose
observers recover the specification of their own world from their records
and realize it, so that the description and the described system are one
object. Every quantity that
has both a construction-side reading and a readback-side reading must take
the same value once a typed bridge proves that the two readings denote one
invariant. The equality is forced by self-identity. Constructing that bridge
and the return map is part of the physics. Existence, uniqueness, and
stability are separate determinacy tests on the resulting equation. The name
given to such an equation carries no mathematical weight: after the typed
identification, unequal readings would describe two systems rather than the
single self-referential universe.

The closed object is the complete inhabited history, whose internal events
and responses to interventions are its proposed intrinsic computation.
This reading requires no external computer or playback time. It is distinct
from fixedness of a repair state, a public record, or a probability law, and
does not itself construct the history or prove a final cosmic endpoint.

The local screen-grain proposal
\[
P=\varphi+\sqrt{\pi}/A_T(P),
\]
with \(A_T(P)\) the inverse coupling returned by the declared source map,
has one interval-certified fixed point per declared map on its analytic domain.
Interpreting that root as the laboratory coupling additionally requires a
target-independent choice of the physical map, a typed bridge identifying
the two sides as readings of one quantity, and same-scheme hadronic spectral
transport to the Thomson endpoint.
The displayed map's certified fixed point sits at $P=1.6309720959$ and
returns $A_T=136.9948352$; the certified gauge-width variant's fixed
point sits at $P=1.6309682414$ and returns $137.0356601$. Through the
same outer equation the measured 2022 CODATA coupling
$\alpha^{-1}=137.035999177(21)$ reads back the grain value
$P=1.6309682094$. The comparisons stand at $3.0\times10^{-4}$ relative
for the displayed map and at $2.5\times10^{-6}$ for the gauge-width
variant, as diagnostics of the declared source maps. The latter absolute
residual is $3.3904\times10^{-4}$, compared with experimental standard
uncertainty $2.1\times10^{-8}$; this is not agreement at measured precision.
The gauge width supplies 99.2 percent of the gap between the bare root and the
measured value; the residual $3.3904\times10^{-4}$ is addressed to same-scheme
hadronic transport to the Thomson endpoint.

The direct global proposal is
\[
N=\log M_0(\mathfrak U_N).
\]
The bounded completion class defined by base agreement, positivity, and the
carrier bound does not select a unique cosmic value. Universal all-rung
membership in a complete A1--A3 capacity-source contract and the corresponding
executable-to-Lean bridge are absent. The incomplete antecedent determines no
direct \(N\), and no theorem extends the bounded-class conclusion to the
complete source class. A
positive result must complete the source antecedent and prove one physical
zero, either within the A1--A3 source contract or through a separately
named stronger source law. A separate
common-load branch starts from
\[
N_0=\pi\exp\!\left(\frac{6\pi}{P\,\alpha_U(P)}\right).
\]
On the finite collar branch, the declared total reserve expectation \(P/4\)
and six-class equidistribution give presence probability \(P/24\) for each
declared class. If one class is physically selected as the blocked event, its
scalar-weighted presence receipt is discharged, and that normalized
collar-survival factor is proved to act on the global capacity, the
corresponding candidate is
\[
N_{\mathrm{pres}}=N_0\left(1-\frac{P}{24}\right),
\qquad
\ln\frac{N_{\mathrm{pres}}}{\pi}
=
\frac{6\pi}{P\,\alpha_U(P)}
+\ln\left(1-\frac{P}{24}\right).
\]
The exponential alternative \(N_{\mathrm{Pois}}=N_0e^{-P/24}\) requires a
separate mean-count or continuum carrier. Exact neutral, one-class, and
six-class-total actions share the declared local datum, obey positive
composition and finite cut-count regrouping, and disagree globally. The
finite source therefore selects no action or blocked-event semantics. The
source class determines no named-law value, and its horizon branch has no
capacity object. A stronger source-derived action would
require the screen/electroweak bridge, physical common-load bridge, physical
\(\mathbb Z_6\) seam action, scalar-weighted receipt, and horizon-record
identification. On the source-forward numerical
branch,
\[
N_0=3.5321\times10^{122},
\]
so the finite-presence candidate is
\(3.2921\times10^{122}\) and the exponential candidate is
\(3.3001\times10^{122}\). The full weighted Planck
base-\(\Lambda\)CDM chain gives the comparison coordinate
\(3.3129\times10^{122}\). The respective residuals are
\(-0.63\) percent and \(-0.39\) percent. Both comparisons were exposed before
the branch choice and carry no predictive weight.


Collar recoverability, generalized-entropy structure, and quantum focusing lie outside A1--A3:
they are classified gravity interfaces with their own theorems, countermodels, and
required inputs, consumed field by field where the Einstein composition needs them. Sector
content is not selected by A1--A3: the conditional rank-45 tensor
candidate does not establish physical families or a local Spin carrier, while
scalar multiplicity and extra-sector completeness lack physical attachments.
Each theorem below lists the
declared completions it consumes.

\subsubsection{Closure values and theorem-local technical data}\label{14-assumptions-and-external-inputs}

The Newton constant of (G6) enters through one selected no-\(G\) scale certificate,
\[
\gamma_\star\equiv\frac{\ell_\star\nu_{\mathrm{Cs}}}{c},
\qquad
B_\star\equiv\frac{3\pi}{\ell_\star^2},
\qquad
G_{\mathrm{SI}}=\frac{c^3\ell_\star^2}{\hbar}.
\]
Dimensionless capacity and pixel coordinates fix dimensionless geometry, not an SI scale. On
the de Sitter branch, with capacity \(N_\star\) and pixel ratio
\(P_\star=a_{\mathrm{cell}}/\ell_\star^2\),
\[
\Lambda_\star\ell_\star^2=\frac{3\pi}{N_\star},
\qquad
\Lambda_\star a_{\mathrm{cell}}=\frac{3\pi P_\star}{N_\star},
\]
and these relations survive the common rescaling
\(\ell_\star^2\mapsto\lambda^2\ell_\star^2\),
\(a_{\mathrm{cell}}\mapsto\lambda^2a_{\mathrm{cell}}\),
\(\Lambda_\star\mapsto\lambda^{-2}\Lambda_\star\). The SI area \(\ell_\star^2\) and
\(G_{\mathrm{SI}}\) therefore require the scale certificate, which supplies
\(B_\star=\Lambda_\star N_\star\) and \(\ell_\star^2=3\pi/B_\star\). The branch certificate
used for the SI gravity row is
\[
B_\star=
3.60787391468\ldots
\times10^{70}\,\mathrm{m^{-2}},
\qquad
\ell_\star^2=
2.61228030237\ldots
\times10^{-70}\,\mathrm{m^2}.
\]
Rounded capacity displays such as \(N_\Lambda\simeq3.313\times10^{122}\) are not
high-precision inputs for the Newton row; with the Planck-2018 curvature benchmark, using only
the rounded capacity would shift the Newton row at the \(8.4\times10^{-4}\) level. The capacity
normalization uses the de Sitter static-patch entropy~\cite{gibbonsHawking77},
\(N_{\mathrm{scr}}=S_{\mathrm{dS}}=A_{\mathrm{dS}}/(4\ell_P^2)=3\pi/(\Lambda\ell_P^2)\).
Its selection from the source and the global-capacity closure belong to the cosmology
layer~\cite{ophdesittershock}.

The controlled scaling step for the Lorentz, null-modular, and Einstein statements is the
support-visible BW scaling theorem. Fixed-cutoff regulator bookkeeping, quasi-local propagation,
endpoint-Lipschitz control, and collar decomposition are cited from their preceding theorems.

\textbf{Local information projection at the regulator scale.} At the regulator scale \(\ell_{\mathrm{UV}}\), consider the premise-A3 subbranch carrying a global-extension interface and identity-proportional local references for which the state-determining weighted local objective is represented by one finite density matrix \(\omega\). The selected state maximizes the corresponding weighted local entropy subject to:

\begin{enumerate}
\def\labelenumi{\arabic{enumi}.}
\tightlist
\item
  A finite set \(\{O_a\}\) of gauge-invariant local operator densities, each supported on a ball of radius \(\le r_0 = O(\ell_{\mathrm{UV}})\).
\item
  Constraint equations \(\bigl\langle \sum_x O_a(x) \bigr\rangle = C_a\), one homogeneous global sum per density label \(a\), where \(x\) runs over the cells of the UV lattice. The constraints are the \(N_{\mathrm{con}}\) global sums, not one independent constraint per cell, so the number of independent constraints and of Lagrange multipliers is \(N_{\mathrm{con}}\) (plus the optional global constraints below) on every lattice, independent of the cell count.
\item
  Optionally, a finite number of global constraints (total energy, charge).
\end{enumerate}

This is the finite-matrix subbranch needed to derive the local Gibbs form of Lemma 2.6 from premise A3 with a cutoff-independent multiplier count. The global extension and the objective representation are explicit branch premises.

\textbf{Clarification (conditional maximum randomness and thermal equilibrium).} premise A3 is local state selection relative to an exact reference and a complete visible constraint set. It does not state that the universe is in thermal equilibrium. Non-equilibrium physics is modeled in the enlarged \emph{local-multiplier} family obtained by replacing the homogeneous multipliers \(\lambda_a\) with slowly varying fields \(\lambda_a(x)\), equivalently by imposing per-cell constraints \(\langle O_a(x)\rangle = c_a(x)\). That enlarged family has dimension proportional to the number of UV cells, so it is a function-space approximation regime with gradient error terms. The refinement-stable branch and the induced maps \(R_{\ell\to L}\) below concern the homogeneous \(N_{\mathrm{con}}\)-parameter branch under the separate optimizer-pushforward interface. Controlled violations of exact Markov additivity appear under the explicit mixing hypotheses stated later in Section 2.3. Equilibrium is an approximation regime with explicit error terms.

\textbf{Rotationally invariant constraint branch.} Constraint sets are \(\mathrm{SO}(3)\)-invariant on \(S^2\).

\textbf{Gauge as overlap redundancy at fixed cutoff.} Overlap identifications are not unique; the freedom that leaves overlap observables invariant forms a local groupoid.

\textbf{Central-defect subbranch.} Assume the overlap centers are identified with one fixed
abelian unitary coefficient group \(Z_\Sigma\) on which overlap transport acts trivially. Choose
unitary overlap implementers with \(U_{ji}=U_{ij}^{-1}\). On triple overlaps, when the failure of
strict coherence is central, define

\[
\Omega_{ijk}:=U_{ij}U_{jk}U_{ki}=z_{ijk},
\qquad z_{ijk}\in Z_\Sigma,
\]

equivalently \(U_{ij}U_{jk}=z_{ijk}U_{ik}\). The defect is data of the implementer lift; writing only
\(\mathrm{Ad}(z_{ijk})\) would erase it because conjugation by a central element is the identity.

\textbf{Collar double-scaling hypothesis.} There is one refinement family with
\(\delta\to0\), \(\ell_{\mathrm{UV}}\to0\), and uniform finite-range and strong
conditional Gibbs mixing constants such that

\[
I(A_\delta:D_\delta \mid B_\delta)_\omega
\le c|\partial C|_{\mathrm{UV}}e^{-\delta/\xi},
\qquad
\frac{\delta}{\xi}-\log|\partial C|_{\mathrm{UV}}\longrightarrow+\infty.
\]

The second condition controls the complete boundary-prefactored envelope. The weaker ratio
\(\delta/\ell_{\mathrm{UV}}\to\infty\) does not imply it. Section 2.3 gives the
finite-range theorem, explicit constants, a sufficient logarithmic schedule, and the finite-receipt
fields. On the exact central-interface branch the CMI is zero without this mixing premise.

\textbf{Geometric-subnet BW lift.} The continuum Lorentz statement is asked only on the support-visible extracted geometric subnet for caps, together with the controlled collar/refinement limit that carries the Markov and recovery remainders explicitly. The finite cap regulators are type-I algebras, but the scaling-limit observer algebra may leave that class, so the theorem is automorphism-level on the emitted geometric cap pair and the special type-I generator form is \(K_C=2\pi B_C+Z_C\) with \(Z_C\) central. The spacetime and Einstein paper proves the support-visible BW cap automorphism only when one cofinal branch carries the finite cap-normal certificate and the complete multiresolution algebra-state comparison package. It also proves the cap-normal \(H^3\) chart: \(q(\Omega)=(1,\Omega)\), \(n_C=(\cot\alpha,\csc\alpha\,\mathbf c)\), \(n_{gC}=\Lambda_gn_C\), and \(H^3\simeq\mathrm{SO}^+(3,1)/\mathrm{SO}(3)\), with cap half-spaces but no event-position consequence. Bare confluence underdetermines topology, dimension, and modular normalization by explicit countermodels. Separately, authenticated provenance generates the finite informational order and canonical source height. The exact rank-three source Gram quotient and an independent real axis define the ambient \(1+3\) Lorentz target carrier with an \(S^2\) null boundary. After multiplication by a supplied positive \texttt{timeScale}, source height supplies only the temporal coordinate of an event placement. It is not an execution rank, intrinsic poset dimension, physical clock, or physical event time. The null-net-standardness packet supplies derived positive translations, while the source-causal continuum certificate keeps physical order, spatial placement, history-family custody, count--volume calibration, independent dimension, manifoldlikeness, topology, uniqueness, and curvature convergence explicit. The closure packets prove finite null tomography and finite entropy/MaxEnt identities; local stress, Ward conservation, generalized-entropy stationarity, smooth curvature, and the composed Einstein reading remain conditional on their named analytic and physical supplies. No source-derived family carrying the complete continuum packet is constructed (Theorems 4.3c--4.3h).

\textbf{Derived quasi-local propagation and modular-locality control.} At scale \(\ell_{\mathrm{UV}}\), the automorphism group generated by the MaxEnt generator of Lemma 2.6, or more generally by any branch generator lying in the same bounded-support algebraic closure, has a finite Lieb--Robinson velocity \(v_{\mathrm{LR}}\): for local operators \(A, B\) supported on regions separated by distance \(d\),

\[
\|[A(t), B]\| \le c \|A\| \|B\| \min(|R_A|, |R_B|) e^{-(d - v_{\mathrm{LR}}|t|)/\xi}
\]

for \(d > v_{\mathrm{LR}}|t|\), where \(\xi = O(\ell_{\mathrm{UV}})\).

This is the branch-internal control statement that turns the quasi-local structure of the
local-Gibbs branch into explicit support control for time-evolved operators. It is part of the
premise-A3 branch stated at regulator scale.

\begin{quote}
\textbf{Approximate modular covariance on the fixed-cutoff realized branch.} Under the fixed-cutoff realized presentation, Lemma 2.6 (local Gibbs), the derived quasi-local propagation bound above, and the collar double-scaling plus mixing hypotheses of Section 2.3, the modular flow \(\sigma_t^{\omega,C}\) maps \(\mathcal{A}(R)\) into a slightly thickened region algebra:

\[
\sigma_t^{\omega,C}(\mathcal{A}(R)) \subseteq \mathcal{A}(R^{+v_{\mathrm{mod}}|t|})
\quad \text{up to error } \eta(d - v_{\mathrm{mod}}|t|),
\]

where \(R^{+s}\) denotes the \(s\)-neighborhood thickening, \(v_{\mathrm{mod}}\) is a "modular propagation velocity" controlled by \(v_{\mathrm{LR}}\) and local norm bounds, and \(d\) is the distance from \(R\) to \(\partial C\).

\end{quote}

This statement is a regulator-scale control; the tangent-limit package used in Theorem 4.2 is stated there with its own hypotheses.

\textbf{Refinement-stable information-projection branch.} Choose a family of coarse-graining channels \(\Phi_{\ell\to L}\) between UV scale \(\ell\) and IR scale \(L\) satisfying the separately declared optimizer-pushforward and refinement-tail interface. Then on the realized branch one may write

\[
\Phi_{\ell\to L}\bigl(\omega_{\ell}(\lambda)\bigr)=\omega_{L}\!\bigl(R_{\ell\to L}(\lambda)\bigr),
\]

for an induced map \(R_{\ell\to L}\) on the common finite-dimensional homogeneous multiplier family. The existence of \(R_{\ell\to L}\) as a self-map of that same family is the load-bearing cross-scale premise. Reusing the density labels \(a\) does not establish it. Coarse-graining a finite-range Gibbs family generically generates interactions outside the retained density list. The closure defect, the moment-matching I-projection realizing \(R_{\ell\to L}\), and its trace-norm residual bound are constructed in Definition 2.6a and Lemma 2.6b. Exact vanishing of the defect along the realized branch remains part of the declared optimizer-pushforward interface. Under that interface, the regulator states lie in a shared multiplier family, and the realized branch is a persistent trajectory or invariant subset under the induced refinement maps. This state-side persistence statement does not upgrade fixed-cutoff edge labels to a transportable refinement-persistent sector colimit.

\textbf{Fixed-cutoff realized presentation.} At a UV scale \(\ell_{\mathrm{UV}}\), local patch algebras are type-I with finite-dimensional Hilbert spaces, and gauge-as-gluing is realized as a boundary action on a lifted finite-dimensional presentation. Section 2.3 proves the resulting fixed-cutoff collar package directly from overlap consistency on this realized presentation, while distinguishing the lifted fixed-point algebra from the sector-preserving collar algebra used in the edge-center (EC) theorem.

External mathematical inputs include SSA and recovery theorems (Petz 1986, 1988; Fawzi and
Renner 2015).

The Einstein route also uses the fixed-volume area-variation identity for small geodesic balls and
the local quadratic-polarization argument in the tensor upgrade. The null modular bridge and
small-ball kernel are carried internally on the compact surface rather than imported from a
separate EFT small-ball law.

\textbf{Recovered-core theorem boundary.} The Lorentz, null-modular, and Einstein recovered core
runs through the support-visible automorphism theorem. Fixed-cutoff statements remain regulator
statements.

Full citations appear in the References.

\subsubsection{Notation}\label{15-notation}

\begin{itemize}
\tightlist
\item
  \(\rho_C\): reduced state on cap \(C\).
\item
  \(K_C := -\log \rho_C^{\omega}\): modular Hamiltonian of the reference state.
\item
  \(B_C\): geometric generator of the cap-preserving conformal dilation.
\item
  \(S_{\mathrm{gen}}(C)\): generalized entropy on a cap.
\item
  \(\ell_{\mathrm{UV}}\): UV length scale of the refined screen net.
\item
  \(\delta\): collar width around a cap boundary.
\end{itemize}

\begin{center}\rule{0.5\linewidth}{0.5pt}\end{center}
\subsection{Information-Theoretic Tools}\label{2-information-theoretic-tools}

\subsubsection{Strong subadditivity and Markov states}\label{21-strong-subadditivity-and-markov-states}

For any tripartite state \(\rho_{ABC}\),

\[
I(A:C \mid B) := S(AB) + S(BC) - S(B) - S(ABC) \ge 0.
\]

Exact Markov states satisfy \(I(A:C \mid B) = 0\) and admit a recovery map:

\[
\rho_{ABC} = (\mathrm{id}_A \otimes \mathcal R_{B\to BC})(\rho_{AB}).
\]

\subsubsection{Approximate recovery}\label{22-approximate-recovery}

If \(I(A:C \mid B) \le \varepsilon\) in nats, there exists a CPTP recovery map \(\mathcal{R}\) with

\[
\| \rho_{ABC} - (\mathrm{id}_A \otimes \mathcal R)(\rho_{AB}) \|_1
\le 2\sqrt{1-e^{-\varepsilon}}
\le 2\sqrt{\varepsilon}.
\]

\subsubsection{Collar refinement and sufficient mechanisms}\label{23-collar-refinement-and-sufficient-mechanisms}

Fix a cap \(C \subset S^2\) with boundary circle. For collar width \(\delta\) define

\[
B_\delta := \{x \in S^2 : \mathrm{dist}(x,\partial C) \le \delta\},
\qquad
A_\delta := C \setminus B_\delta,
\qquad
D_\delta := (S^2 \setminus C) \setminus B_\delta.
\]

Then \(S^2 = A_\delta \cup B_\delta \cup D_\delta\) with \(A_\delta\) and \(D_\delta\) interacting only through \(B_\delta\). The collar double-scaling hypothesis is the requirement that \(I(A_\delta:D_\delta \mid B_\delta) \to 0\) in the refinement limit. At fixed regulator scale, the same finite-dimensional realized presentation supplies the patch-net and overlap-gluing data used below. We therefore isolate that realized presentation first, then separate the exact-Markov and quantitative routes.

\textbf{Derived regulator realization.} Choose a finite UV cellulation at scale \(\ell_{\mathrm{UV}}\) and let \(R\) be a finite union of cells. Hilbertize each cell\textquotesingle s finite local data to a finite-dimensional space \(\tilde{\mathcal H}_i\cong \mathbb C^{n_i}\). The extended algebra before quotienting by overlap redundancy is

\[
\widetilde{\mathcal A}(R)=\mathcal B(\tilde{\mathcal H}_R),
\qquad
\tilde{\mathcal H}_R=\bigotimes_{i\subset R}\tilde{\mathcal H}_i.
\]

Overlap-preserving changes of local trivialization act only on cut data. On each finite-dimensional chart, their unitary image has compact closure, so one may represent the boundary gluing redundancy by a compact group \(G_{\partial R}\subset U(\tilde{\mathcal H}_R)\). The lifted boundary-invariant endomorphism algebra is

\[
A_{\mathrm{inv}}(R)=\widetilde{\mathcal A}(R)^{G_{\partial R}}
=\mathcal B(\tilde{\mathcal H}_R)^{G_{\partial R}}.
\]

This is the fixed-cutoff realized presentation used throughout the collar analysis. The EC theorem works on the invariant-state realization and uses the sector-preserving induced collar algebra there; it does not apply to the entire fixed-point algebra on the unreduced tensor product.

The quotient repair law acts on this overlap-invariant data: the global map is the finite
quotient normal-form operator. Its termination, confluence, schedule independence, the
separated cofinal inverse-limit bridge and the distributed-run certificate are proved in the
consensus paper~\cite{ophconsensus}. Those results establish at most a common finite quotient
metric or pseudometric. Its identification with Lorentzian spacetime goes through the
source-causal continuum certificate of Theorem 4.3e and the modular branch below.

\textbf{Theorem 2.3 (EC from regulated overlap gluing).} For a collar \(B_\delta\) around a cap boundary \(\Sigma\), there is a canonical decomposition

\[
H_{B_\delta} = \bigl(\tilde{\mathcal H}_{B_L}\otimes \tilde{\mathcal H}_{B_R}\bigr)^{G_\Sigma}
= \bigoplus_{\alpha}
\bigl(H_{b_L^{\alpha}} \otimes H_{b_R^{\alpha}}\bigr),
\]

and the sector-preserving collar algebra
\[
A_{\mathrm{EC}}(B_\delta):=\bigoplus_\alpha
\bigl(B(H_{b_L^\alpha}) \otimes B(H_{b_R^\alpha})\bigr)
\]

with

\[
Z(A_{\mathrm{EC}}(B_\delta)) = \bigoplus_\alpha \mathbb C\,\mathbf 1_\alpha,
\]

such that \(\mathcal{A}(A_\delta B_\delta)\) acts only on \(H_{b_L^\alpha}\) and \(\mathcal{A}(B_\delta D_\delta)\) acts only on \(H_{b_R^\alpha}\) within each block.

\textbf{Proof.} Split the collar into half-collars \(B_L\) and \(B_R\) meeting on \(\Sigma = \partial C\). By the realized regulator presentation above, the physical collar Hilbert space is the diagonal invariant subspace \((\tilde{\mathcal H}_{B_L} \otimes \tilde{\mathcal H}_{B_R})^{G_\Sigma}\). Decompose each side into irreps:

\[
\tilde{\mathcal H}_{B_L} = \bigoplus_\alpha (V_\alpha \otimes H_{b_L^\alpha}),\qquad
\tilde{\mathcal H}_{B_R} = \bigoplus_\beta (V_\beta^* \otimes H_{b_R^\beta}).
\]

Then

\[
\tilde{\mathcal H}_{B_L} \otimes \tilde{\mathcal H}_{B_R}
= \bigoplus_{\alpha,\beta}
(V_\alpha \otimes V_\beta^*) \otimes
(H_{b_L^\alpha} \otimes H_{b_R^\beta}).
\]

By Schur\textquotesingle s lemma,

\[
(V_\alpha \otimes V_\beta^*)^{G_\Sigma}
\cong
\begin{cases}
\mathbb C, & \alpha=\beta,\\
0, & \alpha\ne\beta.
\end{cases}
\]

Therefore the invariant subspace is

\[
H_{B_\delta} = \bigoplus_\alpha (H_{b_L^\alpha} \otimes H_{b_R^\alpha}),
\]

as claimed. The induced sector-preserving collar algebra is

\[
A_{\mathrm{EC}}(B_\delta) = \bigoplus_\alpha
\bigl(B(H_{b_L^\alpha}) \otimes B(H_{b_R^\alpha})\bigr),
\]

so the center is generated by the block projectors. Adjacent region algebras act on the left or right factor only because the gluing action is supported on \(\Sigma\). QED.

\textbf{Remark.} On the central-defect subbranch, replace \(G_\Sigma\) by its central extension. The sector label \(\alpha\) then ranges over irreps of the extension; the decomposition is unchanged.

We refer to the decomposition in Theorem 2.3 as \textbf{edge-center completion (EC)}.

\textbf{Exact Markov route.} If \(I(A_\delta:D_\delta\mid B_\delta)=0\), the Hayden--Jozsa--Petz--Winter (HJPW) structure theorem yields a blockwise factorization of the state over \emph{some} direct-sum/tensor decomposition of \(\mathcal H_{B_\delta}\); that decomposition is state-dependent, and HJPW does not by itself identify it with the preselected EC factors \(b_L^\alpha,b_R^\alpha\) of Theorem 2.3. The exact identity used for literal Markov-modular equalities below is the \textbf{EC-aligned} normal form

\[
\rho_{A_\delta B_\delta D_\delta}
= \bigoplus_\alpha p_\alpha
\bigl(\rho_{A_\delta b_L^\alpha} \otimes \rho_{b_R^\alpha D_\delta}\bigr),
\qquad
I_\omega(A_\delta:D_\delta \mid B_\delta)=0,
\]

and the identification of the HJPW factors with the EC factors is a genuine additional state condition, the \textbf{Markov-split alignment hypothesis (MSA)}. It is not implied by exact Markovity: with four qubits \(A,B_L,B_R,D\) and \(\rho=\Phi_{AB_R}\otimes\Phi_{B_LD}\) (Bell pairs across the cut, one center block), one has \(I(A:D\mid B)=0\) but \(I(A:B_R)=2\ln 2\), so the state is exactly Markov yet not of the displayed EC-aligned form, which would force \(I(A:B_R)=0\). All literal Markov-modular equalities below therefore assume exact Markovity \emph{plus} MSA, or an explicitly stated idealized limit that supplies both; The spacetime and Einstein paper states MSA formally and records this counterexample. MSA is checkable state by state: The spacetime and Einstein paper proves it equivalent to a splitting of \(\log\rho\) into commuting one-sided terms supported on \(A b_L^\alpha\) and \(b_R^\alpha D\), to the existence of a \(\rho\)-preserving conditional expectation onto either one-sided edge algebra (a Takesaki commuting square over the collar center \cite{takesaki72}), and to blockwise vanishing one-sided mutual information \(I(Ab_L^\alpha:b_R^\alpha D)=0\). On the local information-projection branch this reduces the derivation of MSA to a sharp condition on the effective Hamiltonian: every cross-cut coupling must act through the central sector label, as it does in lattice-gauge-type regulators where the interface energy is a function of the conserved flux. On that declared \textbf{central-interface branch} The spacetime and Einstein paper derives both MSA and exact collar Markovianity as theorems (descent of one-sided invariants and boundary charges through the edge-center quotient), and it adopts the central-interface normal form as an explicit named branch premise, so the literal Markov-modular identities hold on the declared branch without a further state hypothesis. This premise is independent of A1--A3: The spacetime and Einstein paper exhibits a finite regulator package that satisfies observer agreement, schedule-independent transactional repair, and an exactly closed refinement-channel family while its retained constraint family carries a boundary-invariant noncentral cross-cut coupling. Observer agreement and repair do not force the premise. Off the declared branch MSA is carried as an explicit hypothesis.

On the separately declared rate-controlled family obtained by pushing the optimizer forward, Theorem 2.3 supplies the kinematic edge-center decomposition. Exact Markovity is one idealized route. The following lemma and theorem supply the quantitative decay route when the approximation is kept explicit.

\begin{quote}
\textbf{Lemma 2.6 (interior local information projection implies local Gibbs form).} Under premise A3 on the subbranch with identity-proportional local references, state-determining aggregation equivalent to the declared entropy objective, homogeneous global-sum constraints \(\langle\sum_x O_a(x)\rangle = C_a\), and the finite-dimensional regulator realization above, assume also that the constraint vector lies in the relative interior of the attainable moment set, so the selected state is faithful on the retained sector. Then the selected state has the finite-multiplier Gibbs form

\[
\omega = \frac{e^{-H_{\mathrm{eff}}}}{\mathrm{Tr}\,e^{-H_{\mathrm{eff}}}},
\qquad
H_{\mathrm{eff}} = \sum_x \sum_a \lambda_a O_a(x) + H_{\mathrm{sec}},
\]

where the sum runs over UV cells \(x\) and density labels \(a\), with one cell-independent multiplier \(\lambda_a\) per label. For the finite-range collar theorem, \(H_{\mathrm{sec}}\) is either a sum of uniformly bounded finite-range charge densities or is central and fixed on the chosen superselection sector. A genuinely nonlocal noncentral term falls outside the theorem. Under this condition the effective Hamiltonian is finite-range in UV graph units, and the multiplier count matches the constraint count \(N_{\mathrm{con}}\) plus the number of retained sector terms on every lattice.

\textbf{Proof.} On a finite-dimensional algebra, an interior entropy maximizer subject to linear constraints \(\mathrm{Tr}(\rho O_i) = c_i\) is given by finite Lagrange multipliers:

\[
\rho = \frac{e^{-\sum_i \lambda_i O_i}}{\mathrm{Tr}\,e^{-\sum_i \lambda_i O_i}}.
\]

Strict concavity of von Neumann entropy ensures uniqueness. Here the constrained operators are the \(N_{\mathrm{con}}\) global sums \(O_a := \sum_x O_a(x)\) plus the retained sector terms, so the exponent has one shared multiplier per density label. MaxEnt gives the Gibbs form on this faithful/interior branch; the displayed locality restriction on \(H_{\mathrm{sec}}\) is separate and load-bearing. A boundary constraint can instead force a rank-deficient optimizer (for example \(\mathrm{Tr}(\rho\lvert1\rangle\!\langle1\rvert)=0\) on \(\mathbb C^2\)); it has no finite full-algebra multiplier and must be treated on its support or by infinite penalties. Such a boundary branch does not supply the stagewise faithfulness used below. Had the constraints instead been imposed per cell, \(\langle O_a(x)\rangle = c_a(x)\), the multipliers would be cell-dependent fields \(\lambda_a(x)\) and the family would grow with the lattice; that enlarged family belongs to the non-equilibrium approximation regime of the clarification above and is not the branch used in the refinement argument. QED.
\end{quote}

\begin{quote}
\textbf{Definition 2.6a (closure defect and induced refinement map).} Fix successive regulator scales \(\ell\) (finer) and \(L\) (coarser), let \(\{\omega_L(\lambda')\}\) be the homogeneous global-sum exponential family at scale \(L\), and let \(\Phi_{\ell\to L}\) be an admissible coarse-graining channel. For fine-scale multipliers \(\lambda\) set \(\sigma := \Phi_{\ell\to L}(\omega_\ell(\lambda))\) and define the \textbf{closure defect}

\[
\varepsilon_{\ell\to L}(\lambda) := \inf_{\lambda'} D\bigl(\sigma \,\big\|\, \omega_L(\lambda')\bigr),
\]

with \(D\) the quantum relative entropy. When the infimum is attained at a unique \(\lambda^\ast\), set \(R_{\ell\to L}(\lambda) := \lambda^\ast\).

\textbf{Lemma 2.6b (I-projection residual bound).} If the constrained operators at scale \(L\) together with the identity are linearly independent, then \(\lambda' \mapsto D(\sigma\|\omega_L(\lambda'))\) is smooth and strictly convex; for faithful \(\sigma\) the minimizer \(\lambda^\ast = R_{\ell\to L}(\lambda)\) exists, is unique, and is characterized by moment matching \(\langle S_c\rangle_{\omega_L(\lambda^\ast)} = \langle S_c\rangle_\sigma\) for every constrained operator \(S_c\); and

\[
\bigl\|\sigma - \omega_L(R_{\ell\to L}(\lambda))\bigr\|_1 \le \sqrt{2\,\varepsilon_{\ell\to L}(\lambda)},
\]

so the family reproduces every bounded expectation up to \(\|B\|\sqrt{2\varepsilon_{\ell\to L}(\lambda)}\). The separately declared exact optimizer-pushforward interface is the statement \(\varepsilon_{\ell\to L}(\lambda) = 0\) along the realized branch. In that case \(\sigma = \omega_L(R_{\ell\to L}(\lambda))\), and the induced refinement map is this moment-matching I-projection.

\textbf{Proof.} Writing \(\log\omega_L(\lambda') = -\sum_c \lambda'_c S_c - \log Z_L(\lambda')\), the objective is \(-S(\sigma) + \sum_c \lambda'_c \langle S_c\rangle_\sigma + \log Z_L(\lambda')\), and the Hessian of \(\log Z_L\) is the Duhamel (Kubo--Mori) covariance matrix of the constrained operators, positive definite under the stated independence; this gives strict convexity, uniqueness, and the moment-matching first-order condition, while attainment for faithful \(\sigma\) is standard finite-dimensional exponential-family duality~\cite{wichmann63}. The displayed residual bound is the quantum Pinsker inequality \(D(\rho\|\tau) \ge \tfrac12\|\rho-\tau\|_1^2\)~\cite{hiaiOhyaTsukada81} evaluated at \(\lambda^\ast\), and \(\varepsilon_{\ell\to L}(\lambda) = 0\) forces \(\sigma = \omega_L(\lambda^\ast)\) by strict positivity of relative entropy off the diagonal. QED.

\textbf{Acceptance check.} A finite two-lattice test implementing this package (the constraint/multiplier count against the displayed \(N_{\mathrm{con}}\)-dimensional family on successive lattices, the moment-matching projection, the residual bound, a generic decimation channel with strictly positive closure defect, and the exactly closed transverse-field subfamily where \(R_{\ell\to L}\) acts as the identity) is included with the repository sources.

\textbf{Numerical projection and closure are distinct.} For faithful \(\sigma\) on
\(\mathbb C^d\), choose a Hilbert--Schmidt orthonormal basis of the observable
space \(V\) modulo identity, and write its traceless Hamiltonian as
\(H(\theta)=\sum_j\theta_jQ_j\), with \(\rho_\theta=e^{-H}/\operatorname{Tr}e^{-H}\). The invariant moment residual is
\(\eta=\|P_V(\sigma-\rho_\theta)\|_{\mathrm{HS}}\); a raw residual in caller
units is not a convergence criterion. Coercivity and convexity give
\[
D(\rho_*\|\rho_\theta)\le B:=\eta\left(\|\theta\|_2+
\frac{\log d}{\lambda_{\min}(\sigma)}\sqrt{\frac{d-1}{d}}\right),
\qquad \|\rho_*-\rho_\theta\|_1\le\min\{2,\sqrt{2B}\},
\]
where \(\rho_*\) is the unique exact projection. This optimizer error is
separate from the closure defect \(D(\sigma\|\rho_*)\); the Pythagorean
identity adds the two. The implementation evaluates these bounds in floating
point, without outward-rounded enclosures or an exact-zero certificate.
Invertible changes of observable coordinates and identity offsets preserve
the state; scalar energy origins shift \(\log Z\) and must be retained
in partition and free-energy calculations. Compensated scalar accumulation
and a stable relative-entropy remainder avoid cancellation of small finite
effects; neither supplies a spectral support certificate or recovers input
differences lost to rounding.
\end{quote}

\textbf{Strong conditional mixing premise.} Let \(G_\ell\) be the UV cell graph. Assume local
dimension at most \(q\), interaction degree at most \(\Delta_0\), interaction diameter at most
\(r_0\) graph layers, and
\[
\sup_x\sum_{X\ni x}\|\beta\Phi_\ell(X)\|\le J_0
\]
with constants independent of the cut and refinement stage. Let \(m_\ell\) be the
resolved collar width in graph layers and set \(\delta=m_\ell\ell\). For the faithful Gibbs state, define
\[
\mathbf J_\rho(A:D\mid B)
:=\log\rho_{ABD}+\log\rho_B-\log\rho_{AB}-\log\rho_{BD}.
\]
The strong mixing premise is a boundary-anchor expansion
\[
\mathbf J_\rho(A_\delta:D_\delta\mid B_\delta)
=\sum_{z\in\partial^{\mathrm{UV}}_{r_0}C}E_{z,\ell,\delta},
\qquad
\|E_{z,\ell,\delta}\|_\infty
\le\kappa e^{-(m_\ell-r_0)/\zeta},
\]
with \(\kappa,\zeta\) uniform in the cut, boundary condition, and stage. This conditional or
matrix mixing condition is stronger than ordinary two-point exponential clustering. Local Gibbs
form, a Lieb--Robinson bound, or a Hamiltonian spectral gap does not imply it for a general
noncommuting Gibbs state.

\begin{quote}
\textbf{Theorem 2.5 (finite-range conditional Gibbs mixing gives collar recovery).}
Under Lemma 2.6 and the strong conditional mixing premise,
\[
\boxed{
I_\omega(A_\delta:D_\delta\mid B_\delta)
\le
\kappa|\partial C|_{\mathrm{UV}}e^{-(m_\ell-r_0)/\zeta}
\le
c|\partial C|_{\mathrm{UV}}e^{-\delta/\xi_\ell}}
\]
with
\[
\xi_\ell=\zeta\ell,
\qquad
c=\kappa e^{r_0/\zeta}.
\]

\textbf{Proof.} Embedding the marginal logarithms in the tripartite algebra gives
\[
I(A:D\mid B)_\rho
=\operatorname{Tr}\!\left[\rho_{ABD}\mathbf J_\rho(A:D\mid B)\right].
\]
The boundary expansion and trace-norm duality bound this expectation by the sum of the boundary
term norms. The second form follows from \(\delta=m_\ell\ell\). QED.
\end{quote}

For one uniform family, the sharp sufficient limit condition is
\[
\frac{\delta}{\xi_\ell}-\log|\partial C|_{\mathrm{UV}}\longrightarrow+\infty.
\]
If
\[
|\partial C|_{\mathrm{UV}}\le a(\ell_0/\ell)^p,
\qquad
\delta(\ell)\ge\zeta(p+\eta)\ell\log(\ell_0/\ell),
\quad p\ge0,\quad \eta>0,
\]
then \(\delta/\ell\to\infty\), and
\[
I(A_\delta:D_\delta\mid B_\delta)
\le ca(\ell/\ell_0)^\eta\longrightarrow0.
\]
This lower bound on \(\delta\) does not imply \(\delta\to0\).
Physical collar shrinking remains the separate double-scaling hypothesis
stated above. It is realized, for example, by choosing
\(\delta(\ell)=\zeta(p+\eta)\ell\log(\ell_0/\ell)\), or by rounding
its graph-layer width upward with an additional error at most \(\ell\).
For a counterexample to inferring shrinking from the lower bound alone,
take \(\ell_n=1/(4n)\), \(\ell_0=\delta_n=1/4\),
\(p=0\), and \(a=\zeta=\eta=1\). Then
\(\delta_n\ge\ell_n\log n\) and \(\delta_n/\ell_n=n\), but
\(\delta_n=1/4\) does not shrink. A zero-Hamiltonian product Gibbs chain
has zero conditional matrix defect and CMI and satisfies all the local
finite-range and mixing bounds in this example; it fails precisely the
separately declared shrinking-collar hypothesis.
The ratio \(\delta/\ell\to\infty\) alone is insufficient. For
\(\ell_n=\ell_0e^{-n^2}\), \(\delta_n=\zeta n\ell_n\), and
\(|\partial C_n|_{\mathrm{UV}}=e^{n^2}\), that ratio diverges while the upper envelope grows as
\(e^{n^2-n}\).

A finite receipt records the interaction range and norm bound, treatment of global sector terms,
boundary-cell count, graph separation, regional CMI in nats, matrix-defect norm, predeclared
\(\kappa,\zeta\), log-envelope slack, Fawzi--Renner error, held-out cuts, and the rate margin
\(\delta/(\zeta\ell)-\log|\partial C|_{\mathrm{UV}}\). Fitted constants and local random-triplet
CMI remain finite proxies. They do not certify a cofinal limit.

This bound is the quantitative hinge for constructive gluing.

\subsubsection{Conformal-modular fixed-point microphysics}\label{27-conformal-modular-fixed-point-microphysics-cmfp}

On the faithful/interior local finite-constraint branch of premise A3 with identity-proportional reference, the logarithm of the selected state is a quasi-local UV generator. A separately declared optimizer-pushforward and refinement-tail interface places the branch inside one common finite-dimensional multiplier family. At each regulator stage the patch and cap algebras are type-I matrix algebras. The finite regulator class is not refinement-closed, so the scaling-limit observer algebra may leave that class. The fixed-cutoff analysis proves the local-Gibbs result and propagation. Endpoint-Lipschitz control of \emph{reduced} modular Hamiltonians requires the additional exact-Markov/central-interface and endpoint-control branch stated below; it does not follow from finite-range locality alone. The cross-scale interface supplies refinement stability. The geometric cap pair with geometric modular flow is obtained in the support-visible sense: the target algebra is the extracted geometric subnet, the support-visible modular matrix elements converge under regularization, and weak-\(*\)/GNS extraction plus support-readable modular covariance and ordered cut-pair rigidity gives the geometric cap automorphism.

\textbf{Local finite-constraint information-projection branch.} The constraint family \(\mathcal C\) is generated by finitely many gauge-invariant local densities \(\{O_a(x)\}\) of UV range \(O(\ell_{\mathrm{UV}})\), constrained through their homogeneous global sums \(\langle\sum_x O_a(x)\rangle = C_a\). premise A3 selects the fixed-stage information projection. Retention of the same finite label set under refinement belongs to the separate optimizer-pushforward and refinement-tail interface.

\textbf{Theorem 2.6 (Interior local constraints imply a local-Gibbs form).} If the MaxEnt constraints are expectations of finitely many quasi-local operators \(\{O_a\}\) with bounded support size at scale \(\ell_{\mathrm{UV}}\), and the prescribed moment vector lies in the relative interior of the attainable moment set so that the optimizer is faithful on the retained sector, then the entropy maximizer is
\[
\omega \propto \exp\!\left(-\sum_a \lambda_a O_a\right),
\]
so the MaxEnt generator \(H_{\mathrm{MaxEnt}}=-\log\omega\) is a UV-range quasi-local sum. This is exactly the local-Gibbs form used later.

\textbf{Proof.} The finite-dimensional interior exponential-family theorem gives finite Lagrange multipliers \(\lambda_a\). Boundary optima require restriction to their support or extended infinite multipliers and are not included in this faithful branch. QED.

\textbf{Derived propagation control on the same branch.} Because \(H_{\mathrm{MaxEnt}}\) is a finite-range or quasi-local sum on a finite type-I regulator net, standard Lieb--Robinson estimates~\cite{liebRobinson72} apply to the automorphism group it generates, or to any branch generator lying in the same bounded-support algebraic closure. Thus there is a finite propagation velocity \(v_{\mathrm{LR}}\) and constants \(C,\xi\) such that for local observables \(A_X,B_Y\),
\[
\bigl\|[\tau_t(A_X),B_Y]\bigr\|
\le
C\,\|A_X\|\,\|B_Y\|\,\min(|X|,|Y|)\,
e^{-(d(X,Y)-v_{\mathrm{LR}}|t|)/\xi}.
\]
\textbf{Additional endpoint-control branch.} Locality of the global Gibbs generator does not by
itself make reduced interval modular Hamiltonians local; the four-qubit counterexample in
Theorem 4.3d rules out that inference. On the exact-Markov/central-interface branch, assume in
addition a specified endpoint-to-regulator-cell map and a uniform local norm/site-density bound
such that, after subtracting the central endpoint block, the operator-norm sum of terms changed
between allowed intervals \(I,I'\) is at most \(L|I'\Delta I|\). On this explicitly named branch,
the local split and that geometric norm bound give bounded-interval endpoint-Lipschitz matrix elements,
\[
\left|\langle\psi,(K[I']-K[I])\phi\rangle\right|
\le
C_{\psi,\phi,I_{\max}}\,|I'\Delta I|,
\]
used later in the null-modular bridge. Thus quasi-local propagation follows on the faithful local-Gibbs branch, while reduced-modular endpoint control is a separate exact-Markov/central-interface receipt with the stated endpoint geometry and uniformity.

\textbf{Refinement-stable multiplier branch.} The separately declared optimizer-pushforward and refinement-tail interface preserves the same finite constraint family under coarse-graining and places the regulator states in one common finite-dimensional multiplier family. This is a substantive renormalization condition because coarse-graining a finite-range Gibbs family generically generates interactions outside a fixed density list. The ``refinement-stable'' language used later means persistence along the stable or fixed branch of this multiplier family. This state-side notion is enough to compare realized states across cutoffs. FixedCutoffBosonicSectorCategory constructs the fixed-stage tensor-generated category and its forgetful fiber. RefinementFunctorAndFiberDescent constructs the refinement pullbacks and compatible fibers only on a cofinal tail carrying the compact-gauge refinement receipt. The cofinal compact-gauge witness uses the same receipt. The separate finite Standard Model route uses complete reversible response and endogenous overlap transport to force the local Standard Model gauge Lie algebra. Incidence and target-blind readback derive the carrier module and \(R=-J\). The finite conditional fixture supplies the charge-conjugate rank-15 projector pair, hypercharge lattice, common \(\mathbb Z_6\) kernel, and maximal faithful matter image. Source derivation of that current, matter action, and global quotient is not supplied. Under the complete-band and operational-cost premises, the exact screen theorem selects the rank-three response band. Physical family attachment, extra-sector completeness, and the charged-lepton and electroweak quantitative branches carry their own premises. premise A3 supplies neither the finite current realization nor the sector-colimit data or laboratory attachment.

\textbf{Scaling limit and algebraic type.} The regulator presentation gives a family of finite type-I algebras
\[
\mathcal A_{\ell}(C)\cong \mathcal B(\mathcal H_{C,\ell}).
\]
A scaling limit of this family need not be type I. In the continuum-QFT case of interest one expects the local limit algebra \(\mathcal A_\infty(C)\) to be non-type-I, typically type III. Accordingly the fundamental modular datum in the limit is the automorphism group of the pair \((\mathcal A_\infty(C),\omega_\infty^C)\), not a density matrix inside \(\mathcal A_\infty(C)\).

\textbf{Corollary 2.6 (Geometric modular action on paired-certified caps).} For any support-visible
extracted scaling-limit geometric cap pair
\((\mathcal A_\infty^{\mathrm{geo,sv}}(C),\omega_\infty^{\mathrm{geo},C})\) emitted on a branch
carrying both the finite cap-normal certificate and the independently complete same-branch
multiresolution algebra-state comparison package, let \(\alpha_{\lambda_{\widehat C}(s)}\) be the automorphism
induced by the independently normalized BW-framed cap flow. Then
\[
\sigma_t^{\omega_\infty^{\mathrm{geo},C}}
=
\alpha_{\lambda_{\widehat C}(2\pi t)}.
\]
If the limit cap algebra happens to be type I, this may be written as
\[
K_C=2\pi B_C+Z_C,\qquad Z_C\in Z(\mathcal A(C))_{\mathrm{sa}}.
\]
In the generic continuum case, the automorphism identity is the complete statement.

\textbf{Proof.} This is The spacetime and Einstein paper's finite cap-net BW theorem applied to the extracted
geometric cap pair. The finite cap-normal input supplies BW framing, support-order faithfulness,
held-out oriented cross-ratio rigidity, the geometric support flow, independently normalized
geometric \(2\pi\)-KMS comparison, and wrong-normalization controls. The separate
multiresolution algebra-state comparison input supplies the complete algebra-state reference tower,
cap-family-uniform mixed-GNS convergence, comparison maps, and modular support covariance. QED.

\textbf{Support envelope from net regularity (not an alternative derivation).}
The outer-regularity / minimal-support condition used later lets one read a
support envelope for the modular image from the scaling-limit net. By itself it
does not establish modular covariance; that equality comes from the
paired certificate in Corollary 2.6, or must be assumed as an additional
two-sided net-automorphism property.

\textbf{(NR) Outer regularity / minimal support.} For any operator \(O\), the intersection of all connected regions \(P\) with \(O\in\mathcal A(P)\) is again a connected region, denoted \(\mathrm{supp}(O)\).

\textbf{Proposition 2.7 (Modular support-envelope inclusion).} Under (NR), define for any region \(R\subset C\)
\[
f_t^C(R)
:=
\bigcup_{O\in \mathcal A(R)}
\mathrm{supp}\!\left(\sigma_t^{\omega,C}(O)\right).
\]
Then
\[
\sigma_t^{\omega,C}(\mathcal A(R))\subseteq\mathcal A(f_t^C(R)).
\]
Equality does not follow from (NR). It requires, for example, a two-sided
order-isomorphism of the local net compatible with \(\sigma_t\) and
\(\sigma_{-t}\), together with the relevant additivity/local-generation
property. Those are substantive covariance data, not consequences of knowing
minimal operator supports.

\textbf{Proof.} For every \(O\in\mathcal A(R)\), (NR) assigns
\(\sigma_t^{\omega,C}(O)\) a minimal support contained in \(f_t^C(R)\); isotony
therefore places that operator in \(\mathcal A(f_t^C(R))\). Taking all such
operators gives the stated inclusion. Surjectivity onto the entire envelope
algebra is not supplied by the fact that \(\sigma_t\) is an automorphism of the
global algebra. QED.

\textbf{Null-surface modular structure.} On the same support-visible scaling branch and extracted geometric-subnet branch, the null-surface modular machinery narrows as follows:
\begin{itemize}
\item
\textbf{Fixed-cutoff null-strip bridge.} The null-strip package proves transferred cut-center data, the theorem-local inherited left/right strip-split condition needed for the spatial-collar-type tensor decomposition, exact-or-controlled four-term strip additivity on one inherited strip model, and renormalized endpoint-Lipschitz control up to the weak tail generator.  On the scaling branch the geometric action supplies the half-sided relation; the separately registered common-vector and relative-commutant standardness receipts make it a standard half-sided modular inclusion whose Borchers--Wiesbrock consequence is an explicit positive null-translation generator on its Stone domain with the affine half-line modular relation. The same half-line family fixes the generator's endpoint-derivative identity internally; its local null-stress reading consumes the separately supplied effective-locality and same-source bridge.
\item
\textbf{Derived half-sided modular relation.} Nested null half-line algebras satisfy the half-sided covariance and containment relation on the declared geometric scaling branch by Corollary 5.2e.  To obtain a \emph{standard} half-sided modular inclusion, the common reference vector must additionally be cyclic and separating for both half-line algebras and cyclic for their relative commutant.  On that separately registered standardness receipt, Borchers--Wiesbrock identifies the positive null-translation generator on its Stone domain together with the affine half-line modular relation~\cite{wiesbrock93}.
\item
\textbf{Weak continuity and finite variation.} The bounded-interval and half-line endpoint-Lipschitz control follows on the exact-Markov/central-interface branch only after the specified endpoint-to-cell map and uniform norm/site-density bound of the endpoint-control clause. It then defines the weak tail generator at fixed cutoff. The support-visible scaling theorem together with the extracted geometric cap pair supplies the continuum null-generator setting in which that weak-tail data can be matched to the geometric null modular action of the relativistic phase.
\end{itemize}

\textbf{Constraint set specification.} On the local finite-constraint branch, the ``correct fixed-cap constraint set'' becomes explicit: constraints are the local conserved charges of the symmetries used in the derivation:
\begin{enumerate}
\def\labelenumi{\arabic{enumi}.}
\item
\textbf{Edge/cap label constraints:} fix the distribution of collar-sector labels, equivalently \(\langle L_C\rangle\) for each cap size, giving the area term.
\item
\textbf{Gauge charges:} fix boundary flux or charge operators.
\item
\textbf{Geometric (conformal) charges:} fix the expectation of the conformal Killing charges that preserve the cap, i.e.\ the generator \(B_C\) or its microscopic lattice approximation.
\end{enumerate}

MaxEnt therefore selects the unique finite-stage invariant state compatible with those conserved charges. The support-visible BW scaling theorem is the continuum statement: if the limit algebra is non-type-I, the geometric modular action on that branch is outer rather than an inner density-matrix identity.

\textbf{Focusing boundary.} QNEC has rigorous QFT proofs in broad settings~\cite{bousso16,balakrishnan19}, with hypotheses (null-surface locality and analyticity, modular-flow/causality/defect-OPE structure) that have not been verified on the reconstructed net, and QFC is strictly stronger than QNEC~\cite{bousso15qfc}. Quantum focusing is an optional gravity interface and is not consumed by the Einstein theorem:
\begin{itemize}
\item
EC plus the declared state, edge-reference, and entropy-split premises identify \(S_{\mathrm{gen}}=S_{\mathrm{bulk}}+\langle L_C\rangle\) on the finite branch (Section 5.4).
\item
The null-stress and Einstein theorem branch supplies the kinematic inputs required by an internal QNEC/QFC proof; that derivation is not supplied.
\end{itemize}

\textbf{Summary.} The CMFP package therefore separates the dependency structure into two layers:
\begin{itemize}
\item
\textbf{Fixed-stage and cross-scale consequences:} premise A3 with the declared local constraint and reference package gives the local-Gibbs form and quasi-local propagation on the faithful/interior branch. Bounded-interval reduced-modular endpoint-Lipschitz control additionally consumes the exact-Markov/central-interface split and the specified endpoint-to-cell uniform bound. The separate optimizer-pushforward interface supplies the realized state-side refinement branch along which later transport questions are asked.
\item
\textbf{Support-visible scaling statement:} the scaling limit emits the support-visible geometric cap pair and ordered cut-pair rigidity holds on it. On that branch the limit algebra may be non-type-I, the modular action may be outer, and the \(2\pi\) normalization follows.  The later null half-sided-inclusion bridge additionally consumes its registered common-vector and relative-commutant standardness receipts.
\end{itemize}

The theorem route works on the support-visible quotient. The realized transported cap-local system packages the geometric cap-local test family, the projectively compatible transported marginal family, and the asymptotic transport-equivalence certificate. The regularized transport theorem below replaces the unavailable full-algebra lower spectral floor. Local weak-\(*\) extraction and GNS gluing then emit the support-visible scaling-limit cap pair, support-readable modular covariance reads the modular group as a cap-local support map, and ordered cut-pair rigidity collapses the residual cap-preserving conformal freedom to the unique hyperbolic subgroup.

\textbf{Proposition 2.6a (Recoverability is not modular geometry: common-floor collapse countermodel).} Exact or asymptotically exact Markov recovery at finite cutoff does not by itself imply a full-algebra lower spectral floor along refinement. In a two-dimensional matrix algebra, let
\[
\rho_n=\begin{pmatrix}e^{-n}&0\\0&1-e^{-n}\end{pmatrix}.
\]
Each \(\rho_n\) is faithful at finite \(n\), and tensoring it with any fixed finite exact-Markov collar factor gives a full-rank exact-Markov collar family. Nevertheless \(\lambda_{\min}(\rho_n)=e^{-n}\to0\), so there is no positive lower bound along the refinement family. On the off-diagonal matrix unit \(E_{12}\), the modular generator carries the logarithmic gap
\[
[ -\log\rho_n, E_{12} ] = n\,E_{12}+O(1)E_{12},
\]
and the unregularized modular transport has no finite common-floor limit on that direction. Thus collar Markovity and finite-stage faithfulness are not enough for the false stronger full-algebra BW lift. The observer-facing theorem uses the support-visible regularized replacement, with the exact-Markov comparison family checked only on fixed collars and made replacement-independent after weak-\(*\)/GNS extraction.

\textbf{Theorem 2.6b (regularized support-visible modular transport).} Fix a local collar model of finite dimension \(d_{m,\delta}\). Let \(\rho_n\) be the transported physical collar marginal, \(\widehat\rho_n\) the exact-Markov comparison marginal on the same collar model, and \(\Delta_n=\|\rho_n-\widehat\rho_n\|_1\). For \(a>0\), set \(K_a(\rho)=-\log(\rho+a\mathbf 1)\). For every bounded collar observable \(O\) in the support-visible algebra,
\[
\left|\operatorname{Tr}\rho_n O\bigl(K_a(\rho_n)-K_a(\widehat\rho_n)\bigr)\right|
\le
\|O\|\left(
  \frac{4\Delta_n}{a}
  + d_{m,\delta}a
  +4\Delta_n|\log a|
\right).
\]
Consequently, if \(a_n\downarrow0\) is chosen with
\[
\Delta_n/a_n\to0,\qquad
d_{m,\delta}a_n\to0,\qquad
\Delta_n|\log a_n|\to0,
\]
then the regularized support-visible modular matrix elements converge on that fixed collar model.

\textbf{Proof sketch.} On the spectral interval \([a,\infty)\), the logarithm is operator-Lipschitz at the scale used above by its integral representation. Splitting the comparison into the support above \(a\), the \(a\)-tail, and the trace-distance error gives the displayed bound. The three displayed conditions make the three error terms vanish. QED.

\textbf{Definition 2.6c (explicit multiresolution reference branch).}\label{def:explicit-multiresolution-reference-branch}
The finite-to-continuum bridge is evaluated on a declared reference branch, not on an
unstructured sequence of finite plots. A regulator is
\[
r=(m,L,b),
\]
where \(m\) is the refinement depth, \(L\) is the finite support or volume scale, and \(b\)
is a boundary/phase label. The screen branch uses nested icosahedral subdivisions of \(S^2\);
Euclidean branches use dyadic mesh \(a_m=a_0 2^{-m}\), with an additional temporal spacing
\(a_t\) when a transfer matrix is present. Each added cell, collar, edge-center, or
support-visible channel \(j\) carries
\[
D_j=\bigoplus_{\alpha\in S_j}M_{d_{j,\alpha}}(\mathbb C)
\]
and a faithful detail state \(\tau_j\). At regulator \(r\),
\[
\widetilde M_r=\bigotimes_{j\in\mathcal J_r}D_j,\qquad
\widetilde\omega_r=\bigotimes_{j\in\mathcal J_r}\tau_j .
\]
A finite-depth local presentation circuit \(\Gamma_r:\widetilde M_r\to M_r\) changes from
multiresolution coordinates to the OPH patch, port, collar, and gauge coordinates. For
\(r\preceq s\), write
\[
\widetilde M_s=\widetilde M_r\otimes D_{s\setminus r},\qquad
\tau_{s\setminus r}=\bigotimes_{j\in\mathcal J_s\setminus\mathcal J_r}\tau_j .
\]
The canonical refinement and coarse-graining maps are
\[
\iota_{rs}(A)=
\Gamma_s\!\left[\Gamma_r^{-1}(A)\otimes\mathbf1\right],
\]
and
\[
Q_{sr}(X)=
\Gamma_r\!\left[
(\operatorname{id}\otimes\tau_{s\setminus r})
\bigl(\Gamma_s^{-1}(X)\bigr)
\right].
\]
The embedded conditional expectation is \(E_{sr}=\iota_{rs}Q_{sr}\). The reference state is
\(\widehat\omega_r=\widetilde\omega_r\circ\Gamma_r^{-1}\).

\textbf{Theorem 2.6d (exact modular-compatible reference tower).}\label{thm:explicit-oph-reference-tower}
The data
\[
(M_r,\widehat\omega_r,\iota_{rs},E_{sr})_{r\preceq s}
\]
form a directed tower. The maps \(\iota_{rs}\) are unital injective
\(*\)-homomorphisms with \(\iota_{st}\iota_{rs}=\iota_{rt}\). The \(Q_{sr}\) are faithful
state-preserving completely positive coarse-grainings with \(Q_{tr}=Q_{sr}Q_{ts}\). The maps
\(E_{sr}\) are faithful conditional expectations onto \(\iota_{rs}(M_r)\), and
\[
\widehat\omega_s\circ\iota_{rs}=\widehat\omega_r,\qquad
\widehat\omega_r\circ Q_{sr}=\widehat\omega_s,\qquad
\widehat\omega_s\circ E_{sr}=\widehat\omega_s .
\]
The finite reference modular groups are also exactly compatible:
\[
\sigma_t^{\widehat\omega_s}(\iota_{rs}A)
=
\iota_{rs}(\sigma_t^{\widehat\omega_r}(A)).
\]
This theorem uses no full-algebra spectral floor uniform in \(r\); in bare coordinates it is only
the identity \(A\mapsto A\otimes\mathbf1\), partial expectation against faithful detail states,
and conjugation by \(\Gamma_r,\Gamma_s\).

\textbf{Theorem 2.6e (canonical renormalization and correlation Cauchy bound).}\label{thm:martingale-renormalization}
Let \(\mathfrak A\) be the \(C^\ast\)-inductive limit with embeddings \(j_r:M_r\to\mathfrak A\).
There is a unique state \(\widehat\omega\) with
\(\widehat\omega\circ j_r=\widehat\omega_r\), and the maps \(Q_{sr}\) induce
state-preserving conditional expectations
\[
\mathbb E_r:\mathfrak A\to j_r(M_r)
\]
such that \(\|\mathbb E_r(A)-A\|\to0\) for every \(A\in\mathfrak A\). The finite-regulator
realization of a continuum test observable is
\[
\mathcal R_r(A)=j_r^{-1}(\mathbb E_rA).
\]
For \(\|A_k\|\le M\),
\[
\left|
\widehat\omega_r\!\left(\prod_{k=1}^n\mathcal R_r(A_k)\right)
-
\widehat\omega\!\left(\prod_{k=1}^n A_k\right)
\right|
\le
M^{n-1}\sum_{k=1}^n\|A_k-\mathbb E_rA_k\|.
\]
Thus the operator mixing matrix and additive subtraction at finite cutoff are outputs of
\(\mathcal R_r\), not free fits to the target continuum answer. Independent lattice-spacing,
volume, and boundary sweeps must publish the corresponding martingale-tail or physical-error
envelopes before finite regulator data are read as continuum evidence.

\textbf{Theorem 2.6f (compact-time modular convergence on fixed local algebras).}\label{thm:compact-time-modular-convergence}
Let \(\rho^{\rm phys}_{s\downarrow r}\) be the density matrix of the physical state at fine stage
\(s\), restricted to the embedded fixed local algebra \(M_r\), and let
\[
\epsilon_{s,r}:=\|\rho^{\rm phys}_{s\downarrow r}-\widehat\rho_r\|_1 .
\]
For a regularizer \(\eta_s>0\), define
\[
K_{s\downarrow r}^{(\eta_s)}
=-\log(\rho^{\rm phys}_{s\downarrow r}+\eta_s\mathbf1).
\]
If \(\lambda_r=\lambda_{\min}(\widehat\rho_r)>0\), then for every \(A\in M_r\) and every
\(T<\infty\),
\[
\sup_{|t|\le T}
\left\|
e^{-itK_{s\downarrow r}^{(\eta_s)}}A e^{itK_{s\downarrow r}^{(\eta_s)}}
-
\sigma_t^{\widehat\omega_r}(A)
\right\|
\le
2T\|A\|
\left[
\frac{\epsilon_{s,r}}{\eta_s}
+\log\!\left(1+\frac{\eta_s}{\lambda_r}\right)
\right].
\]
Hence \(\eta_s\to0\) and \(\epsilon_{s,r}/\eta_s\to0\) give uniform-on-compact-time convergence
on each fixed support-visible local algebra. For a countable cofinal local test tower, pointwise
\(\epsilon_{s,r}\to0\) admits a diagonal subsequence and one cutoff schedule. The phrase
``finite-stage modular action'' below means this transported, regularized local action; it does
not mean the unregularized full finite-cap modular group.

\textbf{Cap-family uniformity clause.} For application to the finite cap-net BW theorem, the
compact-time estimate is required uniformly over the declared finite nondegenerate cap family and
the fixed separating support-visible test tower. The operator-norm estimate then implies the
quadratic mixed-GNS Cauchy residual because \(\|XD\|_{2,\omega}\le\|X\|\|D\|_{2,\omega}\). The
reference tower supplies the analytic engine; the finite cap-normal certificate adds the
support-order, BW-frame, oriented cross-ratio, geometric \(2\pi\)-KMS, and wrong-scale controls.

\textbf{Theorem 2.6g (finite positive transfer and Lorentzian handoff condition).}\label{thm:positive-transfer-gate}
On Euclidean repair branches, let \(P_{C,r}\) be weighted resampling inside the fibers of the
complete repaired visible datum for active collar \(C\).  Ref.~\cite{observableNormalForms}
proves that this finite operator is the orthogonal conditional expectation and gives the exact
support, equal-fiber-row, and weighted detailed-balance matrix receipt that recognizes it.  A
concrete active collar repair kernel is identified with \(P_{C,r}\) only when its independently
derived transition matrix passes that receipt; constructing the matrix from the target formula is
not a verification.  Choose \(c_{C,r}\ge0\), and set
\[
L_r^{\rm rep}=\sum_C c_{C,r}(I-P_{C,r}),\qquad
T_r(a_t)=e^{-a_tL_r^{\rm rep}} .
\]
Then \(L_r^{\rm rep}=L_r^{\rm rep\,*}\ge0\), \(0<T_r(a_t)\le I\), \(T_r(a_t)\mathbf1=\mathbf1\),
and the stationary Euclidean path measure is reflection positive:
\[
\mathbb E_r[(\Theta F)F]\ge0
\]
for positive-time cylinder functions \(F\). A Lorentzian unitary statement requires this finite
positivity plus either exact shell-additive direct-limit semigroup identities or a Mosco /
strong-resolvent substitute, including vacuum-projection and intertwiner convergence. Weak-\(*\)
state convergence alone is not a transfer theorem.

\textbf{BW-side closure boundary.} Proposition 2.6a blocks the stronger unregularized full-algebra
conclusion from the proved finite-cutoff package. Theorems 2.6c--2.6g supply the explicit
reference-tower certificate needed for the finite-to-continuum reading: continuum correlations
come from the conditional-expectation renormalization map, compact-time modular group convergence
is only for transported regularized physical marginals on fixed support-visible local algebras,
and Lorentzian transfer needs reflection positivity or an equivalent positive-transfer
certificate. Theorem 4.2 then uses exactly the observer-facing support-visible content to close
the BW/geometric cap-pair automorphism statement. Without the multiresolution regulator
certificate, finite regulator data remain regulator evidence. A black-box AQFT/BW reconstruction
route would require separately verifying conformal-net properties such as locality, additivity,
duality, standardness, positive energy, and suitable modular inclusions.
\subsection{Overlap Consistency and Gluing}\label{3-overlap-consistency-and-gluing}

The constructive part of overlap consistency is the tree-gluing theorem below. The structural part is the origin of the gluing redundancy itself. On the ordinary or central-defect branch, gauge-as-gluing is the finite-regulator overlap redundancy of local chart presentations: once the finite-dimensional overlap algebras and their allowed chart-change groupoid are put in the block-complete spatial presentation specified below, overlap-consistent rechartings of a cut form a compact unitary transition system. The collar theorem of Section 2.3 should therefore be read with its boundary group \(G_\Sigma\) understood as shorthand for this spatialized compact boundary action, not as an automatic consequence of an arbitrary algebra representation. On the genuinely noncentral branch, the corresponding declared weak gluing data are encoded by a compact crossed-module change system. Block permutations and other groupoid data are retained rather than silently treated as inner automorphisms.
A second structural point is UV underdetermination. If each patch Hilbert space is tensored with an inert finite ancillary factor and the observable patch algebras are embedded as \(\mathcal A(P)\otimes \mathbf 1\), then the physical observables, collar conditional mutual information, carried Markov errors, and quotient normal forms are unchanged. So the physical UV branch is fixed only modulo gauge or implementation hiding together with such ancillary stabilization, not a unique microscopic presentation; the unique theorem-grade object is the induced family of terminal expectation functionals on the declared physical observable algebras, not a preferred microscopic representative.

\subsubsection{Constructive gluing on tree covers}\label{31-constructive-gluing-on-tree-covers}

\textbf{Theorem 3.1 (tree gluing).} Let a rooted tree of patches satisfy a tree-ordered overlap structure and a tripartite factorization \((A_k, B_k, C_k)\) at step \(k\). If a target state \(\rho^{\ast}\) obeys \(I(A_k:C_k \mid B_k) \le \varepsilon_k\), then there exist recovery maps \(\mathcal{R}_k\) such that

\[
\| \rho^{\ast}_{A_k B_k C_k} - (\mathrm{id}_{A_k} \otimes \mathcal{R}_k)(\rho^{\ast}_{A_k B_k}) \|_1
\le \delta_k,
\]

with

\[
\delta_k = 2\sqrt{1-e^{-\varepsilon_k}}
\le 2\sqrt{\varepsilon_k},
\]

The iteratively glued state \(\hat{\rho}\) satisfies

\[
\| \hat{\rho} - \rho^{\ast} \|_1 \le \min\left(2, \sum_{k=2}^n \delta_k\right).
\]

\textbf{Proof.} Induct on \(k\). The recovery error contracts under CPTP maps, so the errors add. QED.

\subsubsection{Gauge-as-gluing and loops}\label{32-gauge-as-gluing-and-loops}

At finite regulator scale, the fixed-cutoff gauge-as-gluing package identifies the overlap-consistency redundancy of local chart presentations. Choose block-complete faithful finite-dimensional spatial presentations of the overlap algebras on a connected cut \(\Sigma\), and require the allowed overlap-chart isomorphisms to be represented by their unitary normalizers. A finite-dimensional \(C^*\)-algebra is a direct sum of full matrix blocks: its star automorphisms permute isomorphic central blocks and act inner within each block. They are spatial in the canonical block-complete presentation, but need not be inner in the algebra itself (an abelian two-block swap is the elementary counterexample). Fixing a reference chart and retaining any allowed block permutations, the compact closure of the subgroup generated by the spatial recharting unitaries is the boundary gluing group \(K_\Sigma\); before fixing the reference chart, the same data form a compact unitary groupoid.

\textbf{Proposition 3.2a (Spatialized gauge-as-gluing at finite regulator).} Let a finite regulator chart be chosen for the patches meeting along a connected interface \(\Sigma\), together with the block-complete spatial presentations and unitary-normalizer realization of the allowed chart-change groupoid just stated. Then overlap consistency determines a compact unitary transition system on the cut data. On the ordinary or central-defect branch, after a reference chart is fixed, this transition system reduces to a compact boundary group \(K_\Sigma\), and when the triple-overlap defect is central its projective composition law lifts to a genuine action of a compact central extension \(\widehat K_\Sigma\). On this realized presentation, gauge names the overlap redundancy; the spatialization data are explicit hypotheses rather than consequences of finite dimensionality alone.

\textbf{Proof.} The structure theorem for finite-dimensional \(C^*\)-algebras decomposes each chart algebra into full matrix blocks. An allowed star isomorphism permutes isomorphic central blocks and is inner on each matched block; the chosen block-complete presentation implements the combined operation by a unitary normalizer. The subgroup generated by those unitaries has compact closure inside a finite-dimensional unitary group. If triple-overlap defects are central, the resulting projective composition law lifts to a central extension. QED.

\textbf{Lemma 3.2b (trees vs loops).} If the patch adjacency graph is a tree, one can choose local charts \(h_i\) so that the overlap labels satisfy \(g_{ij} = h_i^{-1} h_j\) on all edges. If loops exist, the loop holonomy

\[
H(\gamma) = g_{i_1 i_2} g_{i_2 i_3} \cdots g_{i_n i_1}
\]

transforms by conjugation at the chosen basepoint under local frame changes.
Its conjugacy class, and therefore every class-function value (including every
one-dimensional character value), is invariant. In particular, nontriviality
is invariant and obstructs global trivialization. QED.

\textbf{Proposition 3.2ba (Finite path ratio, character phase, and cyclic
quantization).} For reversal-compatible group-valued edge labels on a typed
finite path system, the ordered transport ratio of two paths with common
endpoints is exactly the holonomy of the loop obtained by following the first
and reversing the second. Every one-dimensional unitary character sends this
identity to an exact relative phase. A vertex rechart conjugates based
holonomy, so holonomy itself is invariant for an abelian structure group and
the character phase is invariant for an arbitrary group. There is an explicit
four-vertex path/face control whose two declared triangular faces have trivial
holonomy while a separate undeclared, unfilled loop and its two-arm relative
character phase are nontrivial. Finally, if the loop holonomy factors through
a supplied \(\mathbb Z_n\) sector, every character phase is an \(n\)th root of unity.

\textbf{Scope.} The result is machine checked for supplied finite paths, edge
labels, character, faces, and cyclic-sector factorization. The path/face
control declares exactly two triangular faces; it neither claims flatness on
undeclared triangles nor proves a puncture or noncontractibility statement.
No theorem here constructs the connection from observer records or identifies
the phase with a physical gauge flux, spacetime loop, charge, laboratory arm,
detector, or
interference fringe.

\textbf{Corollary 3.2c (Collar consequence on the ordinary or central-defect branch).} Let \(B_\delta = B_L \cup B_R\) be a collar around a cap boundary \(\Sigma\), and set \(\widehat K_\Sigma = K_\Sigma\) on the ordinary branch. Then the EC theorem of Section 2.3 has the derived interpretation

\[
H_{B_\delta}
=
(\tilde{\mathcal H}_{B_L}\otimes \tilde{\mathcal H}_{B_R})^{\widehat K_\Sigma}
\cong
\bigoplus_\alpha
(H_{b_L^\alpha}\otimes H_{b_R^\alpha}),
\]

with

\[
Z(A_{\mathrm{EC}}(B_\delta)) = \bigoplus_\alpha \mathbb C\,\mathbf 1_\alpha.
\]

The right half-collar carries the contragredient action because it sees the inverse transport across the same cut. Exact Markov normal forms used later require the additional state hypothesis \(I_\omega(A_\delta:D_\delta \mid B_\delta)=0\) \emph{together with} the Markov-split alignment hypothesis of Section 2.3, or the explicitly stated idealized recoverability limit supplying both; the block decomposition itself is kinematic and follows from the derived boundary action.

\textbf{Proof sketch.} Decompose the left boundary data into irreps \((V_\alpha \otimes H_{b_L^\alpha})\) of \(\widehat K_\Sigma\) and the right boundary data into the dual modules \((V_\beta^* \otimes H_{b_R^\beta})\). Then Schur's lemma leaves a singlet only when \(\alpha=\beta\), producing the displayed direct sum. QED.

\subsubsection{Loop obstruction class (central defect)}\label{33-loop-obstruction-class-central-defect}

On the central-defect subbranch, identify the overlap centers with one fixed abelian unitary
coefficient group \(Z_\Sigma\) and assume overlap transport acts trivially on it. Choose unitary
overlap implementers normalized by \(U_{ji}=U_{ij}^{-1}\) and define central defects
\(z_{ijk}\) at the implementer level by

\[
\Omega_{ijk}:=U_{ij}U_{jk}U_{ki}=z_{ijk}
\in Z_\Sigma,
\]

or equivalently \(U_{ij}U_{jk}=z_{ijk}U_{ik}\). This implementer-level relation records the
projective multiplier; the automorphism-level expression \(\mathrm{Ad}(z_{ijk})\) alone would be
the identity and would not determine \(z_{ijk}\).

Then \(\{z_{ijk}\}\) is an untwisted \v{C}ech 2-cocycle with coefficients in \(Z_\Sigma\), and its cohomology class \([z]\) is gauge invariant. Central defects do not obstruct the collar block decomposition: they only replace \(K_\Sigma\) by its central extension \(\widehat K_\Sigma\). The condition \([z]=0\) is exactly the strictifiability of the projective triangle defect. It is not by itself global path independence: at least one resulting strict edge \(1\)-cocycle must also have trivial represented holonomy around every closed overlap loop. With a nontrivial action on varying centers, the corresponding statement uses the twisted/local-system \v{C}ech differential instead. (A full proof for the fixed-coefficient subbranch appears in Section 6.4 below, in the algebra-net language.)

\subsubsection{Non-central obstruction (2-group cocycle)}\label{34-non-central-obstruction-2-group-cocycle}

When defects are not central, the natural coefficient data is a crossed module \((H \to G)\) with an action of \(G\) on \(H\) by conjugation. Here \(G\) is the reconstructed gauge group, and \(H\) is the unitary group acting on edge multiplicity spaces, with boundary map \(\partial: H \to G\).

A crossed module is a homomorphism \(\partial: H \to G\) together with an action of \(G\) on \(H\) such that

\[
\partial(g \triangleright h) = g\,\partial(h)\,g^{-1},
\qquad
\partial(h) \triangleright h' = h h' h^{-1}.
\]

On a good cover \(\{P_i\}\), a weakly coherent gluing is encoded by:

\[
g_{ij}: P_{ij} \to G,\qquad h_{ijk}: P_{ijk} \to H,
\]

obeying the 2-cocycle conditions

\[
g_{ij} g_{jk} = \partial(h_{ijk}) g_{ik},
\]

and on quadruple overlaps,

\[
h_{jkl} h_{ijl} = (g_{ij} \triangleright h_{ikl}) h_{ijk}.
\]

Gauge changes act by 1- and 2-cochains in the standard way for crossed-module cohomology.

Write
\[
q_\Sigma=[(g,h)]\in\check H^2(N_\Sigma,H_\Sigma\to G_\Sigma)
\]
for the full crossed-module gluing orbit. It classifies the equivalence orbit of the pair \((g,h)\)
but does not in general determine a unique ordinary \(G_\Sigma\)-valued \(1\)-cocycle class after
strictification, because the map below need not be injective. Let
\[
\begin{aligned}
\iota_\Sigma:\check H^1(N_\Sigma,G_\Sigma)
&\longrightarrow\check H^2(N_\Sigma,H_\Sigma\to G_\Sigma),\\
[g]&\longmapsto[(g,1)]
\end{aligned}
\]
be the natural strict-locus map. Define the separate \(\{0,1\}\)-valued strictification obstruction by
\[
\begin{aligned}
o^{(2)}_\Sigma(g,h)=0
&\quad\Longleftrightarrow\quad q_\Sigma\in\operatorname{im}(\iota_\Sigma)\\
&\quad\Longleftrightarrow\quad [(g,h)]\text{ has a representative }(g^{\mathrm{str}},1)\\
&\hspace{7.2em}\text{with }g^{\mathrm{str}}_{ij}g^{\mathrm{str}}_{jk}=g^{\mathrm{str}}_{ik}.
\end{aligned}
\]
and set \(o^{(2)}_\Sigma(g,h)=1\) otherwise.

\textbf{Theorem 3.4 (non-central higher-defect strictification).} The higher associator defect of
\((g_{ij},h_{ijk})\) is removable iff \(o^{(2)}_\Sigma(g,h)=0\). In that case a higher-gauge change
gives \(h_{ijk}=1\) and a genuine \(G\)-valued edge \(1\)-cocycle. Strict endpoint-only transport
additionally requires that at least one allowed strict representative have trivial represented
holonomy. The full orbit \(q_\Sigma\) need not determine a unique ordinary \(H^1\) class.

\textbf{Proof sketch.} Removing the higher associator corresponds to \(h_{ijk}=1\) and
\(g_{ij}g_{jk}=g_{ik}\). Crossed-module coboundaries preserve the full orbit, so this
strictification exists exactly when that orbit meets the image of the ordinary \(G\)-valued
\(1\)-cocycle locus, which is the definition of \(o^{(2)}_\Sigma=0\). The map \(\iota_\Sigma\)
need not be injective: \(H\)-valued edge changes can relate strict representatives and alter their
ordinary holonomy through \(\partial\). Thus the invariant residual test is whether at least one
allowed strict representative has trivial represented holonomy, as in Theorem 3.4b. QED.

The central-defect case is the abelian truncation with \(H\) central and trivial action, which reduces to Section 3.3. A genuinely noncentral class is the point where the fixed-cutoff collar theorem upgrades from the ordinary-group package to the higher-gauge one below.

\textbf{Corollary 3.4a (Fixed-cutoff higher-gauge EC and transportability).} For a finite regulator chart on a connected cut \(\Sigma\), the genuinely noncentral branch admits a compact crossed-module change system
\[
\mathcal T_\Sigma
=
C^1(N_\Sigma,H_\Sigma)\rtimes C^0(N_\Sigma,G_\Sigma).
\]
The physical higher-gauge collar is
\[
\mathcal H_B^{2g}
=
(\widetilde{\mathcal H}_{B_L}\otimes \widetilde{\mathcal H}_{B_R})^{\mathcal T_\Sigma}
\cong
\bigoplus_\lambda
(\mathcal H_{b_L^\lambda}\otimes \mathcal H_{b_R^\lambda}),
\]
with
\[
Z(\mathcal A_{2g}(B))
=
\bigoplus_\lambda \mathbb C\,\mathbf 1_\lambda,
\]
and the full orbit \(q_\Sigma\) is invariant under local rechartings and classifies the fixed-cutoff
genuinely noncentral gluing datum. Its separate invariant property \(o^{(2)}_\Sigma=0\) holds iff
the higher associator is removable. The resulting strict representatives can have different
ordinary \(G_\Sigma\)-valued \(1\)-holonomies, so transportability is the separate existential test
of Theorem 3.4b rather than a uniquely defined holonomy attached to \(q_\Sigma\).

\textbf{Proof sketch.} Pair-overlap rechartings are inner, while triple-overlap associators strictify to compact crossed-module data. Finite-dimensional unitary \(\mathcal T_\Sigma\)-modules decompose semisimply, and Schur matching leaves only diagonal left/right blocks. Crossed-module cohomology classifies the higher defect; Theorem 3.4b separately tests the ordinary loop holonomy that remains after strictification. QED.

\textbf{Theorem 3.4b (TransportabilityFromOverlapGluing).} Fix a connected cut \(\Sigma\) and a finite regulator charting nerve \(N_\Sigma\). Transport of a collar charge is defined by the path composite in the overlap recharting groupoid: along \(p=(i_0\to\cdots\to i_m)\), the charge block is moved by \(U_p=U_{i_{m-1}i_m}\cdots U_{i_0i_1}\). This construction uses the overlap unitary transition system itself, not an added DHR transportability assumption.

After a central or crossed-module strictification, let \(U^{\mathrm{str}}\) denote a resulting
genuine edge \(1\)-cocycle. Write \(\mathcal C_\alpha\) for the full transported collar-charge block,
including its boundary carrier and multiplicity/intertwiner data, and define
\[
\operatorname{Hol}_{\alpha,U^{\mathrm{str}}}:\pi_1(|N_\Sigma|,i_*)\longrightarrow U(\mathcal C_\alpha),
\qquad
[\gamma]\longmapsto U^{\mathrm{str}}_\gamma\big|_{\mathcal C_\alpha}.
\]

On the ordinary branch, transport is strictly path-independent iff every closed overlap loop has trivial holonomy on the collar-sector block. On the central branch, elementary triangle moves accumulate the central cocycle \(z_{ijk}\). Strict path-independent transport exists iff \([z]_\Sigma=0\) and there is a central \(1\)-cochain strictification for which \(\operatorname{Hol}_{\alpha,U^{\mathrm{str}}}=1\). The first condition removes the projective triangle defect; the second removes the residual ordinary loop action. If \([z]_\Sigma=0\) but the second condition fails, the lifts form a genuine \(1\)-cocycle and transport is path dependent.

On the genuinely noncentral branch, path comparisons are \(H_\Sigma\)-valued \(2\)-morphisms in
the crossed module \(H_\Sigma\to G_\Sigma\). Strict ordinary transport exists iff
\(o^{(2)}_\Sigma=0\) and at least one crossed-module strictification gives a genuine
\(G_\Sigma\)-valued \(1\)-cocycle with trivial represented holonomy on \(\mathcal C_\alpha\). If
\(o^{(2)}_\Sigma\ne0\), the sector has higher transport rather than ordinary compact-group DR
transport. If \(o^{(2)}_\Sigma=0\) but every allowed strict representative has nontrivial
represented \(G_\Sigma\)-holonomy, the higher associator has strictified but ordinary path
dependence remains. The full orbit \(q_\Sigma=[(g,h)]\) need not determine a unique ordinary
\(H^1\) class and is not itself the strictification test. Thus
``zero obstruction'' below means the combined associator-strictification and trivial-holonomy
criterion, not vanishing of triangle or higher-associator data alone. QED.

\textbf{Triangle-free cycle check.} Let \(N_\Sigma=C_n\), \(n\ge4\), have no filled
\(2\)-simplices. On \(\mathcal C_\alpha=\mathbb C^2\), put identity transport on every oriented
cycle edge except one, where the sector action is the non-scalar unitary
\(V=\operatorname{diag}(-1,1)\), and use inverse transports on reversed edges. All central
\(2\)-data and higher associator data are vacuously trivial, so \([z]_\Sigma=0\) and
\(o^{(2)}_\Sigma=0\). Its ordered cycle holonomy is \(V\ne1\), which no central rephasing removes.
For a genuinely noncentral instance, take the crossed module
\[
\SU(2)\hookrightarrow\U(2)
\]
with conjugation action and the standard action on \(\mathcal C_\alpha\). The band
\(\U(2)/\SU(2)\cong\U(1)\) is detected by the determinant, and the cycle has band holonomy
\(\det V=-1\). An \(\SU(2)\)-valued edge change has determinant one and a vertex-frame change only
conjugates the based loop, so no allowed strict representative has trivial holonomy. The two routes
between the endpoints of the distinguished edge differ, and the combined criterion rejects both
the central and noncentral assignments.

\begin{center}\rule{0.5\linewidth}{0.5pt}\end{center}
\subsection{Modular Flow and Lorentz Kinematics}\label{4-modular-flow-and-lorentz-kinematics}

\subsubsection{Modular additivity in the Markov collar limit}\label{41-modular-additivity-in-the-markov-collar-limit}

Consider a collar tripartition \(A:B:D\) around a cap boundary, with the EC decomposition of Section 2.3 understood. Define, for a faithful reference state \(\omega\),
\[
\Delta K(\omega)
:=
K_{ABD}(\omega)-K_{AB}(\omega)-K_{BD}(\omega)+K_B(\omega).
\]

Exact modular additivity is not a consequence of EC alone. It is available only on the EC-aligned exact Markov set, or along a controlled fixed-cutoff family that approaches that set on one fixed collar model. Four quantities must therefore be kept separate:
\begin{enumerate}
\item the raw collar conditional mutual information \(I(A:D\mid B)\);
\item the constructive Fawzi--Renner comparison error
\[
r_{\mathrm{FR}}(\varepsilon):=
2\sqrt{1-e^{-\varepsilon}}
\le
2\sqrt{\varepsilon};
\]
\item the fixed-collar exact-Markov replacement modulus
\[
\delta^{\mathrm M}_{A:B:D}(\varepsilon)
:=
\sup\left\{
\inf_{\sigma\in\mathfrak M_{A:B:D}}\|\rho-\sigma\|_1:
I(A:D\mid B)_\rho\le \varepsilon
\right\},
\]
where
\[
\mathfrak M_{A:B:D}
:=
\left\{
\sigma_{ABD}: I(A:D\mid B)_\sigma=0
\right\}.
\]
\item the EC-aligned subset and a family-specific aligned replacement modulus,
\[
\mathfrak M^{\mathrm{EC}}_{A:B:D}
:=
\{\sigma\in\mathfrak M_{A:B:D}:\sigma\text{ satisfies MSA on the preselected EC split}\},
\]
\[
\delta^{\mathrm{M,EC}}_{\mathcal F;A:B:D}(\varepsilon)
:=
\sup\left\{
\inf_{\sigma\in\mathfrak M^{\mathrm{EC}}_{A:B:D}}\|\rho-\sigma\|_1:
\rho\in\mathcal F,\ I(A:D\mid B)_\rho\le\varepsilon
\right\}
\]
for the declared controlled collar family \(\mathcal F\). Convergence of the
generic modulus in item 3 does not imply convergence of this aligned modulus;
\(\delta^{\mathrm{M,EC}}_{\mathcal F}\to0\) is a separate family hypothesis.
\end{enumerate}

\textbf{Lemma 4.1a (EC-aligned exact Markov implies exact additivity up to a central term).} If \(I(A:D \mid B)_\omega = 0\) and \(\omega\) satisfies the Markov-split alignment hypothesis of Section 2.3, then the state takes the EC-aligned HJPW block form
\[
\omega_{ABD}
=
\bigoplus_{\alpha}
p_{\alpha}\,
\omega^{(\alpha)}_{A b_L^{\alpha}}
\otimes
\omega^{(\alpha)}_{b_R^{\alpha} D},
\]
and \(\Delta K(\omega)\) is central. On the canonical HJPW block model one may take
\[
\Delta K(\omega)=0.
\]
Equivalently, there exists a central operator \(K_{\partial,ABD}(\omega)\in Z(\mathcal A(B))\) such that
\[
K_{ABD}(\omega)
=
K_{AB}(\omega)+K_{BD}(\omega)-K_B(\omega)+K_{\partial,ABD}(\omega).
\]

\textbf{Proof.} By alignment, the HJPW blocks are the EC blocks. On each block the modular Hamiltonians of \(ABD\), \(AB\), \(BD\), and \(B\) are the logarithms of tensor-product states with the same classical block weight \(p_\alpha\). The blockwise logarithms therefore cancel exactly in the combination \(K_{ABD}-K_{AB}-K_{BD}+K_B\). If one keeps the blockwise endpoint-label bookkeeping explicit rather than absorbing it into the canonical block identification, the remainder is a direct sum of block constants, hence central in \(Z(\mathcal A(B))\). Without alignment, exact Markovity cancels the four logarithms on the state's own HJPW blocks, but the resulting block constants need not lie in \(Z(\mathcal A(B))\), so the central form used here is unavailable. QED.

\textbf{Proposition 4.1b (Controlled exact-Markov replacement on a fixed collar).}
Because the state space is compact and conditional mutual information is continuous in finite dimension,
\[
\delta^{\mathrm M}_{A:B:D}(\varepsilon)\longrightarrow 0
\qquad
(\varepsilon\downarrow0).
\]
Thus small collar CMI converges to the exact HJPW normal form only in this controlled fixed-cutoff sense; the manuscript does not claim a dimension-free one-shot trace-norm bound from small \(I(A:D\mid B)\) directly to an exact Markov state. The limit points are exact Markov states in the CMI-zero sense; identifying their factorization with the preselected EC factors is the Markov-split alignment hypothesis of Section 2.3, which is carried separately and is not produced by small CMI (the Bell-pair counterexample of Section 2.3 has zero CMI at fixed distance from every EC-aligned state).

\textbf{Theorem 4.1b' (Finite-collar Markov replacement stability).} On one fixed faithful
finite-dimensional collar model, the qualitative modulus above can be made collar-local and
quantitative. If all relevant marginals stay above a chosen floor \(\lambda_\ast>0\), then there are
constants \(C_{A:B:D,\lambda_\ast}>0\) and \(\theta_{A:B:D,\lambda_\ast}>0\), depending only on that
collar model and floor, such that
\[
d_{\mathrm M}(\rho)
\le
C_{A:B:D,\lambda_\ast}\,
I(A:D\mid B)_\rho^{\theta_{A:B:D,\lambda_\ast}}.
\]
This is a fixed-collar Lojasiewicz-type rate for the analytic function
\(I(A:D\mid B)\) on the faithful state manifold. It is not a dimension-free stability theorem for
arbitrary tripartite quantum systems.

\textbf{Corollary 4.1c (Carried EC-aligned collar defect operator).} Let \(\omega_\varepsilon\in\mathcal F\) satisfy \(I(A:D\mid B)_{\omega_\varepsilon}\le \varepsilon\), assume the declared family has an EC-aligned replacement modulus tending to zero, and choose \(\sigma_\varepsilon\in\mathfrak M^{\mathrm{EC}}_{A:B:D}\) with
\[
\|\omega_\varepsilon-\sigma_\varepsilon\|_1
\le
\delta^{\mathrm{M,EC}}_{\mathcal F;A:B:D}(\varepsilon).
\]
Define the carried defect operator
\[
\mathfrak D_{A:B:D}(\omega_\varepsilon,\sigma_\varepsilon)
:=
\Delta K(\omega_\varepsilon)-\Delta K(\sigma_\varepsilon).
\]
Then
\[
K_{ABD}(\omega_\varepsilon)
=
K_{AB}(\omega_\varepsilon)+K_{BD}(\omega_\varepsilon)-K_B(\omega_\varepsilon)
+K_{\partial,ABD}(\sigma_\varepsilon)
+\mathfrak D_{A:B:D}(\omega_\varepsilon,\sigma_\varepsilon),
\]
where \(K_{\partial,ABD}(\sigma_\varepsilon)=\Delta K(\sigma_\varepsilon)\) is central. For every bounded observable \(X\) supported on \(A\cup B\),
\[
\left|
\operatorname{Tr}\!\left[X(\omega_\varepsilon-\sigma_\varepsilon)\right]
\right|
\le
\|X\|_\infty\,\delta^{\mathrm{M,EC}}_{\mathcal F;A:B:D}(\varepsilon).
\]
If the relevant marginals of \(\omega_\varepsilon\) and \(\sigma_\varepsilon\) are uniformly faithful with lower spectral bound \(\lambda_\ast>0\), then
\[
\|\mathfrak D_{A:B:D}(\omega_\varepsilon,\sigma_\varepsilon)\|_\infty
\le
4\lambda_\ast^{-1}\,\delta^{\mathrm{M,EC}}_{\mathcal F;A:B:D}(\varepsilon).
\]

Independently, Fawzi--Renner recovery guarantees a channel acting on the collar and a comparison state \(\omega^{\mathrm{rec}}_\varepsilon\) with squared Uhlmann fidelity at least \(e^{-\varepsilon}\), for CMI \(\varepsilon\) in nats, and hence
\[
\|\omega_\varepsilon-\omega^{\mathrm{rec}}_\varepsilon\|_1
\le
r_{\mathrm{FR}}(\varepsilon),
\qquad
r_{\mathrm{FR}}(\varepsilon):=2\sqrt{1-e^{-\varepsilon}}\le 2\sqrt{\varepsilon}.
\]
This \(O(\varepsilon^{1/2})\) term bounds observable errors for that recovered state; it is not a guarantee for an arbitrary unrotated Petz map. The generic modulus \(\delta^{\mathrm M}_{A:B:D}(\varepsilon)\) justifies replacement only by some exact Markov state on its state-dependent HJPW split. The separately assumed \(\delta^{\mathrm{M,EC}}_{\mathcal F;A:B:D}(\varepsilon)\) is the quantity that justifies the EC-aligned normal form used in later geometric modular arguments.

Accordingly, the manuscript's later exact collar formulas take one of two forms:
\begin{enumerate}
\item literal exactness when the reference state is exact Markov and EC-aligned on the relevant collar; or
\item a controlled collar family for which \(\delta^{\mathrm{M,EC}}_{\mathcal F;A:B:D}(\varepsilon_\delta)\to0\), with the constructive Fawzi--Renner remainder and the carried operator defect \(\mathfrak D_{A:B:D}\) kept explicit at finite stage.
\end{enumerate}

\textbf{Theorem 4.1d (Finite-stage modular-defect propagation and dimension quarantine).}
For any fixed downstream branch calculation that uses finitely many collar or strip
modular-additivity identities, replacing each exact identity by its controlled finite-stage form
changes every bounded downstream modular observable \(O\) by at most
\[
\mathcal P_O\!\left(
\{r_{\mathrm{FR}}(\varepsilon_j)\},
\{\delta^{\mathrm{M,EC}}_{\mathcal F_j;j}(\varepsilon_j)\},
\{\eta^{\mathrm{reg}}_j\}
\right),
\]
where \(\mathcal P_O\to0\) as all listed errors vanish. The polynomial-continuity modulus depends
only on the declared fixed collar models, support-visible dimensions, bounded observable class,
faithful floors or regularization schedules, and bounded modular-parameter intervals. No BW/Lorentz,
null-modular, or Einstein scaling step uses a dimension-free trace-norm stability theorem from
small CMI directly to exact Markov normal form.
\subsubsection{\texorpdfstring{Theorem: \(\mathrm{BW}_{S^2}\) on the extracted geometric subnet}{4.2 Theorem: \textbackslash mathrm\{BW\}\_\{S\^{}2\} on the extracted geometric subnet}}\label{42-theorem-mathrmbw_s2-from-markov-locality-symmetry-regularity}

The collar analysis proves a fixed-cutoff statement on the finite type-I regulator net: the reduced cap state has a literal density matrix, its modular Hamiltonian exists, and its nonadditive part is confined to a shrinking collar up to carried errors. The Lorentz claim is therefore not a fixed-cutoff matrix-algebra identity and not the slogan ``finite cells imply Lorentz invariance.'' Its target is the support-visible refinement-limit geometric cap pair \((\mathcal A_\infty^{\mathrm{geo,sv}}(C),\omega_\infty^{\mathrm{geo},C})\). The branch theorem is the conditional implication
\[
\mathsf{FiniteCapBWCertificate}
\; + \; \mathsf{ModularAlgebraStateData}
\Longrightarrow
\sigma_t^{\omega_\infty^{\mathrm{geo},C}}
=
\alpha_{\lambda_{\widehat C}(2\pi t)}.
\]
The finite cap-normal certificate consists of cap-normal density and nondegeneracy, BW framing,
support-order faithfulness, geometric support-flow group and continuity control, held-out oriented
cross-ratio convergence, independently normalized geometric \(2\pi\)-KMS comparison, and
wrong-normalization separation. The independent multiresolution algebra-state comparison package supplies the
algebra-state reference tower, common comparison maps, cap-family-uniform regularized modular
transport, mixed-GNS convergence, and modular support covariance. The support-visible theorem below
uses both inputs; it does not require a type-I continuum algebra or a full-algebra unregularized
common spectral floor. Bare finite consensus implies neither input.

Fix a cap \(C\subset S^2\) and a shrinking collar family \((A_\delta,B_\delta,D_\delta)\) around \(\partial C\). Write
\[
\varepsilon_\delta:=I(A_\delta:D_\delta\mid B_\delta)_\omega,
\qquad
r_{\mathrm{FR}}(\varepsilon_\delta):=
2\sqrt{1-e^{-\varepsilon_\delta}}
\le 2\sqrt{\varepsilon_\delta},
\]
and, on each fixed faithful collar model,
\[
\eta_\delta^{\mathrm M}
:=
4\lambda_\ast^{-1}\,
\delta^{\mathrm{M,EC}}_{\mathcal F_\delta;A_\delta:B_\delta:D_\delta}(\varepsilon_\delta),
\]
where \(\lambda_\ast>0\) is the lower spectral bound for comparing modular Hamiltonians. All exact collar formulas in this subsection are therefore to be read either literally at EC-aligned exact Markovity or along a controlled collar family satisfying
\[
\delta^{\mathrm{M,EC}}_{\mathcal F_\delta;A_\delta:B_\delta:D_\delta}(\varepsilon_\delta)\to 0
\qquad(\delta\downarrow 0),
\]
with \(r_{\mathrm{FR}}(\varepsilon_\delta)\) and \(\eta_\delta^{\mathrm M}\) carried explicitly.

For each cap \(C\), let \(\lambda_C(s)\subset \mathrm{Conf}^+(S^2)\) denote the standard cap-preserving conformal one-parameter subgroup, normalized so that the null blow-up near a smooth cut acts by \(v\mapsto e^{-s}v\), and let \(\alpha_{\lambda_C(s)}\) denote the induced automorphism of the scaling-limit cap net.

\textbf{Theorem 4.2 (Finite \(\mathrm{BW}_{S^2}\) cap-net scaling).} For a nondegenerate BW-framed cap
\(\widehat C=(C,n_C,p_C^-,p_C^+)\), consider an OPH realization
on an observer-supporting refinement branch that satisfies premises A1--A3 and carries the named
fixed-cutoff regulator, collar, consensus, and optimizer-pushforward interfaces. Assume the finite
cap-normal BW certificate above and, independently, the complete multiresolution reference tower
and algebra-state comparison package of Theorems 2.6c--2.6f on the same branch. The support-visible
geometric cap net then admits a weak-\(*\)/GNS scaling-limit cap pair
\((\mathcal A_\infty^{\mathrm{geo,sv}}(C),\omega_\infty^{\mathrm{geo},C})\), whose modular automorphism
group is
\[
\sigma_t^{\omega_\infty^{\mathrm{geo},C}}
=
\alpha_{\lambda_{\widehat C}(2\pi t)}.
\]
At finite regulator stage the nonadditive part of
\[
K_C^{(\delta)}:=-\log \rho_C^{(\delta)}
\]
is confined to the shrinking collar up to the carried errors \(r_{\mathrm{FR}}(\varepsilon_\delta)\), the EC-aligned fixed-collar replacement bound \(\eta_\delta^{\mathrm M}\), and the regularized support-visible modular transport bound. The support-visible modular limit is independent of the EC-aligned exact-Markov replacement sequence after projective compatibility and GNS quotienting, and support-readable modular covariance turns the limiting modular group into a cap-local support map before the certificate's framed-cap and oriented cross-ratio rigidity clauses are applied. No separate cap-isotropy/\(\mathrm{SO}(2)\)-equivariance selector, finite-cell Lorentz-invariance premise, or full-algebra unregularized common floor is used. Record/pointer and interface-inert auxiliary registers are outside the extracted geometric subnet. If the scaling-limit cap algebra is type I, the same automorphism identity may be represented by
\[
K_C=2\pi B_C+Z_C,\qquad Z_C\in Z(\mathcal A(C))_{\mathrm{sa}};
\]
otherwise the automorphism identity is the full statement and the action is outer. Consequently,
under the paired same-branch finite cap-normal and complete multiresolution algebra-state comparison inputs, the
transported regularized modular action on every fixed support-visible geometric local algebra
converges uniformly on compact modular-parameter intervals to the geometric cap-dilation action.
Without both inputs, this theorem makes no BW or finite-to-continuum modular claim. Even with
them, it does not assert compact-time convergence of finite-regulator modular groups beyond the
declared support-visible certificate.

\textbf{Proof.} Markov locality localizes the modular defect to the shrinking collar, with carried finite-stage errors \(r_{\mathrm{FR}}(\varepsilon_\delta)\) and \(\eta_\delta^{\mathrm M}\). Proposition 2.6a shows why a full-algebra unregularized floor is unavailable in general. Theorem 2.6b supplies the matrix-element replacement on fixed collar models, and Theorem 2.6f supplies the stronger compact-time group estimate only for transported regularized local actions with the stated reference-tower certificate. Projective replacement compatibility and the support-visible convergence theorem ensure that the emitted modular limit is independent of the chosen EC-aligned exact-Markov comparison family. Support-visible cap-pair extraction on the local GNS support quotient uses weak-\(*\) compactness of finite-stage state spaces and the consensus refinement mechanism to supply local limit states, then applies GNS construction to obtain the scaling-limit cap pair. Support-readable modular covariance reads the modular group as a cap-local support map. The finite cap-normal certificate supplies support-order faithfulness, BW framing, held-out oriented cross-ratio rigidity, and compact-time equicontinuity, so the support flow is the framed cap-preserving conformal flow up to normalization:
\[
\sigma_t^{\omega_\infty^{\mathrm{geo},C}}
=
\alpha_{\lambda_C(\kappa_C t)}
\]
for some \(\kappa_C>0\), with the certificate-selected frame understood in \(\lambda_C\). The
certificate then passes the independently normalized geometric \(2\pi\)-KMS condition to
\(s\mapsto\alpha_{\lambda_{\widehat C}(s)}\). KMS uniqueness, together with wrong-\(\beta\)
separation and nontriviality, fixes \(\kappa_C=2\pi\). Thus
\[
\sigma_t^{\omega_\infty^{\mathrm{geo},C}}
=
\alpha_{\lambda_{\widehat C}(2\pi t)}.
\]
If the limit algebra is type I, this automorphism identity is represented by
\(K_C=2\pi B_C+Z_C\) with \(Z_C\) central; otherwise the automorphism identity is the full
statement and the action is outer. QED.

\textbf{Definition 4.2a (BW-branch observer-relative modular ordering).}
On the branch satisfying the hypotheses of Theorem 4.2, the modular automorphism
parameter \(t\) of the extracted cap pair
\((\mathcal A_\infty^{\mathrm{geo,sv}}(C),\omega_\infty^{\mathrm{geo},C})\) supplies a
dimensionless ordering for that observer's accessible algebra-state pair. Physical time requires
an observer-readable transition, event correspondence, and calibrated clock instrument. No claim
about arbitrary operational clocks, global time, or the full problem of time follows from the
geometric modular-flow theorem alone.

\textbf{Scope note.} The theorem fixes the \(2\pi\) normalization and the carried refinement limit on the support-visible extracted geometric cap pair under the paired finite cap-normal and independently complete same-branch multiresolution algebra-state comparison inputs. The UV-side scaffold is the realized transported geometric cap-local system, regularized support-visible modular transport on fixed local collars, support-visible cap-pair extraction on the local GNS support quotient, support-readable modular covariance, BW framing, oriented cross-ratio rigidity, and geometric KMS normalization. No separate cap-isotropy input is used, and the conformal support-map clause is the explicit scaling-limit branch condition rather than a finite-regulator Lorentz premise. Compact-time convergence of finite-regulator modular groups is available only through the transported-local certificate of Theorem 2.6f; a black-box conformal-net BW reconstruction would be an additional validation route, not a hidden premise.
\subsubsection{\texorpdfstring{Theorem: cap-normal \(H^3\) reconstruction}{4.3 Theorem: cap-normal H3 reconstruction}}\label{43-theorem-mathrmbw_s2-implies-lorentz-kinematics}

\textbf{Theorem 4.3 (Cap-normal \(H^3\) reconstruction, imported from the spacetime and Einstein derivation).}
On the support-visible round-cap BW branch, write the celestial null section as
\[
q(\Omega)=(1,\Omega),\qquad \Omega\in S^2.
\]
For a nondegenerate oriented round cap
\[
C(\mathbf c,\alpha)=\{\Omega:\mathbf c\cdot\Omega\geq\cos\alpha\},
\qquad 0<\alpha<\pi,
\]
define
\[
n_C=(\cot\alpha,\csc\alpha\,\mathbf c).
\]
Then, for the time-first Lorentz form \(\eta=-dt^2+d\mathbf x^2\),
\[
\eta(n_C,n_C)=1,
\qquad
\eta(n_C,q(\Omega))
=
\frac{\mathbf c\cdot\Omega-\cos\alpha}{\sin\alpha}.
\]
Thus \(\eta(n_C,q(\Omega))=0\) exactly on \(\partial C\), with positive
sign in the cap interior and negative sign outside. For
\(g\in\mathrm{Conf}^+(S^2)\) with Lorentz representative \(\Lambda_g\),
\[
\Lambda_gq(\Omega)=\omega_g(\Omega)q(g\Omega),
\qquad
n_{gC}=\Lambda_gn_C.
\]
Consequently
\[
\mathrm{Conf}^+(S^2)\cong \mathrm{PSL}(2,\mathbb C)\cong \mathrm{SO}^+(3,1),
\]
oriented round caps form the de Sitter normal space
\[
\operatorname{Cap}^{\mathrm{or}}_{\mathrm{round}}(S^2)
\cong
dS_3^{\mathrm{cap}}
\cong
\mathrm{SO}^+(3,1)/\mathrm{SO}^+(2,1),
\]
and future unit timelike observer frames form the distinct hyperbolic space
\[
H^3\simeq \mathrm{SO}^+(3,1)/\mathrm{SO}(3),
\qquad
\dim H^3=3.
\]
Each cap normal determines the geodesic half-space
\[
\mathcal H_C^+=\{u\in H^3:\eta(n_C,u)\geq0\},
\]
not a preferred observer point.

\textbf{Proof.} This is the cap-normal reconstruction theorem in The spacetime and Einstein paper, which supplies the projective-null-cone model, the cap-normal bijection, Lorentz equivariance, the \(H^3\) observer-frame hyperboloid, the ideal-boundary identification, and the cap/half-space duality. The BW theorem supplies the cap-preserving conformal subgroup and its \(2\pi\) modular normalization; the cap-normal theorem supplies the explicit projective Lorentz representation and does not rederive BW from bare finite consensus. QED.

\textbf{Corollary 4.3a (Frame-space boundary).} The theorem proves the exact canonical Lorentz/\(H^3\) observer-frame hyperboloid on the stated branch. Its points are future unit timelike frames, not event positions. It does not populate a physical spacetime, construct a chart-blind neutral bulk, derive a physical curvature radius \(R_H\), produce stress-energy, or close the Einstein branch-entry umbrella. Record-conditioned cap responses can estimate frame data under the certificate imported next; event location, object families, neutral bulk, scale, stress, and Einstein branch entry are separate receipts. Theorems 4.3c--4.3h below import The spacetime and Einstein paper's geometry-producer, null-net-standardness, source-causal precursor and continuum, stress, entropy-bridge, and composed branch-entry packets. Their composition is conditional on one source-derived common-domain tower with certified cofinal tails, universal coupling, a vacuum reference, and independent physical scale readouts. No source-derived tower carrying that complete packet is constructed here.

\textbf{Theorem 4.3b (Conditional record-conditioned \(H^3\) frame estimate).} Fix a quotient-visible record token \(i\) and clock slice \(t\). On the frame space of Corollary 4.3, suppose its record-conditioned modular cap responses
\[
R_i(C,t,O)=\omega_{i,O}\!\left(\sigma_t^{C,O}(M_{C,0,O})\right),
\qquad
\omega_{i,O}(A)=\frac{\omega_O(P_{i,O}AP_{i,O})}{\omega_O(P_{i,O})},
\]
admit the calibrated frame-local factorization
\[
y_{i,j}=F_j(X_i(t))+e_{i,j},\qquad
F_j(X)=\psi_j\!\left(\frac{\eta(X,n_j)}{R_H}\right),
\]
on a compact frame region \(\Omega\subset H^3_{R_H}\). Assume the finite response map has quantitative observability
\[
\alpha \bar d(X,Y)\le |F(X)-F(Y)|_W\le L\bar d(X,Y),\qquad \alpha>0,
\]
and the total error obeys \(|e_i|_W\le\sigma_i\). If \(\mathcal N_\varepsilon\) is an \(\varepsilon\)-net of \(\Omega\) and \(\widehat X_i(t)\) is a residual minimizer up to tolerance \(\tau_i\), then the conditioned feasible frame support obeys the conservative enclosure
\[
S_i(t)\subseteq B_H\!\left(\widehat X_i(t),
R_H\left[\frac L\alpha\varepsilon+\frac2\alpha\sigma_i+\frac1\alpha\tau_i\right]\right).
\]
The theorem does not assert that every point of this output ball fits the
data.
Exact continuous data have a unique minimizer. Finite noisy data produce the displayed frame ball; a unique finite frame value is licensed only when the certified residual gap \(\Delta_{\mathrm{loc},i}(t)>0\). If that gap fails, the fail-closed protocol returns \(\mathrm{AMBIGUOUS}\), meaning that no unique frame value is certified; it is not a proof that the mathematical data have multiple minimizers. This theorem does not locate an event, populate the event base, resolve unlabeled mixtures, identify particle species, produce stress-energy, or establish chart-blind neutral bulk.

\textbf{Theorem 4.3c (Quotient-intrinsic geometry producer and dimension-selection boundary).} On a refinement tower of repaired quotient normal forms, form at each level the support-visible incidence complex: vertices are the quotient-visible patch classes and simplices include every jointly supported patch set, including all higher simplices. A filled tetrahedron is not its spherical boundary. Each level is computable from its normal form and is invariant under repair schedule and gauge. Refinement generally changes the literal complex. Topology preservation requires a certified subdivision or homeomorphism: a complete midpoint subdivision has such a realization, whereas its endpoint-valued simplicial projection is only homotopic to it. Simplicial maps or CPTP transport alone do not establish topology preservation. If this compatible tower carries the computable receipts (i) spherical incidence (connected closed orientable combinatorial surface with Euler characteristic \(2\)), (ii) disk/mesh nondegeneracy, (iii) modular cross-ratio Cauchy convergence, and (iv) independently normalized geometric \(2\pi\)-KMS comparison with wrong-normalization separation, then the refinement limit produces the round \(S^2\), its conformal structure, the oriented round-cap family with cap normals, and every clause of the finite cap-normal BW certificate, with quantitative error envelopes. Conversely, six finite repair systems with one rewrite signature (three patches in the seed record, one repair step) have incidence complexes realizing \(S^2\), \(T^2\), the real projective plane, the Klein bottle, the \(2\)-skeleton of \(\partial\Delta^4\), and a wedge of two spheres, and the scalar KMS receipt admits any modular temperature \(\beta\): repair confluence alone does not select topology, orientability, dimension, or modular normalization. On a connected closed surface the orientability clause of (i) follows from Euler characteristic \(2\), so the two nonorientable systems also fail (i) at that clause. Orientation framing enters as a separately typed receipt: the two coherent orientations of the \(S^2\) system share every confluence-level datum, and a mirror pair of embeddings carries complex-conjugate cross-ratios. Finite cap reconstruction also requires a resolved spacelike boundary-circle plane and an inside witness or equivalent oriented boundary data to distinguish the two sides, including major caps. Cross-ratios determine only conformal data up to M\"obius transformations; a declared observer frame is needed for angular radii, and a physical scale remains separate. The MGNS state enters as a separately typed receipt, and no countermodel artifact is exhibited for it. The screen \(S^2\) is therefore \emph{produced on the receipt branch}; the receipt set is one sufficient certificate, but no semantic-minimality theorem or full-tower countermodel matrix is supplied. Producers that reach an oriented \(S^2\) support limit by another computable function of the normal forms, such as the source-derived incidence nerve of the carrier-support bridge, are untouched by these countermodels. (Spacetime and Einstein paper, geometry-producer packet; machine receipts accompany the public geometry code.)

\textbf{Theorem 4.3d (Null-net standardness, derived translations, and the four-translation assembly).} On the producer branch, the null half-line/interval blow-up net at a produced cap boundary point lives on one common GNS space, and: (i) the reference vector is cyclic and separating for every half-line used, given faithful finite-stage states on the interior MaxEnt branch, the mixed-GNS Cauchy clause, and a cofinal cyclicity receipt uniform over every half-line in the construction; boundary MaxEnt solutions require a support-restricted treatment and do not supply faithfulness by themselves; (ii) the produced geometric cap flow derives the half-sided covariance relation without assuming translation unitaries, but the Borchers--Wiesbrock step additionally consumes standardness for every nested pair---cyclicity and separation for both half-line algebras and cyclicity for their relative commutant; interval cyclicity is therefore a registered standard-relative-commutant receipt, not a consequence of one half-line receipt; (iii) weak additivity of the null net is a separate cofinal receipt used for the generated net, rather than a consequence of finite-stage faithfulness; (iv) on the exact-Markov/central-interface branch, with a specified endpoint-to-cell map and a uniform local norm/site-density bound, the reduced interval modular Hamiltonians split into local terms plus uniformly bounded interface blocks (endpoint-Lipschitz), while an explicit four-qubit nearest-neighbour Gibbs witness shows finite-range locality alone does \emph{not} imply this (machine-verified); (v) on the standard half-sided-inclusion and modular-intersection receipt branches, the derived positive null translations and modular dilations generate M\"obius covariance and the projective bounded-interval action; and (vi) the per-direction positive generators commute, transform with the null-vector weight under the produced Lorentz action, satisfy the dependent-family linearity relation, and assemble into one four-parameter translation representation with joint spectrum in the closed future cone: a positive-energy Poincar\'e representation. (Spacetime and Einstein paper, null-net standardness packet.)

\textbf{Theorem 4.3e (Source-derived finite \(1+3\) precursor and effective-continuum criterion).} An authenticated semantic-event log supplies its finite event carrier and strict generated read-from order. Its canonical source height is computed from that parent relation: roots have height zero and every nonroot has one plus the maximum height of its direct parents. Lean proves strict increase on parent edges and generated ancestry, bounds every authenticated direct-parent chain ending at an event by that event's height, and constructs an attaining chain. Thus source height is exactly the longest authenticated-parent-chain length. Let \(V_{\mathrm{src}}\) be the exact rank-three positive source Gram quotient. Then \(W_{\mathrm{src}}=\mathbb R\oplus V_{\mathrm{src}}\) has real dimension four, and \(Q(t,x)=t^2-g_{\mathrm{src}}(x,x)\) has inertia \((1,3)\). Every source-unit direction gives the future-null vector \((1,u)\), realizing the exact equivalences \(\mathsf{SourceUnitDirection}\simeq S^2\simeq\mathsf{FutureNullRay}\). For a supplied spatial readback \(x(e)\) and positive height scale \(\tau\), the placement \(F(e)=(\tau h_L(e),x(e))\) carries authenticated parenthood, and hence generated ancestry, into the future cone whenever every parent displacement obeys the source-metric speed bound. If every future-cone comparison between placed events is supported by reflexive generated ancestry, order and cone agree exactly. Lean derives this support from two concrete finite checks: spatial readback separates each equal-height layer, and every increasing-height pair unsupported by ancestry obeys the strict spacelike inequality. Antisymmetry then derives injectivity, every source interval agrees with the corresponding cone interval on the placed image, and composition with the source-carrier equivalence gives a separated one-chart finite causal-chart interface. This global finite chart is not an open manifold atlas. Thus the ambient target-carrier dimension and Lorentz signature are constructed without a freely declared rank-four chart, fitted cone, or supplied placement rank; the intrinsic poset dimension, spatial placement, speed/converse conditions, physical interpretation, and clock remain separate data. The exact-embedding route to a smooth effective manifold uses one cofinal source-selected family carrying physical event/link identification, exact two-way cone/order agreement, refinement compatibility, dense and isotropic source-\(S^2\) link directions, a calibrated \(\#I/\rho\to\operatorname{Vol}(I)\) law, independent dimension/manifoldlikeness and height/profile tests, stable thickened-antichain topology or homology, and a specified distinguishing Lorentzian limit with uniqueness control. On that certificate, causal order fixes conformal geometry and the calibrated measure fixes scale. Smooth Einstein promotion additionally requires either same-family tensor-curvature reconstruction converging to the Einstein tensor or the independently premised continuum small-ball/null-balance identification, together with physical stress, Ward/Bianchi, coupling, and remainder data on the same family. Scalar-curvature convergence alone is diagnostic. No physical instance of this certificate is supplied. (Spacetime and Einstein paper; Lean theorems and independently verified simulation evidence.)

\textbf{A source-record order and count limit.} In the continuous rank-three
carrier, conservative Fibonacci records give $q^3$ distinct spatial sites.
Repeat this population at every layer and let each event read all preceding-layer
sites inside a radius $a_q$, including itself. With spatial fill distance
$h_q=o(a_q)$, $a_q\to0$, and model duration $a_q/c$, actual read paths
approach the speed-$c$ causal cone. Equal-volume cell assignments with vanishing displacement error make
normalized raw counts converge to $dt\,d^3x$; the ordering fraction in
interior timelike diamonds tends to $1/10$. This is a flat effective
spacetime theorem for a specified population, read law and model clock
(Spacetime and Einstein paper). It supplies neither their native dynamical
selection nor their identification with a matter action or laboratory clock.

The same histories admit a retrospective clock. For intervals approaching
fixed interior timelike diamonds, the fourth root of their inclusive count
ratio converges to their proper-time ratio.
A reference interval fixes the unit, and complete ancestry access supplies
the count. This readout needs no oscillator or numerical timestamps.
Operational covariance gives a separate cone criterion: a nonzero closed
convex pointed cone invariant under the source icosahedral rotations and
all boosts along one axis must be one of the two Lorentz cones. A time
orientation selects the future cone. The covariance must act on possible
influences and readouts. Algebraic direction maps alone do not establish
that operational premise\cite{ophEinstein}.

For Einstein promotion there are two distinct sufficient routes at this step:
an explicit same-family tensor-curvature reconstruction converging to the
smooth Einstein tensor, or a separately premised continuum
small-ball/null-balance identification. Convergence of a scalar-curvature
estimator is an independent diagnostic and is insufficient by itself.

Every finite semantic log also admits an explicit auxiliary separated one-way
placement in the source (1+3) carrier. Enumerating events along one source
axis and taking height scale one larger than the event count makes the event
map injective, sends generated precedence into the future cone, and makes
distinct same-height events spacelike. The enumeration is arbitrary and
generally introduces cone comparisons unsupported by ancestry; it is not an
order-reflecting, source-selected, or physical placement.

The provenance grammar has the complementary exact expressivity theorem.
Every finite decidable irreflexive transitive relation $R$ compiles into an
abstract authenticated semantic log whose direct parent relation and generated
strict order are both exactly $R$; predecessor cardinality supplies the rank.
The generic parent relation is all of $R$ rather than its Hasse reduction alone.
The relation is supplied to the compiler, and the per-event snapshots are not
proved to arise from one threaded executor. This is a universal finite-poset
representation result, not a dynamics, manifoldlikeness, or continuum theorem.

An exact four-event Boolean diamond inhabits the order-faithful finite interface:
its parent edges are null and its two independent branches are spacelike.
This is a machine-checked non-chain causal-set-like witness, not a physical
continuum or an empirical dimension measurement.

The finite provenance order is causal-set-like rather than a physical causal
set. Lean packages reflexive authenticated ancestry as a partial order and
proves every interval finite at each finite cutoff, exactly realizing the
abstract order-theoretic causal-set axioms. Standard causal set theory starts from a locally finite order interpreted
as proto-causality \cite{bombelli1987,surya2019}. The exact-embedding continuum route
requires an order embedding whose image has the approximately Poisson count
density needed for a causal-set faithful embedding, together with
manifoldlikeness and dimension
\cite{bombellimeyer1989,brightwellgregory1991,majorrideoutsurya2007}. Under
the relevant distinguishing hypotheses, continuum causal order fixes the
conformal geometry \cite{hawkingkingmccarthy1976,malament1977}; the volume
datum fixes the conformal factor. The Lorentz-invariant causal-set
approximation uses Poisson sprinkling \cite{bombellihensonsorkin2009}; the OPH
commit log is not a Poisson sprinkling, and its event count is not a spacetime
volume. The finite source-causal theorem supplies the (1+3) precursor, not a
continuum geometry. It derives neither a physical order embedding nor the
count-to-volume density law that together constitute a causal-set faithful embedding.
It also establishes neither continuum existence nor the causal-set
Hauptvermutung from the finite OPH order.
The finite causal/Lorentz precursor instead constructs the ambient carrier
from an independent real axis and the exact rank-three source carrier. A
spatial event readback, canonical source height, and positive
\texttt{timeScale} then give a concrete
placement; a local speed bound proves one-way cone compatibility, while
equal-height spatial separation plus strict spacelikeness of unsupported
increasing-height pairs supply converse cone support and exact two-way
agreement. Antisymmetry then derives injectivity, while the two-way
equivalence preserves every source interval exactly on the placed image.
The event-count measure is distinct from public capacity \(N\). The finite join
does not constitute spacetime. Physical event/link identification, refinement,
manifoldlikeness, topology, independent dimension tests, count--volume
calibration, and uniqueness are required before the continuum can be read as
emergent rather than fundamental.
The exact source-direction bridge gives
\(\mathsf{SourceUnitDirection}\simeq S^2\simeq
\mathsf{FutureNullRay}\) for the rank-three source Gram quotient and fixes the
finite (1+3) carrier and cone after adjoining an independent real axis.
Positive \texttt{timeScale} times canonical source height enters only through
event placement;
supplied event-local Lorentz frames preserve the rest metric. This is an
algebraic direction and causal-carrier result, not a physical sky or
signal theorem. The effective reading additionally requires the physically
faithful images of authenticated local read-from links to cover this \(S^2\)
densely and isotropically; that coverage is not supplied. The measure calibration separately
requires \(\#I/\rho\to\operatorname{Vol}(I)\). Refinement must preserve order,
directions, and count density while soldering rank-three neighborhoods with
\(\lvert B_{\mathrm{car}}(r)\rvert\asymp r^3\) over a scaling window. This is
compatibility of the four inputs, not a fifth one. Neither the local rank-three
carrier nor the seeded global finite wiring determines event-base dimension.

The theorem-status composition is exact at the finite interface and
conditional thereafter. The source-native causal-chart interface uses semantic
commits as the event carrier and the generated reflexive order definitionally;
it uses no separate population map or duplicate precedence field. The finite
compatibility theorem constructs the ambient (1+3) carrier, proves its rank
and source-unit/null identities, and proves the one-way or two-way cone result
under the explicit placement conditions. Supplied chart transports retain
their bounded algebraic covariance role. No physical placement, count-density
field, or refinement family is constructed. Physical faithfulness,
link-direction coverage, manifoldlikeness, topology, volume, and refinement
conditions remain independently testable.

The exact same-event-type source-order Einstein theorem, in a separately
supplied \(3+1\) tensor interface, replaces an all-null
balance hypothesis by nine supplied balances on fixed algebraic
source-direction representatives. Given supplied symmetric geometry/stress
fields indexed by the finite event type, four discrete step maps, Ward/Bianchi
identities, step connectedness, a base event, and an algebraic coupling, exact
nine-direction tomography derives all-null balance, the pointwise metric
ambiguity, its constancy, and the Einstein-form relation. The minimal premise
bundle contains no normalization, vacuum-calibration, Newton-constant, or
universal-source field. The generated order and the independently derived \(1+3\) carrier use
the same event type, but no provenance link selects a tensor field, step, or
balance. The theorem does not identify tensor-coordinate differences with
those of the constructed carrier or identify the supplied fields with smooth
curvature or physical stress-energy. That reading requires the source-causal
continuum, count--volume, convergent curvature/stress, small-ball, coupling,
and scale receipts on one family.

\textbf{Proposition 4.3e0 (Bounded algebraic event-frame soldering).}
Let raw record representatives carry an actual coincidence setoid and a readback into
the intrinsic Hermitian Lorentz module that is constant on coincidence
classes. The readback descends uniquely. Given one time-oriented affine
Lorentz overlap cocycle and, separately, a supplied chart-coordinate family
satisfying its affine overlap law, coordinate displacements transform through
the cocycle's linear part, their Lorentz intervals are chart independent, and
the reverse transition and triple-overlap laws follow from the same cocycle. Future-null
displacements descend to covariant celestial directions. One base-chart
future-unit frame selection induces compatible frames in every chart, and
Lorentz-orthogonal rest projections transport by exact rest-metric
isometries. The rank-three source Gram quotient is linearly and isometrically
equivalent to the standard internal rest fiber, as a candidate local readout.
The standalone quotient-descent theorem does not construct this atlas
coordinate family or identify the descended readback with it.
A finite two-event/two-chart control has nonzero affine translation and a
nonzero future-null displacement. A separate three-record control proves that a
reflexive symmetric pairwise-overlap relation need not be transitive, so it
cannot replace the coincidence setoid. This bounded algebraic adapter
assumes the supplied chart-coordinate overlap law, coincidence setoid,
separation, physical cone attachment, refinement naturality, and operational
clock. It can be used after a source-causal continuum certificate has
selected an effective atlas, but it does not construct the event carrier and
does not supply physical causality, manifoldlikeness, topology, count--volume
calibration, or smooth curvature.

\textbf{Proposition 4.3ea (Finite signed-seam operator realization).}
A declared sign-twisted derivative on a finite seam complex with unit-counting
measure gives a positive semidefinite kinetic operator. Its kernel consists
of parallel sections; on a connected frustrated component that kernel is
zero. The finite realization with \(8{,}662\) sites and \(11{,}816\) seams
has \(38\) frustrated triangles. Exact boundary-rank elimination gives
lift-ambiguity rank \(3{,}117\), and deterministic integer sections check
adjoint, kinetic, covariance, gluing and boundary identities.
The separate rank-three response and declared fifteen-state charge assumptions
define a conditional \(45\)-dimensional internal factor. The declared operator
\(D_\sigma\otimes I_{45}\) retains the gap of \(D_\sigma\) with multiplicity
\(45\). These are properties of the specified finite
operator and matter table. They select neither a spacetime chart nor a
physical matter action, and do not identify the domain with the twelve-port
Spin construction. The microphysics paper gives the finite algebra and
independent replay~\cite{ophfpe}.

\textbf{Theorem 4.3f (Finite null tomography and conditional local-stress upgrade).} A symmetric bilinear form on \((\mathbb R^4,\eta)\) is determined by its null-null values up to a multiple of \(\eta\), and a directional scalar family arises from a tensor iff it satisfies the finite dependent-family linearity relations; nine null directions with resolved rank nine reconstruct the trace-free part. Nine readings always interpolate and therefore do not test tensoriality. Additional dependent directions test consistency within a declared total measurement and numerical error budget; they do not constrain unsampled rays without a family-level premise. For a full-rank design using rays normalized to \(k^0=1\) and a Frobenius-orthonormal trace-free basis, the tensor Frobenius error from charge noise is bounded by \(\delta/s_{\min}\), separately from geometry and solve errors; floating-point singular values are estimates, not interval certificates. The released code verifies this finite linear algebra on declared and synthetic charge families. Values violating a dependent-family relation admit no exact rank-two reconstruction, and scalar collar CMI alone cannot define a rank-two source.

The physical upgrade is conditional. If one common source-derived endpoint-charge family on the event region satisfies the linearity relations, admits the declared same-source identification with local null-stress charges, carries one common normalization and Poincar\'e-covariant local transformation law, and satisfies the weak Ward identity, then tomography defines a symmetric semiclassical \(\langle T_{ab}\rangle\), with null projections equal to those charges and ambiguity \(\phi g_{ab}\) plus the declared improvement terms. Locality, covariance, the charge/stress identification, and weak conservation are separate weak-Ward and null-stress matching supplies; they are not consequences of finite tomography or of the synthetic reconstruction receipt. Universal coupling to the geometric branch is a further named premise. (Spacetime and Einstein paper, stress packet; finite tomography receipts accompany the public geometry code.)

\textbf{Theorem 4.3g (Repaired entropy bridge and absolute Einstein equation).} On the central-interface branch, write the state as a direct sum of bulk states tensored with maximally mixed edge factors. Its entropy splits exactly into \(S_{\mathrm{bulk}}=H(p)+\sum_\alpha p_\alpha S(\rho_\alpha^{\mathrm{bulk}})\) and \(S_{\mathrm{edge}}=\sum_\alpha p_\alpha\log d_\alpha\). At a faithful reference \(\sigma\), the tangent first law is \(\delta S=2\pi\delta\langle B_C\rangle+\delta\langle Z_C\rangle\); for a finite change at fixed \(K_\sigma=2\pi B_C+Z_C=-\log\sigma\), the exact identity is \(S(\rho)-S(\sigma)=2\pi\Delta\langle B_C\rangle+\Delta\langle Z_C\rangle-D(\rho\Vert\sigma)\), with \(Z_C=\sum_\alpha(\log d_\alpha)P_\alpha\). Hence \(\delta S_{\mathrm{bulk}}=2\pi\delta\langle B_C\rangle\) and \(\delta S_{\mathrm{edge}}=\delta\langle Z_C\rangle\) throughout the declared central-interface class, including variations of the sector weights. The supplied nonconstant Gibbs family has the analytic MaxEnt envelope identity \(dS/dt=\lambda\); setting \(\lambda=2\pi\) and identifying \(t\) with physical energy require the declared branch normalization, giving the Clausius balance \(\delta S_{\mathrm{edge}}=\lambda\delta t-2\pi\delta\langle B_C\rangle\) on the coupled variation class. Reaching the geometric balance additionally requires the declared fixed-cap generalized-entropy stationarity interface, the conditional same-source local-stress and weak-Ward bridge, the continuum diamond kernel, and the fixed-volume small-ball area identity. One common scaling family must carry the conditional-mixing, boundary-count, fixed-collar, sharp-rate, and \(o(\ell^4)\) remainder bounds. Coverage of all local timelike directions follows from population plus the produced Lorentz orbit. The vacuum-reference premise fixes the first-variation integration constant. Independent scale readouts with trivial positive-rescaling stabilizer identify the structural coupling. Under all of these premises, \(G_{ab}+\Lambda g_{ab}=8\pi G\langle T_{ab}\rangle\) per connected component. Without the vacuum reference, the result fixes the residue tensor only up to one covariantly constant metric term. (Spacetime and Einstein paper, entropy-bridge and small-diamond packets.)

\textbf{Theorem 4.3h (Conditional Einstein branch composition).} Bare finite consensus is not Einstein-complete: the same finite-consensus reduct admits extensions in which the Einstein equation holds and extensions in which it fails. The valid gravity result is conditional. Let one source-derived repaired quotient tower carry deterministic geometry, modular, causal-order, stress, entropy, and scale readouts on a common domain, with commuting refinement diagrams. Assume the finite source-causal precursor, geometry and null-net receipts; cap-interior modular data; one physical source-causal continuum certificate on a cofinal family with either same-family tensor-curvature reconstruction or the independently premised continuum small-ball/null-balance identification, for which every relevant remainder is \(o(\ell^4)\); the universal-coupling and vacuum-reference receipts; and two independent scale readouts with trivial positive-rescaling stabilizer. The stress-tomography, entropy, curvature-identification, timelike-polarization, Ward, and Bianchi results then compose to
\[
G_{ab}+\Lambda g_{ab}=8\pi G\langle T_{ab}\rangle
\]
on each connected event component. At the curvature-identification step, the
causal-set route requires a same-family tensor-curvature reconstruction
converging to \(G_{ab}\); the alternative continuum route instead premises the
small-ball/null-balance identification independently. Convergence of a
scalar-curvature estimator alone establishes neither route. The exact finite layer includes the nine-direction tomography frame with determinant \(8192/27\) and the entropy and MaxEnt identities. Conditional on the continuum diamond-kernel formula and the smooth fixed-volume small-ball area identity, exact arithmetic gives the coefficients \(8\pi^2\ell^4/15\) and \(-4\pi\ell^4/15\); the continuum formulas themselves remain explicit analytic premises. The finite operator of Proposition 4.3ea does not supply the geometric modular, stress, entropy, scale, or asymptotic data of this theorem. No source tower satisfying the joint continuum antecedent is constructed here, and no semantic-minimality theorem for that antecedent is claimed. (Spacetime and Einstein paper, composed branch-entry theorem.)

The register surface of this composition is inhabited. For every declared
repair law and every horizon record over it, an explicit value of the typed
premise bundle is machine-checked: its modular flow supply is the finite
thermodynamic first-law datum on the simplex tangent space at the declared
reference, with the boost and central charges of the record, so the first-law
clause of the composed conclusion at this value is the first-order
entanglement first law of the horizon-thermality surface on the same
variation space. The other supplies are synthetic witnesses at the demo tier:
the two-patch consensus tower with six typed readouts extending the committed
event and geometry fragment, a sixteen-point four-direction torus chart with a
nonconstant divergence-free stress field and the geometry defined as
\(\kappa T_{ab}+\Lambda\eta_{ab}\), a strict audit manifest evaluated on that
chart, a shrinking family with \(\ell^{5}\) error channels, and small-ball
reals solved from the three small-ball identities. The Einstein clause, the
small-ball arithmetic, the universal-source clause, and the vacuum reference
hold at this value by definition of those witnesses and select nothing. A
nonvacuity receipt exercises the clauses the data can fail on nonzero data,
and a counterexample shows that the split of the horizon record is the
load-bearing input: with a uniform reference and zero central charge in place
of the split, the same boost charge fails the modular-generator clause on the
same variation. The value certifies joint satisfiability of the premises and
places the composed branch on the thermodynamic variation space; it attaches
no physical metric, stress tensor, horizon, temperature, or Newton constant,
and realized nonemptiness on a source-generated tower is not established.

\subsection{Gravity from Fixed-Cap Generalized-Entropy Stationarity}\label{5-gravity-from-entanglement-equilibrium}

\subsubsection{Cap first law}\label{51-cap-first-law}

Section 4.2 fixes the cap modular statement first at the automorphism level. On the extracted geometric subnet, the scaling-limit cap modular flow is geometric with the standard \(2\pi\) normalization,
\[
\sigma_t^{\omega_\infty^{\mathrm{geo},C}}
=
\alpha_{\lambda_C(2\pi t)}.
\]
If the realized scaling-limit cap algebra happens to be type I, one may choose a density matrix \(\rho_C^{\omega_\infty^{\mathrm{geo},C}}\) and modular Hamiltonian
\[
K_C:=-\log\rho_C^{\omega_\infty^{\mathrm{geo},C}},
\]
and then for a perturbation \(\rho(\varepsilon)\) with \(\rho(0)=\omega_\infty^{\mathrm{geo},C}\) the first-law identity reads
\[
\delta S_C=\delta\langle K_C\rangle
=
2\pi\,\delta\langle B_C\rangle.
\]
This displayed variation is the special type-I/fixed-central-sector form of
\(K_C=2\pi B_C+Z_C\), where the central term has no first-order contribution on the chosen
sector. In the generic continuum/non-type-I case, Section 4.2 does not supply an inner operator
identity; the theorem there is the geometric modular automorphism statement, so the cap
first-law formula is restricted to this special realization.

\subsubsection{Null-surface modular bridge}\label{52-null-surface-modular-bridge}

This subsection extends the fixed-cutoff weak-tail-generator boundary to the derived positive null-translation stage, the exact endpoint-derivative identity for its half-line generator, and a conditional local-charge reading. At fixed cutoff one first transfers cut-center data to regulated null strips, then imposes the extra inherited split condition needed for the spatial-collar-type tensor decomposition, obtains exact or controlled strip additivity on one inherited strip model, and derives a weak tail generator for the renormalized half-line family. On the scaling-limit geometric-cap branch of Theorem 4.2, the null half-line blow-up net then inherits geometric dilation and therefore the half-sided covariance relation. Together with the separately supplied common-vector and relative-commutant standardness receipts, this is a standard half-sided modular inclusion, so Borchers--Wiesbrock supplies an explicit positive null-translation generator on its Stone domain. Bounded-interval formulas are discharged by the E0.5 affine/projective kernel below, and the tensor upgrade carries the null-invisible metric ambiguity. Accordingly the null modular bridge consists of:
\begin{itemize}
\tightlist
\item
  null-cut center transfer yielding the central cut-label decomposition, together with the extra inherited split condition needed for the spatial-collar-type tensor decomposition;
\item
  exact or controlled four-term strip additivity on one fixed inherited strip model;
\item
  endpoint-Lipschitz control for the renormalized half-line family and the resulting weak tail generator;
\item
  the geometrically derived half-sided relation plus the registered standardness receipts on the null half-line blow-up net, and the resulting Borchers positive translation generator.
\end{itemize}
The later density-upgrade template is kept below as an explicit downstream template rather than as a fixed-cutoff theorem proved in this subsection. Lemma 5.2f makes the positive null-translation generator itself explicit, and Theorem 5.2g records the Borchers unitary group, positivity and self-adjointness on the Stone domain, the corrected affine commutator relation, and the exact half-line modular-Hamiltonian identity \(K_a=K_0-2\pi aP_\Omega\). Identifying that endpoint derivative with a local null-stress charge additionally consumes an effective-locality and same-source matching premise. Bounded-interval formulas are downstream.

\textbf{Proposition 5.2a (Null-cut center transfer and inherited split).} Fix a regulated null tripartition
\[
I_-=(v_1,v_2),\qquad
J=(v_2,v_3),\qquad
I_+=(v_3,v_4),
\]
with cuts \(\Gamma_-:=\{v=v_2\}\) and \(\Gamma_+:=\{v=v_3\}\). Under the fixed-cutoff realized presentation, assume that the two null cuts inherit the ordinary or central-defect boundary-redundancy data used in the spatial collar branch, so that in a compatible type-I regulator presentation one has
\[
\tilde{\mathcal H}_{I_-}
\cong
\bigoplus_{\alpha_-}
W_{\alpha_-}\otimes \mathcal H_{i_-^{\alpha_-}},
\]
\[
\tilde{\mathcal H}_{J}
\cong
\bigoplus_{\alpha_-,\alpha_+}
W_{\alpha_-}^{*}\otimes
\mathcal H_{j^{\alpha_-,\alpha_+}}\otimes
W_{\alpha_+},
\]
\[
\tilde{\mathcal H}_{I_+}
\cong
\bigoplus_{\alpha_+}
W_{\alpha_+}^{*}\otimes \mathcal H_{i_+^{\alpha_+}},
\]
where opposite sides of each cut carry inverse transport. If the strip algebra is the commutant of the transported cut actions, then
\[
\mathcal A(J)
\cong
\bigoplus_{\alpha_-,\alpha_+}
\mathcal B(\mathcal H_{j^{\alpha_-,\alpha_+}}),
\qquad
Z(\mathcal A(J))
=
\bigoplus_{\alpha_-,\alpha_+}\mathbb C\,P_{\alpha_-,\alpha_+}.
\]
If, in addition, each multiplicity space factors as
\[
\mathcal H_{j^{\alpha_-,\alpha_+}}
\cong
\mathcal H_{j_L^{\alpha_-,\alpha_+}}
\otimes
\mathcal H_{j_R^{\alpha_-,\alpha_+}},
\]
with \(\mathcal A(I_-\cup J)\) acting blockwise only on
\(\mathcal H_{i_-^{\alpha_-}}\otimes \mathcal H_{j_L^{\alpha_-,\alpha_+}}\)
and \(\mathcal A(J\cup I_+)\) acting blockwise only on
\(\mathcal H_{j_R^{\alpha_-,\alpha_+}}\otimes \mathcal H_{i_+^{\alpha_+}}\),
then
\[
\mathcal A(J)
\cong
\bigoplus_{\alpha_-,\alpha_+}
\mathcal B(\mathcal H_{j_L^{\alpha_-,\alpha_+}})
\otimes
\mathcal B(\mathcal H_{j_R^{\alpha_-,\alpha_+}}),
\]
and
\[
\mathcal H_{I_-\cup J\cup I_+}
\cong
\bigoplus_{\alpha_-,\alpha_+}
\mathcal H_{i_-^{\alpha_-}}
\otimes
\mathcal H_{j_L^{\alpha_-,\alpha_+}}
\otimes
\mathcal H_{j_R^{\alpha_-,\alpha_+}}
\otimes
\mathcal H_{i_+^{\alpha_+}}.
\]

\textbf{Proof.} Complete reducibility gives the displayed decompositions. Commuting with the two transported cut actions forces strip operators to preserve the sector pair \((\alpha_-,\alpha_+)\) and act only on the multiplicity space; Schur's lemma kills the off-diagonal intertwiners and leaves the direct-sum algebra above. Taking invariants across both cuts in the glued tripartition again uses Schur's lemma and leaves exactly the matching sector pairs. The left/right split of the multiplicity spaces is additional structure; it is not forced by the center transfer alone. QED.

\textbf{Corollary 5.2b (Exact or controlled four-term null modular relation on an inherited strip model).} On one fixed finite-dimensional strip model satisfying Proposition 5.2a, define
\[
\mathfrak M_{I_-:J:I_+}
:=
\left\{
\sigma:\ I(I_-:I_+\mid J)_\sigma=0
\right\},
\]
and
\[
\delta^{\mathrm M}_{I_-:J:I_+}(\varepsilon)
:=
\sup\left\{
\inf_{\sigma\in\mathfrak M_{I_-:J:I_+}}\|\rho-\sigma\|_1:
I(I_-:I_+\mid J)_\rho\le\varepsilon
\right\}.
\]
Also define the EC-aligned subset \(\mathfrak M^{\mathrm{EC}}_{I_-:J:I_+}\)
by requiring the strip MSA on the inherited split. For a declared controlled
strip family \(\mathcal F_J\), let
\[
\delta^{\mathrm{M,EC}}_{\mathcal F_J;I_-:J:I_+}(\varepsilon)
:=
\sup\left\{
\inf_{\sigma\in\mathfrak M^{\mathrm{EC}}_{I_-:J:I_+}}\|\rho-\sigma\|_1:
\rho\in\mathcal F_J,\ I(I_-:I_+\mid J)_\rho\le\varepsilon
\right\}.
\]
The generic exact-Markov modulus need not control this aligned one; convergence
of the latter is a separate controlled-family hypothesis.
For any strip state \(\eta\), let
\[
\Delta K_J(\eta)
:=
K_{I_-\cup J\cup I_+}(\eta)-K_{I_-\cup J}(\eta)-K_{J\cup I_+}(\eta)+K_J(\eta).
\]

If the strip reference state \(\omega\) is exact Markov and EC-aligned on the inherited split (the strip form of the Markov-split alignment hypothesis of Section 2.3), then \(\Delta K_J(\omega)\) is central, and on the canonical inherited HJPW model one may take
\[
\Delta K_J(\omega)=0.
\]
Equivalently, there exists a central operator \(K_{\partial,J}(\omega)\in Z(\mathcal A(J))\) such that
\[
K_{I_-\cup J\cup I_+}(\omega)
=
K_{I_-\cup J}(\omega)+K_{J\cup I_+}(\omega)-K_J(\omega)+K_{\partial,J}(\omega).
\]

If instead \(\omega\in\mathcal F_J\) and \(I(I_-:I_+\mid J)_\omega\le\varepsilon\), assume the aligned modulus tends to zero and choose \(\widetilde\omega_J\in\mathfrak M^{\mathrm{EC}}_{I_-:J:I_+}\) with
\[
\|\omega-\widetilde\omega_J\|_1
\le
\delta^{\mathrm{M,EC}}_{\mathcal F_J;I_-:J:I_+}(\varepsilon).
\]
Define
\[
\mathfrak D_J(\omega,\widetilde\omega_J)
:=
\Delta K_J(\omega)-\Delta K_J(\widetilde\omega_J).
\]
Then
\[
K_{I_-\cup J\cup I_+}(\omega)
=
K_{I_-\cup J}(\omega)+K_{J\cup I_+}(\omega)-K_J(\omega)
+K_{\partial,J}(\widetilde\omega_J)
+\mathfrak D_J(\omega,\widetilde\omega_J),
\]
with \(K_{\partial,J}(\widetilde\omega_J)=\Delta K_J(\widetilde\omega_J)\in Z(\mathcal A(J))\). Every bounded observable on \(I_-\cup J\) or \(J\cup I_+\) then differs from the exact-Markov reference by at most
\[
\|O\|_\infty\,\delta^{\mathrm{M,EC}}_{\mathcal F_J;I_-:J:I_+}(\varepsilon).
\]
If the relevant strip marginals are uniformly faithful with lower spectral bound \(\lambda_\ast>0\), then
\[
\|\mathfrak D_J(\omega,\widetilde\omega_J)\|_\infty
\le
4\lambda_\ast^{-1}\,\delta^{\mathrm{M,EC}}_{\mathcal F_J;I_-:J:I_+}(\varepsilon).
\]
Independently, Fawzi--Renner gives a recovered comparison state with trace-norm error
\[
r_{\mathrm{FR}}(\varepsilon)=2\sqrt{1-e^{-\varepsilon}}\le 2\sqrt{\varepsilon}.
\]
Thus the exact four-term strip relation is available only at EC-aligned exact Markovity or in a controlled strip family on one fixed inherited strip model with \(\delta^{\mathrm{M,EC}}_{\mathcal F_J;I_-:J:I_+}(\varepsilon_J)\to0\), while \(\mathfrak D_J\) is carried explicitly at finite stage.

\textbf{Proof.} On the inherited split of Proposition 5.2a, the EC-aligned exact-Markov case is the same HJPW block calculation as in the spatial collar theorem: alignment identifies the HJPW blocks with the inherited blocks, the four logarithms cancel blockwise on the canonical model, and any residual bookkeeping term is central. Without alignment, HJPW factorizes the state only over its own state-dependent split, and the centrality of \(\Delta K_J\) in \(Z(\mathcal A(J))\) is not available. The controlled replacement is the same fixed-model compactness argument used for ordinary collars, after relabeling \(A,B,D\mapsto I_-,J,I_+\); the operator-norm bound on \(\mathfrak D_J\) is the corresponding relabeling of the modular-transport estimate. QED.

\textbf{Definition (Renormalized null modular functional).} For a null interval \(I\) on generator \(\Omega\), let \(K_\partial(I,\Omega)\) denote the central endpoint-label term singled out by Corollary 5.2b on the inherited strip model, or by its controlled exact-Markov replacement when that model is used as reference. Define
\[
\widetilde K[I,\Omega]
:=
K[I,\Omega]-K_\partial(I,\Omega),
\qquad
\widetilde K_a(\Omega):=\widetilde K[(a,\infty),\Omega].
\]

\textbf{Proposition 5.2c (Endpoint-Lipschitz null modular families and weak tail generator).} Work on the exact-Markov/central-interface inherited-strip branch, and assume the endpoint-to-regulator-cell map and uniform local norm/site-density bound of the endpoint-control clause in Section 2.7, uniformly on the allowed compact endpoint window. For renormalized modular Hamiltonians on one null generator,
\[
\widetilde K[I,\Omega]
:=
K[I,\Omega]-K_\partial(I,\Omega),
\qquad
\widetilde K_a(\Omega):=\widetilde K[(a,\infty),\Omega],
\]
the additional branch-internal endpoint-control gives:
\[
\bigl|
\langle\psi,(\widetilde K[(a',b'),\Omega]-\widetilde K[(a,b),\Omega])\phi\rangle
\bigr|
\le
C_{\psi,\phi,\Omega}\bigl(|a'-a|+|b'-b|\bigr)
\]
for bounded intervals in a compact endpoint window, and
\[
\bigl|
\langle\psi,(\widetilde K_{a'}(\Omega)-\widetilde K_a(\Omega))\phi\rangle
\bigr|
\le
C_{\psi,\phi,\Omega}|a'-a|
\]
for half-lines. Hence
\[
f_{\psi,\phi}(a):=\langle\psi,\widetilde K_a(\Omega)\phi\rangle
\]
is locally Lipschitz and therefore absolutely continuous, and its distributional derivative defines the weak tail generator
\[
\langle\psi,q(a,\Omega)\phi\rangle
:=
-\frac{1}{2\pi}\,\partial_a f_{\psi,\phi}(a).
\]
For \(a<b\),
\[
\langle\psi,(\widetilde K_b(\Omega)-\widetilde K_a(\Omega))\phi\rangle
=
-2\pi\int_a^b \langle\psi,q(v,\Omega)\phi\rangle\,dv.
\]
\(q(a,\Omega)\) is a weak tail generator rather than a local operator-valued density; its identification with the positive self-adjoint Borchers generator occurs only after Corollary 5.2e and Lemma 5.2f.

\textbf{Proof.} The assumed endpoint-control estimate applies to the renormalized family after removal of the central endpoint term. For bounded intervals, the symmetric-difference length is bounded by \(|a'-a|+|b'-b|\); for half-lines it is exactly \(|a'-a|\). Therefore the matrix-element functions are Lipschitz in the endpoints. A Lipschitz function on a compact interval is absolutely continuous, so \(f_{\psi,\phi}(a)\) has an \(L^\infty_{\mathrm{loc}}\) derivative and obeys the fundamental theorem of calculus. Defining \(q\) as the rescaled negative derivative gives the displayed integral relation. QED.

\textbf{Downstream c.12--c.13 boundary.} The fixed-cutoff bridge established above is the end of this subsection's proved content.

\textbf{Lemma 5.2d (Downstream density-upgrade template).} If, in addition, the weak tail generator \(q(a,\Omega)\) is weakly differentiable in \(a\) and both \(f_{\psi,\phi}(a)\) and \(q_{\psi,\phi}(a)\) vanish at \(+\infty\), then one may define a density
\[
\langle\psi,p(a,\Omega)\phi\rangle
:=
-\partial_a\langle\psi,q(a,\Omega)\phi\rangle
=
\frac{1}{2\pi}\,\partial_a^2 \langle\psi,\widetilde K_a(\Omega)\phi\rangle
\]
and obtain the half-line formula
\[
\langle\psi,\widetilde K_a(\Omega)\phi\rangle
=
2\pi\int_a^\infty (v-a)\,\langle\psi,p(v,\Omega)\phi\rangle\,dv.
\]
This is a downstream density-upgrade step and is not proved from the fixed-cutoff strip arguments above.

\textbf{Null half-line blow-up net.} Fix a smooth cut point and choose affine coordinate \(v\) on generator \(\Omega\) so that the cut sits at \(v=0\). For \(a\ge 0\), write
\[
H_a:=(a,\infty),
\qquad
\mathcal M_a(\Omega)
:=
\overline{\bigvee_{a<c<d<\infty}\mathcal A((c,d),\Omega)}.
\]
By isotony of the null interval net,
\[
a\le b
\quad\Longrightarrow\quad
\mathcal M_b(\Omega)\subseteq \mathcal M_a(\Omega).
\]

\textbf{Corollary 5.2e (Derived half-sided modular relation on null half-lines).} On the scaling-limit geometric-cap branch of Theorem 4.2, assume that the common reference vector is cyclic and separating for every half-line algebra used.  The blow-up modular action near a smooth entangling cut acts on the null coordinate by
\[
v\mapsto e^{-2\pi t}v.
\]
Therefore, for every \(a\ge 0\),
\[
\sigma_t^\omega\bigl(\mathcal M_a(\Omega)\bigr)
=
\mathcal M_{e^{-2\pi t}a}(\Omega).
\]
Hence, for every \(a>0\),
\[
\sigma_t^\omega\bigl(\mathcal M_a(\Omega)\bigr)\subseteq \mathcal M_a(\Omega)
\qquad (t\le 0),
\]
so the inclusion \(\mathcal M_a(\Omega)\subset \mathcal M_0(\Omega)\) obeys the half-sided modular relation. After the harmless convention change \(t\mapsto -t\), this has the standard positive-time orientation. The relation is derived from the null-interval structure, isotony, and the scaling-limit geometric action.  It becomes a \emph{standard} half-sided modular inclusion only when the reference vector is also cyclic for \(\mathcal M_a(\Omega)'\cap\mathcal M_0(\Omega)\); that relative-commutant standardness is a separate receipt, not a consequence of the geometric action.

\textbf{Proof.} Blow up the cap modular flow of Theorem 4.2 near a smooth cut and restrict to the chosen null generator. In the tangent limit the cap-preserving flow becomes the null dilation \(v\mapsto e^{-2\pi t}v\). For every bounded interval \((c,d)\subset H_a\), this sends \(\mathcal A((c,d),\Omega)\) to \(\mathcal A((e^{-2\pi t}c,e^{-2\pi t}d),\Omega)\), and taking the von Neumann closure of the interval net yields the displayed identity for \(\mathcal M_a(\Omega)\). If \(t\le 0\), then \(e^{-2\pi t}a\ge a\), so \(H_{e^{-2\pi t}a}\subseteq H_a\), and isotony gives \(\mathcal M_{e^{-2\pi t}a}(\Omega)\subseteq \mathcal M_a(\Omega)\). QED.

\textbf{Lemma 5.2f (Positive null-translation generator).} Fix \(a>0\) and assume that the common reference vector is cyclic and separating for both algebras below and cyclic for \(\mathcal M_a(\Omega)'\cap\mathcal M_0(\Omega)\).  For the resulting standard half-sided modular inclusion
\[
\mathcal M_a(\Omega)\subset \mathcal M_0(\Omega)
\qquad (a>0),
\]
let \(\Delta_0(\Omega)\) be the modular operator of the standard pair \((\mathcal M_0(\Omega),\omega)\) and define \(K_0(\Omega):=-\log \Delta_0(\Omega)\). Borchers--Wiesbrock then yields a unique strongly continuous one-parameter unitary group
\[
U_\Omega(a)=e^{iaP_\Omega},
\qquad a\in\mathbb R,
\]
such that
\[
U_\Omega(a)\omega=\omega,
\qquad
U_\Omega(a)\mathcal M_b(\Omega)U_\Omega(a)^*=\mathcal M_{a+b}(\Omega)
\quad (a,b\ge0),
\]
and whose generator \(P_\Omega\) is positive and self-adjoint on the Stone domain
\[
D(P_\Omega)
:=
\left\{
\psi\in\mathcal H:
\lim_{a\to0}\frac{U_\Omega(a)\psi-\psi}{ia}\ \text{exists}
\right\}.
\]
In addition,
\[
\Delta_0(\Omega)^{it}U_\Omega(a)\Delta_0(\Omega)^{-it}
=
U_\Omega(e^{-2\pi t}a),
\qquad
\Delta_0(\Omega)^{it}P_\Omega\Delta_0(\Omega)^{-it}
=
e^{-2\pi t}P_\Omega,
\]
and on the common invariant analytic core of \(K_0(\Omega)\) and \(P_\Omega\),
\[
[K_0(\Omega),P_\Omega]=-\,i\,2\pi P_\Omega.
\]
If \(\Delta_a(\Omega)\) denotes the modular operator of \((\mathcal M_a(\Omega),\omega)\) and \(K_a(\Omega):=-\log \Delta_a(\Omega)\), then
\[
\Delta_a(\Omega)=U_\Omega(a)\Delta_0(\Omega)U_\Omega(a)^*,
\qquad
K_a(\Omega)=U_\Omega(a)K_0(\Omega)U_\Omega(a)^*,
\]
so
\[
K_a(\Omega)=K_0(\Omega)-2\pi a\,P_\Omega
\]
as a quadratic-form identity on \(D(K_0(\Omega))\cap D(P_\Omega)\). Equivalently,
\[
\langle\psi,(K_b(\Omega)-K_a(\Omega))\phi\rangle
=
-2\pi(b-a)\langle\psi,P_\Omega\phi\rangle
\]
for all \(\psi,\phi\in D(K_0(\Omega))\cap D(P_\Omega)\). In the canonical normalization of Corollary 5.2b, the weak endpoint derivative from Proposition 5.2c is therefore exactly the Borchers generator. Bounded-interval modular-Hamiltonian formulas are downstream and require the separate interval-preserving branch recorded in Theorem 5.2g.

\textbf{Proof.} Corollary 5.2e gives the half-sided modular relation, and the hypotheses of this lemma supply standardness of the inclusion. Borchers--Wiesbrock then yields the unique strongly continuous unitary group \(U_\Omega(a)\), its positive self-adjoint generator \(P_\Omega\), and the affine covariance relation with the modular group; Stone's theorem gives the displayed domain formula. Differentiating
\[
\Delta_0(\Omega)^{it}U_\Omega(a)\Delta_0(\Omega)^{-it}
=
U_\Omega(e^{-2\pi t}a)
\]
first at \(a=0\) and then at \(t=0\) on the common analytic core gives
\[
[K_0(\Omega),P_\Omega]=-\,i\,2\pi P_\Omega.
\]
Because \(U_\Omega(a)\omega=\omega\) and \(U_\Omega(a)\mathcal M_0(\Omega)U_\Omega(a)^*=\mathcal M_a(\Omega)\), uniqueness of modular data for the translated standard pair gives the displayed formulas for \(\Delta_a(\Omega)\) and \(K_a(\Omega)\). Differentiating \(K_a(\Omega)=U_\Omega(a)K_0(\Omega)U_\Omega(a)^*\) in \(a\) yields \(dK_a/da=-2\pi P_\Omega\) as a quadratic-form identity, and integrating from \(0\) to \(a\) gives the stated half-line modular-Hamiltonian relation. QED.

\textbf{Theorem 5.2g (Half-line generator and conditional EFT null-stress reading).} In the canonical normalization of the preceding half-line construction,
\[
\langle\psi,P_\Omega\phi\rangle
=
-\frac{1}{2\pi}\frac{d}{da}\Big|_{a=0^+}
\langle\psi,\widetilde K_a(\Omega)\phi\rangle
\]
for all \(\psi,\phi\in D(K_0(\Omega))\cap D(P_\Omega)\). If that same renormalized half-line family admits an effective local-stress description and a same-source identification, its null-stress charge is represented by the identical endpoint derivative,
\[
Q_{kk}(\Omega)
:=
-\frac{1}{2\pi}\frac{d}{da}\Big|_{a=0^+}\widetilde K^{\mathrm{eff}}_a(\Omega),
\]
so
\[
P_\Omega = Q_{kk}(\Omega)
\]
as a quadratic-form identity on the common domain. The OPH half-line theorem fixes \(P_\Omega\); reading it as a local null-stress charge consumes the effective-locality and same-source identification premise rather than deriving that premise. Bounded-interval transport through the affine-covariant kernel \(g_I(v)\) is discharged by E0.5 below, and reconstruction of a full tensor from the directional charges is subject to the null-invisible metric ambiguity.

\textbf{Boundary of the null bridge.} The null bridge is an explicit theorem branch rather than an automatic fixed-cutoff identity. Its required data are:
\begin{itemize}
\tightlist
\item
  transferred cut-center data alone give the central sector decomposition of the null strip rather than the left/right HJPW split;
\item
  the extra decomposition-inheritance condition of Proposition 5.2a is exactly what upgrades those sectors to the same collar-type block structure used in the spatial branch;
\item
  the fixed-cutoff null-strip bridge is exact for exact Markov strip states on that inherited decomposition, and otherwise uses a controlled limit with the carried defect operator \(\mathfrak D_J\);
\item
  the exact-Markov/central-interface branch, together with its specified endpoint-to-cell map and uniform norm/site-density bound, gives endpoint-Lipschitz control strong enough to define a weak tail generator for renormalized half-lines; the scaling-limit geometric-cap branch of Theorem 4.2 then derives the half-sided relation on the half-line blow-up net, while common-vector and relative-commutant standardness remain separate receipts;
\item
  with those standardness receipts, the Borchers--Wiesbrock theorem supplies \(P_\Omega\geq0\), its Stone domain, and \(K_a=K_0-2\pi aP_\Omega\);
\item
  the half-line generator and its endpoint-derivative identity are fixed inside the bridge itself; the local null-stress reading requires the separately supplied effective-locality and same-source bridge, while bounded-interval transport and the later tensor upgrade are downstream.
\end{itemize}
This is the route from null-strip factorization to the Einstein branch, with the positive half-line generator proved inside the null modular bridge and its physical stress reading conditional on the named EFT bridge.
\subsubsection{Conditional modular-energy stress-tensor reading on null half-lines}\label{53-modular-energy-as-stress-tensor-charge-uv-cft}

The effective-theory input here is the standard stress-tensor representation that names the operator fixed by the null bridge. In the canonical normalization of the preceding null-half-line construction, the positive null generator is the endpoint derivative of the renormalized half-line modular family:
\[
\langle\psi,P_\Omega\phi\rangle
=
-\frac{1}{2\pi}\frac{d}{da}\Big|_{a=0^+}
\langle\psi,\widetilde K_a(\Omega)\phi\rangle
\]
for all \(\psi,\phi\in D(K_0(\Omega))\cap D(P_\Omega)\).

If that same half-line family admits the supplied effective local-stress description and same-source identification, its local null-stress charge is represented by the identical endpoint derivative,
\[
Q_{kk}(\Omega)
:=
-\frac{1}{2\pi}\frac{d}{da}\Big|_{a=0^+}\widetilde K^{\mathrm{eff}}_a(\Omega),
\]
so on the common quadratic-form domain
\[
P_\Omega = Q_{kk}(\Omega).
\]

Thus the null generator has the displayed local null-stress reading conditional on an effective description that admits a local stress tensor and identifies the endpoint derivative with the same source family. The algebraic half-line construction alone does not supply that effective-locality or source-identification bridge. In a UV CFT regime, or another local Lorentzian regime with local modular Hamiltonians, the conditional reading is the standard null-stress charge associated with the chosen null generator.

For comparison, in a CFT vacuum on a ball one has the familiar local formula
\[
H_{\zeta}=\int_{\Sigma}T_{ab}\zeta^b\,d\Sigma^a,
\]
so, under that EFT premise, the conditional identification is the null-half-line version of the same modular-energy locality.

The downstream boundary is narrower. The E0.5 bounded-interval kernel transports this half-line statement to bounded null intervals with affine-covariant kernel \(g_I(v)\), and the null-to-tensor upgrade determines \(T_{ab}\) only up to the familiar ambiguity
\[
T_{ab}\mapsto T_{ab}+\phi g_{ab}.
\]
The generator is fixed inside the null bridge; its physical local-stress interpretation remains conditional on the stated EFT bridge.

\subsubsection{Localized generalized entropy from Markov + MaxEnt}\label{54-localized-generalized-entropy-from-markov--maxent}

Using the collar decomposition, exact Markovity, and MSA, the declared central-interface branch has the following edge-aligned Markov normal form. On the nonexact branch, Theorem 2.5 supplies a recovered comparison state and a vanishing CMI envelope only under strong conditional mixing; convergence to an exact Markov state is a fixed-collar statement, and alignment with the displayed edge factors remains a separate MSA condition. Under those conditions MaxEnt selection within each edge sector gives

\[
\rho_C = \bigoplus_{\alpha} p_{\alpha}
\left(\rho_{\mathrm{bulk},C}^{\alpha} \otimes \frac{\mathbf 1_{\mathrm{edge}}^{\alpha}}{d_{\alpha}}\right).
\]

The entropy splits as

\[
S(\rho_C) = H(p_{\alpha}) + \sum_{\alpha} p_{\alpha} S(\rho_{\mathrm{bulk},C}^{\alpha}) + \sum_{\alpha} p_{\alpha} \log d_{\alpha}.
\]

\emph{Convention:} Throughout this paper, "log" denotes the natural logarithm (ln), so entropies are measured in \textbf{nats} (1 nat = 1/ln 2 \ensuremath{\approx} 1.443 bits). This is standard in thermodynamics and QFT; the Bekenstein-Hawking formula S = A/4G uses nats. When clarity requires it, we write log\textsubscript{2} explicitly for bits.

Define

\[
S_{\mathrm{bulk}}(C) := H(p_{\alpha}) + \sum_{\alpha} p_{\alpha} S(\rho_{\mathrm{bulk},C}^{\alpha}),
\]

and the central area operator

\[
L_C := \sum_{\alpha} (\log d_{\alpha}) P_{\alpha}.
\]

Then

\[
S_{\mathrm{gen}}(C) := \mathrm{Tr}(\rho L_C) + S_{\mathrm{bulk}}(C).
\]

\textbf{Newton coupling as an edge-entropy area-law readout.}

In the collar double-scaling limit, the edge contribution becomes extensive along the entangling surface \(\Sigma = \partial C\):

\[
\mathrm{Tr}(\rho L_C) \approx N_\Sigma \cdot \bar{\ell}(t),
\qquad
\bar{\ell}(t) := \sum_\alpha p_\alpha \log d_\alpha,
\]

where \(N_\Sigma\) is the number of UV cut elements covering \(\Sigma\) and \(\bar{\ell}(t)\) is the \textbf{single-cell edge entropy} from the heat-kernel distribution (the fixed-cutoff edge-sector law of the gauge-structure paper~\cite{ophGauge}). Similarly, the geometric area is extensive:

\[
A(C) \approx N_\Sigma \cdot a_{\mathrm{cell}},
\]

where \(a_{\mathrm{cell}}\) is the area per UV cut element in the emergent metric.

Matching these expressions gives the geometric area-law coupling:

\[
G_{\mathrm{geom}} = \frac{a_{\mathrm{cell}}}{4 \, \bar{\ell}(t)}
\]

where:

\begin{itemize}
\tightlist
\item
  \(a_{\mathrm{cell}}\) is fixed operationally from the UV correlation/mixing length \(\xi\) via \(a_{\mathrm{cell}} \sim \xi^2\) (from the exponential mixing hypothesis),
\item
  \(\bar{\ell}(t) = \sum_R p_R(t) \log d_R\) is computed from the heat-kernel edge distribution with \(p_R \propto d_R e^{-t\lambda_R}\).
\end{itemize}

Explicitly:

\[
\bar{\ell}(t) = \frac{\sum_R d_R e^{-t\lambda_R} \log d_R}{\sum_R d_R e^{-t\lambda_R}}.
\]

On the OPH gravity branch this formula does not back-solve the scale. The selected
no-\(G\) scale certificate supplies
\(\gamma_\star=\ell_\star\nu_{\mathrm{Cs}}/c\), equivalently
\(B_\star=3\pi/\ell_\star^2\). The local cell readings are
\[
a_{\mathrm{cell}}=P\ell_\star^2,
\qquad
\bar{\ell}_{\mathrm{shared}}=P/4,
\]
so the Newton area-law readout gives
\[
G_{\mathrm{geom}}
=
\frac{P\ell_\star^2}{4(P/4)}
=
\ell_\star^2.
\]
The SI display is \(G_{\mathrm{SI}}=c^3\ell_\star^2/\hbar\).

\subsubsection{Finite MaxEnt identities and the fixed-cap generalized-entropy stationarity interface}\label{55-entanglement-equilibrium-from-maxent}

The proved finite MaxEnt envelope identities hold within the declared state
family. The calculation does not, by itself, identify
a variation of the state constraints with a geometric area variation or prove
stationarity of their sum.  The Einstein branch therefore assumes the
following joint state--geometry interface: for its admissible variations
preserving cap labels (fixed size and charges),

\[
\delta S_{\mathrm{gen}}(C) = 0.
\]

Once this balance is supplied, the split above and the tangent first law
for the bulk term give

\[
\delta S_{\mathrm{gen}}(C) = \delta \langle L_C \rangle + \delta \langle K_{\mathrm{bulk}} \rangle.
\]

\subsubsection{No-smuggling dependency discharge for the Einstein bridge}\label{56a-no-smuggling-dependency-discharge-for-the-einstein-bridge}

\textbf{Theorem E0 (Einstein bridge dependency discharge).} Write \(\mathsf{E}\) for the Einstein-branch premise package: the premises A1--A3 together with the collar-recovery, edge-normalization, and generalized-entropy interfaces consumed by this branch, each an explicit hypothesis outside A1--A3. Let one source-derived cofinal OPH regulator tower carry the typed common-domain geometry, modular, event, stress, entropy, and scale readouts of Theorem 4.3h, together with its finite, asymptotic, and physical-identification premises. On that branch, \(\mathsf{E}\) and the additional receipts supply the inputs used by the Einstein bridge:

\begin{itemize}
\tightlist
\item
  geometry readout through quotient normal forms via the screen fold \(\chi_{S,r}\circ n_r\circ\pi_r\);
\item
  Lorentz/H3 readout from round cap pairs, since \(\mathrm{Conf}^+(S^2)\cong \mathrm{PSL}(2,\mathbb C)\cong \mathrm{SO}^+(3,1)\) and \(H^3\simeq\mathrm{SO}^+(3,1)/\mathrm{SO}(3)\);
\item
  \(2\pi\)-normalized support-visible BW flow on the extracted geometric subnet;
\item
  the positive null generator on the same branch and, under the separately supplied effective-locality and same-source bridge, its local null-stress charge reading;
\item
  the bounded-interval kernel of Lemma E0.5;
\item
  fixed-cap generalized-entropy stationarity from the separately declared stationarity interface;
\item
  the standard fixed-volume small-ball area identity after geometry readout;
\item
  \(o(\ell^4)\) remainder control by Lemma E0.6; and
\item
  all local timelike directions by the cap-pair coverage of Lemma E0.7.
\end{itemize}

Thus the finite consensus reduct supplies quotient normal forms. The typed branch receipts supply the geometric, modular, stress, entropy, coverage, tail, coupling, vacuum, and scale inputs needed by the Einstein bridge. No source-derived tower carrying all of those inputs is constructed here. Bare finite consensus does not imply the Einstein equation.

\textbf{Lemma E0.5 (Conditional bounded-interval kernel).} On the null blow-up of the support-visible geometric branch, assume the effective-locality and same-source bridge that identifies the proved half-line endpoint derivative with a local null-stress charge.  Together with endpoint-Lipschitz renormalized interval control and affine/projective covariance, this gives
\[
K_{(a,b)}^{\mathrm{null}}
=2\pi\int_a^b \frac{(v-a)(b-v)}{b-a}\,T_{kk}(v)\,dv
+R_{(a,b)} .
\]
The remainder \(R_{(a,b)}\) is the carried endpoint/collar term of the same refinement family. In the small-ball limit this is the ball modular kernel used below.

\textbf{Lemma E0.6 (Uniform remainder control with its transfer bounds).}
Let \((r,\ell_r,\delta_r)\) be one declared cofinal family with
\(\ell_r\to0\) and faithful Gibbs states \(\rho_r\), on which the strong conditional matrix-mixing premise
of Theorem 2.5 holds uniformly over the event region and cap family. Set
\[
\varepsilon_r=c|\partial C_r|_{\mathrm{UV}}e^{-\delta_r/\xi_r}.
\]
The boundary-anchor expansion directly controls the operator norm:
\[
\|\Delta K_r(\rho_r)\|_\infty
=\|\mathbf J_{\rho_r}(A_r:D_r\mid B_r)\|_\infty
\le\varepsilon_r.
\]
Here the four canonical marginal logarithms define \(\Delta K\), so
\(\Delta K=-\mathbf J\). Exact faithful Markov references have zero canonical
four-log defect. This direct bound requires no uniform spectral floor;
it does not identify an EC splitting or control each individual modular
Hamiltonian separately.

Assume separately that the needed EC-aligned replacements obey one modulus
\(d_r^{\mathrm{M,EC}}\le C_{\mathrm M}\varepsilon_r^\theta\), with
\(C_{\mathrm M}<\infty\), \(\theta>0\), and that for some
\[
s>\max\{8,4/\theta\},\qquad \omega_r\longrightarrow+\infty,
\]
the same family satisfies
\[
\frac{\delta_r}{\xi_r}
-\log\!\bigl(c|\partial C_r|_{\mathrm{UV}}\bigr)
\ge s\log(\ell_0/\ell_r)+\omega_r.
\]
Then
\[
2\sqrt{\varepsilon_r}=o(\ell_r^4),\qquad
d_r^{\mathrm{M,EC}}=o(\ell_r^4),\qquad
\|\Delta K_r(\rho_r)\|_\infty=o(\ell_r^4).
\]
Indeed \(\varepsilon_r\le(\ell_r/\ell_0)^s e^{-\omega_r}\), and the
exponents are \(s/2>4\), \(s\theta>4\), and \(s>4\), respectively.
For a downstream remainder bounded by a stage-independent constant times
these three quantities, plus a separately controlled \(o(\ell_r^4)\)
kernel residual, the remainder is \(o(\ell_r^4)\). Any growing
transfer constant or further remainder must instead be included in the
same-family tail receipt; the rate calculation does not bound it by itself.

On a weaker route that has only a trace-norm replacement estimate, the
matrix-mixing bound above is unavailable. The logarithm comparison instead
requires the complete rate
\[
4\lambda_r^{-1}d_r^{\mathrm{M,EC}}=o(\ell_r^4),
\]
where \(\lambda_r>0\) bounds the relevant marginals of both states below.
For example, \(\lambda_r\ge\lambda_0(\ell_r/\ell_0)^p\), with
\(\lambda_0>0\) and \(p\ge0\), gives the sufficient strengthened margin
\(s>\max\{8,(4+p)/\theta\}\), under uniformly bounded downstream
transfer constants. Stagewise faithfulness alone supplies no such margin.
This weaker route is not a replacement for the stronger hypotheses
used by Theorem 2.5. Mere pointwise collar convergence also supplies neither
common-family rate, and no per-radius diagonal choice is used.

\textbf{Lemma E0.7 (Cap-pair timelike coverage).} On the cap-normal \(H^3\) branch, every future timelike vector \(v\) has \(r=\sqrt{-\eta(v,v)}\) and \(u=v/r\in H^3\). For any unit \(s\in u^\perp\), the vectors
\[
q_\pm=u\pm s
\]
are future null and
\[
v=\frac r2(q_++q_-).
\]
Thus normalized future null directions in observer skies cover every local timelike direction once the observer-frame chart is supplied. This does not say that two unnormalized projective rays alone select a unique observer frame; the observer normalization or equivalent causal-diamond scale data is also required.

\textbf{Proposition E2 (Conditional recovered-core Ward closure).} Suppose the finite tomography hypotheses reconstruct a tensor from one compatible family of null and timelike charges on the recovered core.  Assume separately that its same-source local charge family obeys the weak Ward balance: signed fluxes on opposite faces of infinitesimal recovered-core diamonds cancel with a controlled remainder, and this balance converges in the recovered charts to the distributional divergence law.  Then \(\nabla^a\langle T_{ab}\rangle=0\) weakly, and pointwise on a smooth branch.  Charge compatibility and tensor tomography alone imply neither the face cancellation nor its continuum limit.

\subsubsection{Einstein equation from fixed-cap generalized-entropy stationarity}\label{56-einstein-equation-from-cap-equilibrium}

This subsection is a branch theorem. The finite consensus reduct supplies quotient normal forms; Theorem E0 discharges the recovered-core premise package \(\mathsf{E}\) needed by the Einstein bridge: geometry readout, support-visible BW modular covariance, the conditional local null-stress reading under effective-locality and same-source matching, bounded-interval transport, the declared fixed-cap generalized-entropy stationarity interface, the fixed-volume small-ball area variation, carried-remainder control, and the timelike scalar-to-tensor upgrade.

In the d=4 scaling regime, the small-ball bridge follows conditionally once the effective-locality, same-source, endpoint-control, and affine/projective hypotheses of E0.5 are supplied. The additional standard geometric input is the fixed-volume area-variation identity for a small geodesic ball. The kernel shape then follows on the extracted geometric subnet: a sufficiently small cap is assumed to admit the tangent causal-diamond description \(D_\ell\), whose preserving conformal Killing field is
\[
\xi_{D_\ell}
=
\frac{1}{2\ell}
\Bigl((\ell^2-r^2-t^2)\,\partial_t-2t\,x^i\partial_i\Bigr).
\]
On the t=0 slice its lapse is
\[
\xi_{D_\ell}\!\cdot n=\frac{\ell^2-r^2}{2\ell}.
\]
The bounded-interval kernel combines the geometric cap generator with the affine/projective transport of the half-line stress bridge, so the cap-modular theorem and the stress bridge together yield
\[
\delta S_{\mathrm{bulk}}(C)
=
2\pi\int_{B_\ell}\frac{\ell^2-r^2}{2\ell}\,\delta\langle T_{00}\rangle\,d^3x
+
\delta\langle E^{(\eta)}_{C,\ell}\rangle.
\]
If \(\delta\langle T_{00}\rangle\) is approximately constant across the ball and the carried remainder is \(o(\ell^4)\), then
\[
\delta S_{\mathrm{bulk}}(C)
=
\frac{8\pi^2\ell^4}{15}\,\delta\langle T_{00}\rangle
+O(\ell^5\partial T)+o(\ell^4).
\]

Stationarity of generalized entropy therefore gives the observer-covariant scalar relation
\[
\delta\!\Bigl[(G_{ab}+\Lambda g_{ab})u^a u^b\Bigr]
=
8\pi G\,u^a u^b\,\delta\langle T_{ab}\rangle
\]
for the local diamond rest-frame four-velocity \(u^a\). In the adapted rest frame \(u^a=(1,0,0,0)\), this is
\[
\delta\!\left(G_{00}+\Lambda g_{00}\right)=8\pi G\,\delta\langle T_{00}\rangle.
\]

\subsubsection{Overlaps supply all timelike directions}\label{57-overlaps-supply-all-timelike-directions}

This scalar-to-tensor upgrade is internal once the Lorentz branch is in place. Overlapping observers through the same bulk point realize every local timelike four-velocity u. Define
\[
Y_{ab}:=G_{ab}+\Lambda g_{ab}-8\pi G\,\langle T_{ab}\rangle.
\]
The local-diamond rest-frame relation says
\[
u^a u^b\,\delta Y_{ab}=0
\]
for every such u. In local inertial coordinates,
\[
u^a=\gamma(1,v^i),
\qquad
|\vec v|<1,
\]
so
\[
0
=
u^a u^b\,\delta Y_{ab}
=
\gamma^2\Bigl(\delta Y_{00}+2v^i\delta Y_{0i}+v^iv^j\delta Y_{ij}\Bigr)
\]
for all \(|\vec v|<1\). The polynomial therefore vanishes coefficientwise, hence
\[
\delta Y_{00}=0,\qquad
\delta Y_{0i}=0,\qquad
\delta Y_{ij}=0.
\]
So the full tensor variation vanishes:
\[
\delta Y_{ab}=0.
\]

Along any connected scaling branch of reference states, \(Y_{ab}\) is therefore constant. The only purely local freedom left by the null reconstruction is the metric term, and fixing that term on one maximally symmetric reference state gives
\[
G_{ab}+\Lambda g_{ab}=8\pi G\,\langle T_{ab}\rangle.
\]
Thus the tensor upgrade does not require a separate EFT handoff once the geometric-cap theorem, the null bridge, and the overlap-complete timelike family are in place.

\subsubsection{Numerical constants of the gravity chain}\label{58-non-tunable-numerical-constants}

The gravity chain carries specific numerical constants. One is supplied by a certificate clause; the other is an identity.

\textbf{The \(2\pi\) KMS normalization.} In Theorem 4.2 the standard \(2\pi\) normalization of cap modular flow on the support-visible scaling-limit geometric cap pair is fixed by the certificate's geometric \(2\pi\)-KMS comparison clause, a supplied receipt, with the same normalization as the Unruh and Hawking temperatures. In the Bisognano--Wichmann setting Lorentz covariance and the spectrum condition force it; here the certificate checks it.

\textbf{The geometric coefficient \(\Omega_{d-2}/(d^2 - 1)\).} This coefficient appears in both (a) the CFT-ball modular Hamiltonian weight integral and (b) the geometric area-variation identity. It is an exact integral identity:

\[
\int_{B^{d-1}_\ell} \frac{\ell^2 - r^2}{2\ell} \, d^{d-1}x
= \frac{\Omega_{d-2} \, \ell^d}{d^2 - 1}.
\]

In \(d = 4\):

\[
\frac{\Omega_2}{4^2 - 1} = \frac{4\pi}{15} \approx 0.8377580409572781.
\]

This is the reason prefactors cancel cleanly when going from \(\delta S_{\mathrm{gen}} = 0\) to the Einstein equation (leaving \(8\pi G\) with the \(2\pi\) fixed by Theorem 4.2).

\textbf{What is predicted.} The framework cleanly separates:

\begin{itemize}
\tightlist
\item
  \textbf{Fixed constants}: \(2\pi\) (KMS period, supplied by the certificate clause), \(\Omega_{d-2}/(d^2-1)\) (geometric coefficient, an identity), the existence of the Einstein form under the stated premises.
\item
  \textbf{Global closure constants}: \(G\) is the conversion between edge entropy and geometric area on the emitted local gravity branch. \(\Lambda\) is not fixed by the local reference-state variation data alone; stable correctable-record closure plus identification of the record carrier with de Sitter horizon entropy fixes \(\Lambda_\star\ell_\star^2=3\pi/N_\star\), while the SI curvature scale also uses the selected OPH scale certificate.
\end{itemize}

\subsubsection{Quantitative Markov error and controlled corrections}\label{59-quantitative-markov-error-and-controlled-corrections}

The role of this subsection is to keep three distinct quantities separate:

\begin{enumerate}
\def\labelenumi{\arabic{enumi}.}
\tightlist
\item
the raw collar conditional mutual information
\[
\varepsilon_\delta:=I(A_\delta:D_\delta\mid B_\delta)_\omega;
\]
\item
the constructive recovery error
\[
r_{\mathrm{FR}}(\varepsilon_\delta)
:=
2\sqrt{1-e^{-\varepsilon_\delta}}
\le
2\sqrt{\varepsilon_\delta};
\]
\item
the exact-Markov replacement error on one fixed faithful collar model,
\[
\eta_\delta^{\mathrm M}
:=
4\lambda_\ast^{-1}\,
\delta^{\mathrm M}_{A_\delta:B_\delta:D_\delta}(\varepsilon_\delta),
\]
where \(\lambda_\ast>0\) is the lower spectral bound needed to compare modular Hamiltonians.
\end{enumerate}

The exact identity for the modular defect expectation is unchanged:
\[
\langle \Delta K_\delta\rangle_\omega
=
-I(A_\delta:D_\delta\mid B_\delta)_\omega
=
-\varepsilon_\delta,
\]
with
\[
\Delta K_\delta
:=
K_{A_\delta B_\delta D_\delta}
-K_{A_\delta B_\delta}
-K_{B_\delta D_\delta}
+K_{B_\delta}.
\]

What changes is the interpretation of later ``exact'' collar formulas. For each fixed finite-dimensional collar one may choose an exact Markov reference state \(\sigma_\delta\) with
\[
\|\omega_\delta-\sigma_\delta\|_1
\le
\delta^{\mathrm M}_{A_\delta:B_\delta:D_\delta}(\varepsilon_\delta),
\]
and a constructive recovered comparison state \(\omega_\delta^{\mathrm{rec}}\) with
\[
\|\omega_\delta-\omega_\delta^{\mathrm{rec}}\|_1
\le
r_{\mathrm{FR}}(\varepsilon_\delta).
\]
These two comparisons do different jobs:

\begin{quote}
\textbf{Bounded-observable control.} For every bounded observable \(X\) supported on \(A_\delta\cup B_\delta\),
\[
\left|
\operatorname{Tr}\!\left[X(\omega_\delta-\omega_\delta^{\mathrm{rec}})\right]
\right|
\le
\|X\|_\infty\,r_{\mathrm{FR}}(\varepsilon_\delta),
\]
and
\[
\left|
\operatorname{Tr}\!\left[X(\omega_\delta-\sigma_\delta)\right]
\right|
\le
\|X\|_\infty\,
\delta^{\mathrm M}_{A_\delta:B_\delta:D_\delta}(\varepsilon_\delta).
\]
The first bound is constructive and dimension-free; the second is the fixed-collar route to the exact Markov set.
\end{quote}

\begin{quote}
\textbf{Modular-additivity control.} If the relevant marginals are uniformly faithful with lower spectral bound \(\lambda_\ast>0\), then
\[
\|\Delta K_\delta(\omega)-\Delta K_\delta(\sigma_\delta)\|_\infty
\le
\eta_\delta^{\mathrm M}.
\]
Since \(\Delta K_\delta(\sigma_\delta)\) is central by exact Markovity, the finite-stage modular defect is within \(\eta_\delta^{\mathrm M}\) of the exact splice or additivity value.
\end{quote}

Therefore the manuscript's exact collar identities arise as limits with a carried remainder
\[
\eta_\delta
:=
r_{\mathrm{FR}}(\varepsilon_\delta)
+
\eta_\delta^{\mathrm M},
\]
together with the separate long-wavelength derivative remainders present in the small-ball expansion. Exact identities at finite collar width require exact Markovity; otherwise one works with \(\eta_\delta\)-controlled approximations and then lets \(\eta_\delta\to0\) in the controlled collar limit.

\textbf{Finite-cutoff Einstein remainder bound.} In the small-ball rest-frame calculation, write
\[
\delta S_C
=
\frac{8\pi^2\ell^4}{15}\,\delta\langle T_{00}\rangle
+
\delta\langle E_C^{(\delta)}\rangle
+
O(\ell^5\partial T),
\]
where \(E_C^{(\delta)}\) packages the carried collar remainder. Then the declared fixed-cap generalized-entropy stationarity interface, together with the proved finite MaxEnt identities, gives
\[
\delta\!\left(G_{00}+\Lambda g_{00}\right)
=
8\pi G\,\delta\langle T_{00}\rangle
+
\mathcal E_\delta,
\]
with
\[
\mathcal E_\delta
:=
\frac{15G}{\pi\ell^4}\,
\delta\langle E_C^{(\delta)}\rangle
+
O(\ell\,\partial T),
\]
and, by the bounds above,
\[
\bigl|
\delta\langle E_C^{(\delta)}\rangle
\bigr|
\le
C_C\,\eta_\delta
\]
for some collar-dependent constant \(C_C\) on the fixed faithful model.
Equivalently, before the controlled collar and small-ball limits, one may write
\[
\bigl|\mathcal E_{\ell,\delta}\bigr|
\le
C_1 r_{\mathrm{FR}}(\varepsilon_\delta)
+C_2\delta^{\mathrm M}_{A_\delta:B_\delta:D_\delta}(\varepsilon_\delta)
+C_3\eta^{\mathrm{reg}}_\delta
+C_4\ell\,\|\partial T\|
+o_\delta(1)+o_\ell(1),
\]
with constants depending only on the fixed collar model, the bounded observable class, and the
small-ball chart. Thus the exact Einstein relation is the controlled limit, not a finite-cutoff
identity inherited from approximate Markovity.

If one wishes to package this carried remainder as an effective anomalous energy density, the natural definition is
\[
\delta\langle T_{00}^{\mathrm{anom}}\rangle
:=
\frac{15}{8\pi^2\ell^4}\,
\delta\langle E_C^{(\delta)}\rangle.
\]
Then
\[
\delta\!\left(G_{00}+\Lambda g_{00}\right)
=
8\pi G\,
\delta\!\left(
\langle T_{00}\rangle
+
\langle T_{00}^{\mathrm{anom}}\rangle
\right)
+
O(\ell\,\partial T).
\]
This is a bookkeeping definition of the finite-stage remainder, not an upgrade of the recovered-core theorem. Its significance is exactly that it is controlled:
\[
\bigl|
\delta\langle T_{00}^{\mathrm{anom}}\rangle
\bigr|
\le
\frac{15\,C_C}{8\pi^2\ell^4}\,
\eta_\delta.
\]
On a separate conditional dark/anomaly branch, this carried scalar remainder
may be identified with a repair stress only after a quotient-invariant
source-localization receipt identifies the same collar operator with the
small-ball Einstein remainder, and after a positive causal packet lift supplies
the full tensor. In SI units that continuation has rest-energy normalization
\[
T_A^{ab}u_a u_b
=
\frac{15\hbar c}{8\pi^2\ell^4}R_{r,\ell}
+O(\epsilon_{\rm loc}),
\]
where \(R_{r,\ell}\) is measured in nats and \(\ell\) is the proper geodesic
small-ball radius. The scalar quantity \(I(A:D\mid B)\) alone is not a stress
tensor: packet factorization, causal directions, an allocation rule, a rest
frame, source-localization map, and transport law are additional continuation
data.

The framework carries the Markov error through three distinct controls:
\begin{enumerate}
\def\labelenumi{\arabic{enumi}.}
\tightlist
\item
\(r_{\mathrm{FR}}(\varepsilon_\delta)\) controls bounded observables at one stage;
\item
\(\delta^{\mathrm M}_{A_\delta:B_\delta:D_\delta}(\varepsilon_\delta)\) controls convergence to the exact Markov normal form on a fixed collar;
\item
\(\eta_\delta^{\mathrm M}\) controls modular-Hamiltonian replacements once faithfulness is assumed;
\item
the Einstein-branch remainder is a carried \(O(\eta_\delta)\) term, not a silent exact identity at finite cutoff.
\end{enumerate}

This is the concrete bridge from ``premises about screens'' to ``precision GR predictions plus a disciplined carried remainder''.

\subsubsection{Focusing/QNEC boundary on the null-modular branch}\label{510-focusingqnec-internalization-via-relative-entropy}

The focusing content does not follow from premises A1--A3. The two focusing statements on the null-modular branch require the additional hypotheses stated below. Quantum focusing is an optional gravity interface and the Einstein theorem does not consume it.

\textbf{What the branch supplies.} The fixed-cutoff null bridge of \S{}5.2 supplies transferred cut-center data and exact-or-controlled strip additivity on the inherited strip model.  The geometric scaling limit supplies the half-sided relation; on the separately registered common-vector and relative-commutant standardness receipts, Borchers--Wiesbrock supplies its translation generator and the affine relation \([K,P]=-i2\pi P\). These are kinematic inputs a QNEC statement consumes; they are not a proof of it or an identification of (P) with a local stress charge.

\textbf{Why relative-entropy monotonicity does not yield QNEC.} Monotonicity under restriction orders the relative entropies \(S(\rho_{R(\lambda)}\|\omega_{R(\lambda)})\) along nested null cuts; it fixes the sign of the first \(\lambda\)-derivative and says nothing about the second. A monotone function of the cut need not be convex, so the step from monotonicity to \(d^2 S_{\mathrm{rel}}/d\lambda^2\ge0\) is unavailable. The entanglement first law \(\delta S=\delta\langle K\rangle\) is a first-order state variation at fixed algebra and reference state; the QNEC second null shape variation is a different object, and expanding the first law to second order in the cut does not produce it. The published QNEC proofs supply exactly the missing machinery: null-surface locality and analyticity hypotheses in the free and super-renormalizable settings~\cite{bousso16}, and modular-flow, causality, and defect-OPE inputs in the general QFT setting~\cite{balakrishnan19}. None of those hypotheses is verified on the reconstructed patch net.

\textbf{Why QFC is not derived from QNEC.} The Quantum Focusing Conjecture is strictly stronger than QNEC: the foundational QFC statement obtains QNEC as its locally parallel limit~\cite{bousso15qfc}, and no implication runs the other way. QNEC together with the Einstein branch (Theorem 5.1) and the Raychaudhuri identity controls the diagonal, stationary-cut part of the generalized-expansion variation and leaves the off-diagonal, shear, counterterm, and general nonstationary-cut content of QFC untouched. The optional quantum-focusing interface therefore carries no independent internal support from this branch.

\textbf{Claim boundary.} The null-modular branch supplies the half-sided
modular structure but not an internal QNEC proof. Such a proof additionally
requires the QFT-side hypotheses named above on the reconstructed net and the
off-diagonal QFC content. Any theorem that consumes focusing must name the
optional quantum-focusing interface explicitly.

\textbf{Theorem 5.1 (OPH Einstein-branch closure in the scaling regime).} Under the premise package \(\mathsf{E}\) on a recovered-core, support-visible scaling branch satisfying the E0 dependency discharge above, the fixed-cap generalized-entropy stationarity condition implies the rest-frame first-variation relation

\[
\delta\!\left(G_{00}+\Lambda g_{00}\right)=8\pi G\,\delta\langle T_{00}\rangle.
\]

If this relation holds for all local directions and reference states in the locally Lorentzian scaling regime, the timelike scalar-to-tensor theorem yields the semiclassical Einstein equation modulo the expected \(\Lambda g_{ab}\) ambiguity. The capacity branch closes that metric term globally under stable whole-fiber correctable-record closure, the unique finite-size slack zero, and the horizon--record identification.

\textbf{Proof sketch.} The recovered-core branch supplies the exact hypothesis package required by the Jacobson-type bridge: quotient-safe geometry readout, \(2\pi\)-normalized BW cap flow, the conditional same-source local null-stress charge, bounded-interval ball kernel with \(o(\ell^4)\) remainder, the declared fixed-cap generalized-entropy stationarity interface, the fixed-volume small-ball area identity, and all timelike directions. The cap first law identifies \(\delta S_C\) with \(\delta\langle K_C\rangle\), while the null bridge together with its effective-locality and same-source premise and the bounded-interval kernel relates the bulk modular variation to the required null stress-tensor charge on the diamond. Since the reference state is maximally symmetric, the fixed-volume small-ball area identity enters at first order in the perturbation, and the declared joint state--geometry balance for admissible fixed-cap variations, combined with the proved MaxEnt envelope identities on the cap-label-preserving state family, yields the displayed rest-frame relation. The timelike-coverage lemma supplies all local timelike directions, the scalar-to-tensor theorem upgrades the relation to the tensor equation modulo \(\Lambda g_{ab}\), and the capacity branch closes the remaining metric term under its separately typed stable-capacity and horizon-record premises. QED.

\subsubsection{Black-hole structural statements and continuation-level spectroscopy}\label{511-discrete-horizon-area-spectrum-and-hawking-emission-established}

The structural black-hole package is a four-step chain: fixed-cutoff edge-center collar
decomposition on the exact-Markov or idealized exact-recoverability horizon carrier, the
small-CMI recoverability reading of interior encoding on that carrier, Hawking/KMS
normalization on the geometric modular branch, and the discrete area spectrum with the leading
semiclassical Schwarzschild first-law relation for an actual sector change. Finite
same-boundary records do not supply physical radiation entropy, exterior time, evaporation
flux, a geometric island or an asymptotic waveform readout.

\textbf{Area eigenvalues from edge sectors.} The central area operator (Section 5.4) is

\[
L_C = \sum_\alpha (\log d_\alpha) P_\alpha,
\]

where \(d_\alpha \in \mathbb{N}\) is the dimension of the edge Hilbert space in sector \(\alpha\). With the normalization \(\mathrm{Tr}(\rho L_C) = \langle A
\rangle / 4G_{\mathrm{geom}}\), the area eigenvalues are

\[
A_\alpha = 4G_{\mathrm{geom}} \log d_\alpha = 4\ell_\star^2 \ln d_\alpha,
\]

where \(\ell_\star^2=\hbar G_{\mathrm{SI}}/c^3\) is the scale-readout area displayed as the
Planck area. Since \(d_\alpha\) is a positive integer, \textbf{areas are discretely spaced} with
logarithmic gaps.

\textbf{Hawking emission energy quantization.} For a Schwarzschild black hole with \(A(M) = 16\pi G^2 M^2/c^4\), define the signed change of the hole for a transition from \(d_{\rm before}\) to \(d_{\rm after}\) by

\[
\Delta_{\rm BH} A = 4\ell_\star^2
\ln(d_{\rm after}/d_{\rm before}).
\]

The differential first law, evaluated on the initial background, gives the leading semiclassical relation

\[
\Delta_{\rm BH}(Mc^2) = k_B T_H
\ln(d_{\rm after}/d_{\rm before})
\]

with \(T_H = \hbar c^3/(8\pi G k_B M)\), whose \(2\pi\)
normalization is fixed by the BW\(_{S^2}\) tangent-limit normalization of
Theorem 4.2. This is not an exact finite-difference identity: using the exact
Schwarzschild mass--area relation, an emission with entropy loss \(\ln k\)
has
\(E_{\rm emit}=Mc^2[1-\sqrt{1-\ln k/S_{\rm BH}}]\), whose expansion is
\(k_BT_H\ln k[1+O(\ln k/S_{\rm BH})]\).

\textbf{Continuation templates.} Spectroscopy needs an additional selection rule. If dominant
emission steps reduce the edge dimension by an integer factor \(k\), then
\(d_{\rm before}=k d_{\rm after}\), the signed black-hole entropy change is \(-\ln k\), and to
leading semiclassical order the emitted spectrum is the comb
\(\Delta E_k=k_BT_H\ln k\), \(\Delta f_k=c^3\ln k/(16\pi^2GM)\). For Kerr remnants, observed
hertz must use \(M_{\rm det}=(1+z)M_{\rm source}\), and the \((k-1)/k\) detailed-balance factor
is a within-line net-response factor, not a derived hierarchy of transition probabilities.
The integer-transition rule, the burst and ringdown templates and their frozen tests are
continuation premises of the cosmology layer and are not part of this paper's chain.

\subsubsection{Classical consequences and precision tests on the conditional branch}\label{512-classical-mechanics-from-emergent-gr}\label{513-precision-gravity-predictions-and-experimental-bounds}

If the physical source-causal continuum, curvature/stress, scale and convergence premises
promote the finite Einstein-form identity to the smooth Einstein equation, its ordinary
consequences apply: the contracted Bianchi identity makes stress-energy conservation automatic,
pressureless matter follows geodesics, and the weak-field limit is Newtonian gravity. Light
bending \(4GM/(c^2b)\), perihelion advance \(6\pi GM/(a(1-e^2)c^2)\) per orbit and gravitational
redshift are inherited consequences of that identification, not independent predictions of the
informational order. On the additional pure two-derivative Einstein--Hilbert action branch the
two transverse-traceless modes propagate on the invariant null cone of the Maxwell modes, so
\((c_{\mathrm{GW}}-c_\star)/c_\star=0\) on that branch, and the quadratic operator has no
Fierz--Pauli hard-mass parameter, \(\mu_{\mathrm{FP,hard}}^2=0\); this is an action-content
statement, not a graviton-pole prediction. The published multi-messenger bound
\(-3\times10^{-15}<(c_{\mathrm{GW}}-c)/c<+7\times10^{-16}\) and the model-dependent dispersion
bound \(m_{\mathrm{grav,disp}}\le1.76\times10^{-23}\,\mathrm{eV}/c^2\) are comparisons with that
branch. The pure Einstein linearization has two tensor polarizations and no dipole channel; it
does not prove that every broader diffeomorphism-invariant effective theory lacks extra modes.

\subsubsection{State--observable error propagation and the dynamical bridge}\label{514-theory-side-error-propagation-from-markov-bounds}

On the additional pure-Einstein action/background branch, the classical equalities of Section~\ref{513-precision-gravity-predictions-and-experimental-bounds} hold. Finite Markov/recovery control supplies a separate and more limited quantitative statement. Two error classes must remain separate. The \emph{state--observable error class} concerns expectation values of specified bounded local observables in a recovered state; the Markov/recovery machinery controls this class. The \emph{dynamical-parameter error class} concerns spectral and dispersion properties of an evolution generator: propagation speeds, particle-pole masses, and correlator poles. The trace-distance chain alone does \emph{not} control those quantities.

\textbf{The key quantitative hook.} From Theorem 3.1, if the target state satisfies

\[
I(A_k : C_k \mid B_k) \le \varepsilon_k,
\]

then, writing \(\rho\) for the target state and \(\sigma\) for its recovered comparison state, the recovery map gives the trace-norm bound

\[
\|\rho-\sigma\|_1\le \delta_k
:=2\sqrt{1-e^{-\varepsilon_k}}
\le2\sqrt{\varepsilon_k}.
\]

\textbf{State--observable bound.} Trace distance gives immediate bounds on expectation-value errors for the declared operator class: bounded operators supported in the recovered region, compared at the time of recovery. Using the standard dual norm inequality:

\[
|\langle O \rangle_\rho - \langle O \rangle_\sigma|
\le \|O\|_\infty \|\rho - \sigma\|_1 = 2 \|O\|_\infty D(\rho, \sigma),
\]

where \(D(\rho, \sigma) = \frac{1}{2}\|\rho - \sigma\|_1\) is the trace distance. This inequality is exactly as strong as it looks and no stronger: the norm \(\|O\|_\infty\) is the Lipschitz constant only for bounded \(O\), and the comparison is an equal-time statement about states.

\textbf{Exponential decay from the mixing hypothesis.} Writing \(w\) for collar width, so it is not confused with the recovery error \(\delta_k\), the mixing assumption (Section 2.3) provides

\[
I_\omega(A_w : D_w \mid B_w)
\le C_{\mathrm{mix}} |\partial C|_{\mathrm{UV}} e^{-w/\xi}.
\]

Combining these gives an explicit precision dial for bounded local observables:

\[
\delta_{\mathrm{step}} \lesssim 2\sqrt{C_{\mathrm{mix}} |\partial C|_{\mathrm{UV}}}
\cdot e^{-w/(2\xi)}.
\]

For an operator normalized by \(\|O\|_\infty\le1\), an independently chosen absolute expectation-error target \(\tau_O\) is guaranteed at one step if \(\delta_{\mathrm{step}}\le\tau_O\); the theorem therefore gives the sufficient CMI condition \(\varepsilon\le(\tau_O/2)^2\). The target \(\tau_O\) applies only to the state--observable bound and cannot be set by importing a gravitational-wave speed tolerance.

\textbf{What the dial does not certify: dynamical parameters.} The gravitational-wave propagation speed and any graviton pole or dispersion-mass parameter are not bounded local observables. They are spectral/dispersion properties of the dynamical generator (poles of retarded correlators, or long-time wave-packet phase and group velocities), so the trace-distance inequality supplies neither their norm nor their Lipschitz constant. Two states can be close in trace distance at one time while being vacua of generators with different dispersion relations, and a small equal-time expectation error does not control phase accumulation over the \(\sim 10^8\)-year propagation baseline of GW170817. Consequently, \(\delta \lesssim 10^{-15}\) is \emph{not} a theory-side certificate for the multimessenger bound \(|c_{\mathrm{GW}}/c - 1| \lesssim 10^{-15}\) of Section~\ref{513-precision-gravity-predictions-and-experimental-bounds}, and the recovery error must not be presented as the precision dial for that published bound. The ideal classical relation \(c_{\mathrm{GW}}=c_\star\) and the vanishing pure-Einstein hard parameter \(\mu_{\mathrm{FP,hard}}^2=0\) belong to the additional action-level branch recorded there; neither stability away from that branch nor a quantum graviton pole follows from the trace-distance chain.

\textbf{Required dynamical-parameter bridge.} A sufficient conditional theorem can be stated precisely, but its hypotheses have not been derived from the Markov remainder in this paper. Let \(\mathcal B=[k_{\min},k_{\max}]\), with \(0<k_{\min}<k_{\max}\), be the declared experimental wave-number band. For each tensor helicity, assume a \(C^1\) self-adjoint finite-regulator generator \(H_\nu^{\mathrm{TT}}(k)\) and the GR reference generator \(H_0^{\mathrm{TT}}(k)\), expressed in angular-frequency units (\(\hbar=1\)) on one specified common finite-dimensional helicity fiber or after a declared \(C^1\) unitary identification of the two fibers. Assume that the reference TT tensor-mode eigenvalue \(\omega_0(k)=c_\star k\) is simple in the helicity block and is separated from the rest of the spectrum by a uniform gap \(g_{\mathcal B}>0\). Define

\[
M_{\mathcal B}
:=\frac{1}{c_\star}\sup_{k\in\mathcal B}
\left\|\partial_k H_0^{\mathrm{TT}}(k)\right\|_\infty
\]

and the dimensionless dynamical error

\[
\eta_{\mathrm{dyn},\nu}
:=\max\!\left\{
\frac{1}{g_{\mathcal B}}\sup_{k\in\mathcal B}
\left\|H_\nu^{\mathrm{TT}}(k)-H_0^{\mathrm{TT}}(k)\right\|_\infty,
\frac{1}{c_\star}\sup_{k\in\mathcal B}
\left\|\partial_k H_\nu^{\mathrm{TT}}(k)-\partial_k H_0^{\mathrm{TT}}(k)\right\|_\infty
\right\}.
\]

For \(\eta_{\mathrm{dyn},\nu}<1/4\), let \(\omega_\nu(k)\) be the continued isolated eigenvalue and define its group velocity by \(c_{\mathrm{GW},\nu}(k):=\partial_k\omega_\nu(k)\). The standard isolated-eigenvalue projector estimate and the Hellmann--Feynman formula then give, uniformly over \(\mathcal B\),

\[
\sup_{k\in\mathcal B}
\left|\frac{c_{\mathrm{GW},\nu}(k)}{c_\star}-1\right|
\le L_{\mathcal B}\,\eta_{\mathrm{dyn},\nu},
\qquad
L_{\mathcal B}:=1+4M_{\mathcal B}.
\]

Indeed, if \(P_\nu(k)\) and \(P_0(k)\) are the rank-one spectral projectors, then \(\|P_\nu-P_0\|_\infty\le2\eta_{\mathrm{dyn},\nu}\), hence \(\|P_\nu-P_0\|_1\le4\eta_{\mathrm{dyn},\nu}\). Differentiating the isolated eigenvalue therefore bounds the group-velocity change by the derivative perturbation plus \(4M_{\mathcal B}c\,\eta_{\mathrm{dyn},\nu}\). This displays a dimensionless constant derived from a specified generator and uniform over the declared band. If \(q:=L_{\mathcal B}\eta_{\mathrm{dyn},\nu}<1\), then stationary propagation at fixed group velocity over path length \(D\) obeys the group-delay bound

\[
|\Delta t|\le \frac{D}{c_\star}\frac{q}{1-q}.
\]

The same bound applies to varying propagation only if it holds pointwise along the entire path; on a cosmological path, the redshifted wave number must remain inside \(\mathcal B\). A wave-packet phase bound likewise requires uniform control of the inverse dispersion \(k_\nu(\omega)\) on the corresponding frequency band. A model-dependent particle dispersion-mass bound is a separate spectral estimate at \(k=0\); it is not a corollary of the band-limited speed estimate.

The OPH-specific bridge requires three additional conditions:

\begin{enumerate}
\def\labelenumi{\arabic{enumi}.}
\tightlist
\item
  \(H_\nu^{\mathrm{TT}}\), or an equivalent retarded two-point kernel, is derived from the finite and scaling OPH dynamics;
\item
  a norm-resolvent, quadratic-form, or correlator estimate has the form
  \(\eta_{\mathrm{dyn},\nu}\le C_{\mathcal B}r_N\), where \(r_N:=\min(2,\sum_{k=2}^N\delta_k)\) is the accumulated trace-norm recovery bound of Theorem 3.1 for the associated \(N\)-step construction and \(C_{\mathcal B}\) is derived rather than assumed; and
\item
  the band-uniform dispersion estimate applies to the detector band and propagation baseline, including phase accumulation where used by the likelihood.
\end{enumerate}

Only after the second item is proved may one combine the estimates to obtain

\[
\sup_{k\in\mathcal B}
\left|\frac{c_{\mathrm{GW},\nu}(k)}{c}-1\right|
\le L_{\mathcal B}C_{\mathcal B}r_N.
\]

No estimate \(\eta_{\mathrm{dyn},\nu}\le C_{\mathcal B}r_N\) is proved here. Substituting \(\delta_k\), \(\delta_{\mathrm{step}}\), or any equal-time trace-distance error for \(\eta_{\mathrm{dyn},\nu}\) without that bridge is invalid.

\textbf{Precision summary.} The framework provides:

\begin{enumerate}
\def\labelenumi{\arabic{enumi}.}
\tightlist
\item
  The ideal pure-Einstein linearization statements \(\mu_{\mathrm{FP,hard}}^2=0\) and
  \(c_{\mathrm{GW}}=c_\star\), without a quantum graviton-pole conclusion.
\item
  Translation of those ideal equalities into the specific observables experiments constrain (Section~\ref{513-precision-gravity-predictions-and-experimental-bounds}).
\item
  Explicit bounds on how far expectation values of bounded local observables in the recovered state can drift, using the conditional mutual information \ensuremath{\to} trace distance \ensuremath{\to} observable error chain.
\end{enumerate}

It therefore provides quantitative error control on recovered states and bounded observables. No theory-side certificate is supplied for the published GW-speed or model-dependent particle dispersion-mass bounds, and no derived dynamical-parameter certificate is supplied.

\subsubsection{The cosmological constant and the dark-sector slot}\label{515-dark-matter-as-modular-anomaly}\label{516-de-sitter-holography-static-patch-vs-boundary-at-infinity}

The Einstein implication leaves two cosmological objects to the cosmology layer. A candidate
anomaly-tensor slot \(T_{ab}^{\mathrm{anom}}\) becomes available only after a finite covariant
source packet supplies source localization, rank-two reconstruction, conservation,
normalization and universal coupling; its rest-frame normalization would be
\(\langle T_{00}^{\mathrm{anom}}\rangle=(15/8\pi^2)\,\delta\langle K_C^{(\mathrm{anom})}\rangle/\ell^4\).
The conditional dark-sector continuation built on that slot is developed
in~\cite{ophdarkmatter}; raw scalar collar CMI supplies only a recovery diagnostic. A
compact neutral coherent source has no monopole and therefore cannot close the declared
single-monopole source ansatz. Bound neutral composites, dipole or higher-multipole sources,
and topological or solitonic architectures are not excluded.

\textbf{The mechanism: \(\Lambda\) as global capacity, not local physics.} Once the separately supplied effective-locality and same-source bridge identifies the directional null charges with one symmetric local stress tensor, finite null tomography determines that tensor only up to an additive metric term. Vacuum-energy or cosmological-constant shifts are therefore invisible to this conditional local null-data route. The conditional Einstein equation obtained under the declared fixed-cap generalized-entropy stationarity interface is correspondingly fixed only up to \(\Lambda g_{ab}\). Accordingly, this null-data route cannot fix \(\Lambda\) without independent trace-sensitive or global input. A source-selected vacuum potential or non-null local equations could supply such input.

On the static-patch branch each timelike observer carries a horizon screen \(S^2\) rather than
a boundary at infinity, and \(\Lambda\) is read as a global capacity through
\(N_{\mathrm{scr}}=3\pi/(\Lambda\ell_\star^2)\) of Section 1.4. The source selection of that
capacity is a cosmology-layer join.

\textbf{Fixed-capacity shock sign.}
The pure de Sitter shock parameter obeys
\[
\mu^2=(d-2)\kappa r_c=d-2=\lambda_{\ell=1}(S^{d-2}),
\]
so the cosmological constant cancels and the smooth isometry modes lie at
zero shock eigenvalue. For finite sector dimensions, maximizing
\[
S_{\rm gen}(p,d)=-\sum_i p_i\log p_i+\sum_i p_i\log d_i
\]
over \(p\) gives \(p_i=d_i/M\) and \(S_{\rm gen}^{\max}=\log M\).
The associated logarithmic area observable
\(\mathcal A=\sum_i(d_i/M)\log d_i\) has symmetric-point Hessian
\[
\left.\operatorname{Hess}\mathcal A\right|_{d_i=d}
=\frac1{nd^2}\left(I-\frac2nJ\right)
\]
on the positive-real relaxation of the integer dimensions. Its curvature is
positive on fixed-\(M\) tangent directions and negative in the homogeneous
direction; the gradient at the symmetric point is nonzero. The physical sign
comes from a one-sided transfer. Assigning a fraction \(f\) of a fixed total
horizon--observer capacity to the observer uniformly depletes the horizon and
gives
\[
\Delta S_{\rm gen}^{\max}=\Delta\mathcal A=\log(1-f)<0.
\]
The finite identity applies to admissible integer depletions, and its
positive-real interpolation is strictly decreasing.
Under the independent horizon-area and fixed-budget dictionaries, this is the
de Sitter time-advance sign studied in
Ref.~\cite{chenStanfordTangYang2026}. It is a boundary maximum along the
transfer direction, rather than an interior maximum at fixed horizon
capacity. Identifying the finite capacity transfer with observer mass
is an additional physical dictionary.

On the icosahedral carrier, the spectrum
\[
\{-2,\ 0,\ 1+3/\sqrt5,\ 1+\sqrt5\}
\]
is conditional on two further statements: the \(A_5\) rotation triplet is an
exact gauge sector, and the shock kinetic term is the scaled
nearest-neighbour port or edge-sector Laplacian. The regular line-graph
identity makes the port and edge-sector low spectra equal. It adds the
eigenvalue \(10\) with multiplicity \(18\) on the edge carrier. No shock
coefficient, physical attachment, or generator for sector-dimension shocks
is supplied. The finite heat-bath repair generator acts at fixed sector
dimensions and therefore supplies none of these objects. The sign mechanism
does not rescue the static-patch trace conjecture~\cite{ophdesittershock}.
That obstruction tests the cited positive cyclic trace interpretation.
Nontrace horizon-screen descriptions, separately constructed
dimension-changing generators, and source-derived kinetic operators beyond
the scaled nearest-neighbour graph remain outside these classes. None of
these boundaries disproves OPH; the pure-de-Sitter normalization, finite
entropy and transfer identities, and regular line-graph theorem survive.

\begin{center}\rule{0.5\linewidth}{0.5pt}\end{center}

\subsection{Einstein-chain theorem and falsification boundary}

The finite support and algebraic packets establish repair-invariant incidence,
the spherical branch receipts, cap-normal identities, the Lorentz action on
the celestial screen, exact finite entropy splitting on the central-interface
branch, and the finite tomography identities used to reconstruct a trace-free
symmetric tensor from null projections.

Authenticated semantic read-from provenance supplies the finite event carrier
and its generated precedence order. Its canonical source height is the
attained longest authenticated-parent-chain length and strictly increases on
generated ancestry.
Separately, every finite decidable strict partial order compiles into an
abstract authenticated log whose direct parent relation and generated order
are exactly the supplied relation. Thus the generic parent relation contains
every supplied predecessor rather than Hasse covers alone; this proves grammar expressivity, not
dynamical or physical selection.
Independently, the exact rank-three positive source Gram quotient and a real
axis define the four-dimensional ambient target carrier
\[
 \mathbb R\oplus V_{\rm src},\qquad
 Q(t,x)=t^2-g_{\rm src}(x,x),
\]
with Lorentz inertia \((1,3)\); source-unit directions are exactly its
future-null rays. Canonical source height enters only as the temporal
coordinate of a positively scaled event placement. Every finite log has an
explicit auxiliary injective forward-causal placement obtained from a finite
enumeration along one source axis and a scale larger than its diameter. It is
not order-reflecting or source-selected. A supplied physical-candidate spatial
readback and edge speed bound give
one-way cone compatibility. Equal-height spatial separation and strict
spacelikeness of unsupported increasing-height pairs imply converse support
and exact two-way order--cone equivalence. Antisymmetry derives injectivity;
the two-way equivalence gives exact interval preservation. Physical event identity,
the spatial readback, the order-faithful-placement conditions, and operational
clock calibration remain separately testable.
A controlled effective route uses vanishing causal discrepancy and
compatible count-volume convergence. Conservative Fibonacci-record
populations with a specified complete-neighbour read law have such a flat
$1+3$ limit when their fill error is small relative to the shrinking read
radius. Equal-mass assignments with vanishing displacement error give
normalized event counts converging to volume. On fixed interior timelike
diamonds, fourth roots of interval-count ratios recover proper-time ratios,
with a reference interval fixing the
unit and complete ancestry access supplying the readout. Within the class of
nonzero closed convex pointed displacement cones, operational invariance
under source rotations and every boost along one axis forces the Lorentz
cone up to time orientation. The covariance is an assumption about actual
influences and readouts. Physical selection of the population, read law and
common matter realization is separate from these conditional results.

The source-direction Einstein theorem works in a separately supplied \(3+1\)
tensor interface indexed by the same finite event type and takes balance on nine
fixed algebraic inverse images of the coordinate null-tomography directions,
symmetric geometry/stress fields in that tensor interface, discrete Ward and
Bianchi identities, an algebraic coupling, four step maps, a base event, and
connectedness. It
derives all-null balance, the pointwise and constant metric ambiguities, and
the Einstein-form tensor identity in that supplied tensor interface. The nine directions are
not observed provenance links, physical rays, or sky samples, and the supplied
fields are not thereby identified with curvature or physical stress-energy.
The informational order shares the event type but is not used in the tensor
calculation: no provenance link selects a direction, field, step, or balance.
The distinguished nine-set is not selected by the source dynamics and is not
invariant as a set under Lorentz or \(\mathrm{SO}(3)\) transformations. The
matrix fields and discrete step maps remain supplied tensor data, and their
coordinate differences are not identified with those of the constructed
carrier. Ambient target-carrier dimension and signature are constructed
independently by the source-causal carrier theorem; intrinsic poset dimension
remains a continuum diagnostic. Neither those fields nor the physical premise
bundle is constructed.

On the exact-embedding route, an effective Lorentzian manifold requires the source-causal
continuum certificate of Theorem 4.3e: one physical refinement family with
exact two-way order/cone agreement, calibrated count--volume, independent
dimension and manifoldlikeness tests, stable topology, and a unique
distinguishing Lorentzian limit. The smooth Einstein reading additionally
requires a same-family tensor-curvature reconstruction converging to the
smooth Einstein tensor, or the separately stated continuum
small-ball/null-balance identification, plus modular and scaling control,
physical stress, Ward/Bianchi, universal coupling, a vacuum reference, and
physical \(8\pi G\) scale identification on that same family. A convergent
causal-set scalar-curvature estimator is an independent diagnostic and is
insufficient by itself. Finite certificate programs verify
their stated algebra and controls; no source-derived family carrying the full
continuum antecedent is constructed here.

The finite native-screen Gaussian vacuum provides a direct test of one
proposed entropy attachment. On a three-site Dirichlet radial slab with the
actual twelve-port angular operator, a fixed-proper-volume uniform metric
variation gives regional bare-entropy derivative \(-1.1004136186\) and
two-face-area derivative \(-4/3\). They have the same sign for every
nontrivial variation in this uniform family, so no positive coefficient of
bare area makes their sum stationary. Equal port-local canonical stress
variations can also give distinct regional entropy derivatives. The slab is
not a geodesic small ball; its finite bare entropy has no renormalized
generalized-entropy or physical stress dictionary. This is a counterexample
to that finite attachment, not to the stated small-ball Einstein theorem.

The implication requires its receipts on one cofinal branch, a rank-three
source Gram quotient, tensor-linearity of the directional charges, the
stated curvature or small-diamond remainder rate, and scale readouts with
trivial common rescaling. The chosen causal route must meet its exact
order--cone criterion or its controlled causal and volume error bounds.

\begin{center}\rule{0.5\linewidth}{0.5pt}\end{center}


